\section{Coulomb branch algebras}
\label{sec:coulomb-branch}
% Regime I: Quadratic, $\Eone$-chiral (Convention~\ref{conv:regime-tags}).

The Yangian $Y(\fg)$ becomes visible in computations only after three
normalizations have been fixed: the RTT kernel
$R_\hbar(u)=1-\hbar P/u$, the ordered configuration parameter
$u=z_i-z_j$, and the completion window in which infinitely many RTT
modes converge weight by weight. The examples below keep these
normalizations explicit. Braverman--Finkelberg--Nakajima Coulomb
branches supply shifted Yangians, Schiffmann--Vasserot and
Kontsevich--Soibelman CoHA constructions supply their geometric
representations, and the low-rank RTT bar complexes test the
bar-cobar comparison on finite spectral windows before passage to
completion.

The finite Drinfeld--Jimbo quantum group, the quantum-loop algebra,
and the Yangian occupy different surfaces. The finite quantum group
governs the KZ monodromy side; the spectral loop parameter belongs to
quantum-loop and Yangian evaluation modules after the exponential or
rational comparison. The ordered bar complex records this distinction
by retaining the order of insertions before the symmetric shadow
forgets it.

\subsection{Definition from gauge theory}

\begin{definition}[Coulomb branch algebra]\label{def:coulomb-branch}
\index{Coulomb branch}
For a 3d $\mathcal{N} = 4$ gauge theory with gauge group $G$ and matter representation $N$,
the \emph{Coulomb branch algebra} $\mathcal{A}(G, N)$ is the quantization of the
Coulomb branch $\mathcal{M}_C(G, N)$.

For type A quiver gauge theories, $\mathcal{A}(G, N)$ is a shifted Yangian or
truncated shifted Yangian.
\end{definition}

\begin{theorem}[BFN construction; \ClaimStatusProvedElsewhere]\label{thm:bfn}
(Braverman--Finkelberg--Nakajima \cite{BFN18}) The Coulomb branch algebra is:
\[
\mathcal{A}(G, N) = H^{G_{\mathcal{O}} \times \mathbb{C}^\times}_*(\mathcal{R}(G, N))
\]
the equivariant Borel--Moore homology of the space of triples $\mathcal{R}(G, N)$.
\end{theorem}

\begin{example}[Quiver of type \texorpdfstring{$A_n$}{A_n}]\label{ex:coulomb-an}
For the $A_n$ quiver with dimension vector $(1, 2, \ldots, n)$:
\[
\mathcal{A} = Y_{(n-1, n-2, \ldots, 1, 0)}(\mathfrak{sl}_n)
\]
a shifted Yangian.
\end{example}


\section{Cohomological Hall algebras}
\label{sec:coha}

\subsection{Definition}

\begin{definition}[Cohomological Hall algebra]\label{def:coha}
For a quiver $Q$ with dimension vector $\mathbf{d}$, the \emph{cohomological Hall algebra} (CoHA) is:
\[
\mathcal{H}_Q = \bigoplus_{\mathbf{d}} H^{G_{\mathbf{d}}}_{\mathrm{BM},*}(\mathcal{M}_{\mathbf{d}}^{\mathrm{nil}})
\]
where $\mathcal{M}_{\mathbf{d}}^{\mathrm{nil}}$ is the moduli of nilpotent representations.

The product is defined via the Hecke correspondence:
\[
\mathcal{M}_{\mathbf{d}_1} \times \mathcal{M}_{\mathbf{d}_2} \xleftarrow{p} \mathcal{E} \xrightarrow{q} \mathcal{M}_{\mathbf{d}_1 + \mathbf{d}_2}
\]
as $a \star b = q_* p^*(a \boxtimes b)$.
\end{definition}

\begin{conjecture}[CoHA as \texorpdfstring{$\Eone$}{E1}-chiral; \ClaimStatusConjectured]\label{conj:coha-e1}
The CoHA $\mathcal{H}_Q$ carries an $\Eone$-chiral algebra structure.
The expected structure has three necessary features:
\begin{enumerate}[label=(\roman*)]
\item The Hecke product is associative but not commutative.
\item The factorization structure comes from stratifying by dimension vectors.
\item The chiral enhancement uses configuration spaces of points on curves
parameterizing deformation directions.
\end{enumerate}
A rigorous construction requires specifying the vertex operators and verifying
weak associativity (the Borcherds identity).
\end{conjecture}

\begin{remark}[Scope]
The obstacle is lifting the Hecke product from correspondences to a vertex operator OPE. The Schiffmann--Vasserot construction \cite{SV13} realizes the CoHA action on Hilbert scheme cohomology as a W-algebra action, but extending this to a full vertex algebra structure on $\mathcal{H}_Q$ itself remains open.
\end{remark}

\begin{example}[Jordan quiver CoHA]\label{ex:jordan-coha}
For the Jordan quiver (one vertex, one loop), the CoHA is:
\[
\mathcal{H}_{\text{Jordan}} = \bigoplus_{n \geq 0} H^{GL_n}_{\mathrm{BM},*}(\mathcal{N}_n) \;\curvearrowright\; \bigoplus_{n \geq 0} H^*_T(\mathrm{Hilb}^n(\bC^2))
\]
The CoHA acts on the equivariant cohomology of the Hilbert scheme via Nakajima-type operators (Schiffmann--Vasserot), realizing a Heisenberg/W-algebra structure on the target. The CoHA itself carries a \emph{different} product structure (the shuffle product) from the OPE product on the chiral algebra side.
\end{example}

\begin{remark}[Toric \texorpdfstring{$\mathbb C^3$}{C3} CoHA versus
\texorpdfstring{$\mathcal W_{1+\infty}$}{W1+infinity}]
\label{rem:coha-c3-positive-half}
For the local CY$_3$ chart $\mathbb C^3$, the cohomological Hall
algebra is the shuffle/positive-half algebra
\[
\mathrm{CoHA}(\mathbb C^3)\;\cong\;Y^+_\hbar(\widehat{\mathfrak{gl}}_1),
\]
not the vertex algebra $\mathcal W_{1+\infty}$ itself. The latter
appears only after passing to an additional construction: Drinfeld
double, derived/chiral centre, Fock-module realisation, or evaluation
representation. Thus the Schiffmann--Vasserot action on Hilbert-scheme
cohomology is a representation-theoretic bridge from the CoHA shuffle
algebra to Heisenberg/$\mathcal W$ operators on the target; it is not
an identification of the source CoHA with the full
$\mathcal W_{1+\infty}$ vertex algebra.
\end{remark}


\subsection{Shifted Yangians and Coulomb branch algebras}
\label{sec:shifted-coulomb}
\index{shifted Yangian!Coulomb branch}
\index{Coulomb branch!shifted Yangian}

The BFN construction (Theorem~\ref{thm:bfn}) realizes Coulomb branch
algebras $\mathcal{A}_C(N,G)$ as equivariant Borel--Moore homology of
spaces of triples on the affine Grassmannian. For type $A$ quiver gauge
theories, these algebras are shifted Yangians
$Y_\mu(\fg)$ (Definition~\ref{def:shifted-yangian}), and their
$\Eone$-chiral structure (Conjecture~\ref{conj:shifted-yangian-e1})
interacts with Koszul duality as follows.

\begin{conjecture}[Langlands-dual Koszul duality for shifted Yangians;
\ClaimStatusConjectured]
\label{conj:shifted-yangian-langlands}
\index{shifted Yangian!Koszul dual}
\index{symplectic duality}
For a coweight $\mu \in X_*(\mathfrak{h})$ and simple $\fg$:
\begin{equation}\label{eq:shifted-langlands-dual}
Y_\mu(\fg)^! \;\simeq\; Y_{-\mu}(\fg^\vee)
\end{equation}
as $\Eone$-chiral algebras, where $\fg^\vee$ is the Langlands dual and
$-\mu$ is the negated coweight (viewed as a weight of $\fg^\vee$ via
the canonical identification $X_*(\mathfrak{h}) \cong X^*(\mathfrak{h}^\vee)$).
For unshifted $\mu = 0$, this reduces to the Yangian Koszul duality
$Y_\hbar(\fg)^! \cong Y_{-\hbar}(\fg)$
(Theorem~\ref{thm:yangian-koszul-dual}) when $\fg$ is simply-laced.
\end{conjecture}

\begin{remark}[Scope]
This is an $\Eone$-chiral incarnation of symplectic duality
(Braden--Licata--Proudfoot--Webster \cite{BLPW16}): the shift
$\mu \mapsto -\mu$ reverses monopole boundary conditions, i.e.\
$R_\hbar(u) \mapsto R_\hbar(u)^{-1}\simeq R_{-\hbar}(u)$ with
Langlands involution. A proof requires
Conjecture~\ref{conj:shifted-yangian-e1} and identification of the
bar complex at coweight $\mu$ with the cobar at $-\mu$ for
$\fg^\vee$.
\end{remark}

\begin{example}[Type \texorpdfstring{$A_1$}{A_1} quiver with framing]
\label{ex:coulomb-a1-framing}
\index{Coulomb branch!type $A_1$}
For the $A_1$ quiver with gauge group $G = GL_1$ and framing
$\mathbf{w} = (n)$ (i.e., $n$ fundamental hypermultiplets), the
Coulomb branch is the Kleinian singularity:
\[
\mathcal{M}_C(n, GL_1) = \{(x,y,z) \in \bC^3 : xy = z^n\}
\]
The quantized Coulomb branch algebra
$A_\hbar = \bC\langle x, y\rangle / (xy - (z+\hbar)(z+2\hbar)
\cdots (z+n\hbar))$ is a shifted Yangian $Y_\mu(\mathfrak{sl}_2)$
with $\mu = (n-1)\omega^\vee$ (fundamental coweight). Its bar complex
$\barB(A_\hbar)$ has:
\begin{align*}
\dim \barB^1 &= 3 \quad (x, y, z) \\
\dim \barB^2 &= 9 \quad (\text{all pairs}) \\
d([x|y]) &= xy = z^n + O(\hbar) \\
d([y|x]) &= yx = z^n + n\hbar z^{n-1} + O(\hbar^2)
\end{align*}
The bar cohomology $H^2(\barB(A_\hbar))$ detects the deformation
of the singularity: $\dim H^2 = n - 1$ (the Milnor number of $xy = z^n$).
The Koszul dual is the Higgs branch algebra, which for this quiver is
$A_{-\hbar} = \bC\langle x', y'\rangle / (x'y' = z^n - n\hbar z^{n-1}
+ \cdots)$, the same singularity with reversed quantization parameter.
\end{example}

\subsection{\texorpdfstring{Cohomological Hall algebras as $\Eone$-chiral algebras}{Cohomological Hall algebras as E1-chiral algebras}}
\label{sec:coha-structure}
\index{CoHA!$\Eone$-chiral structure}
\index{CoHA!factorization structure}

The CoHA product (Definition~\ref{def:coha}) is defined via Hecke
correspondences on moduli of quiver representations. Its factorization
structure (the assembly of $\mathcal{H}_Q = \bigoplus_{\mathbf{d}}
H^{G_{\mathbf{d}}}_{\mathrm{BM},*}(\mathcal{M}_{\mathbf{d}}^{\mathrm{nil}})$
over all dimension vectors) is reminiscent of the factorization algebras
underlying vertex algebras.

\begin{conjecture}[CoHA--Yangian Koszul duality; \ClaimStatusConjectured]
\label{conj:coha-koszul}
\index{CoHA!Koszul duality}
For a quiver $Q$ whose underlying graph is the Dynkin diagram of a
simple Lie algebra $\fg_Q$:
\begin{equation}\label{eq:coha-yangian-dual}
\mathcal{H}_Q^! \;\simeq\; Y(\fg_Q)
\end{equation}
as $\Eone$-algebras, where $\mathcal{H}_Q$ is the CoHA of $Q$
\textup{(}Definition~\textup{\ref{def:coha})} and $Y(\fg_Q)$ is the
Yangian \textup{(}Definition~\textup{\ref{def:yangian})}.
\end{conjecture}

\begin{remark}[Scope]
\begin{enumerate}[label=(\roman*)]
\item For $Q = $ Jordan quiver (one vertex, one loop),
$\mathcal{H}_{\mathrm{Jordan}}$ is the shuffle algebra, which
is isomorphic to $U(\widehat{\mathfrak{gl}}_1)^+$
(Schiffmann--Vasserot \cite{SV13}). Since $Y(\mathfrak{gl}_1)$
is the Yangian of $\mathfrak{gl}_1$ and
$\widehat{\mathfrak{gl}}_1 \cong \mathcal{H}_{\mathrm{Heis}}$,
the duality $\mathcal{H}_{\mathrm{Jordan}}^! \simeq Y(\mathfrak{gl}_1)$
reduces to the Heisenberg Koszul duality
$\mathcal{H}_k^! \simeq \mathrm{Sym}^{\mathrm{ch}}(V^*)$
(Theorem~\ref{thm:heisenberg-not-self-dual}; note:
$\mathcal{H}$ is \emph{not} self-dual).
\item The factorization structure of the CoHA over dimension vectors
$\mathbf{d}$ corresponds to the coproduct structure of the Yangian bar
complex: the Hecke correspondence $\mathcal{M}_{\mathbf{d}_1}
\xleftarrow{p} \mathcal{E} \xrightarrow{q}
\mathcal{M}_{\mathbf{d}_1 + \mathbf{d}_2}$ is the geometric analog
of the bar complex coproduct $\Delta\colon \barB^n \to
\bigoplus_{i+j=n} \barB^i \otimes \barB^j$.
\item At the level of generators, the bar cohomology classes of
$Y(\fg_Q)$ correspond to CoHA generators in the proposed equivalence via the identification
of the RTT bar differential (equation~\eqref{eq:yangian-bar-d2}) with
the Hecke product. The quadratic RTT relations map to the quadratic
part of the shuffle product.
\end{enumerate}
A proof requires first establishing the $\Eone$-chiral structure
on $\mathcal{H}_Q$ (Conjecture~\ref{conj:coha-e1}) and then computing
its bar complex.

The conjecture asserts a Koszul duality equivalence between the CoHA of a Dynkin quiver and the corresponding Yangian as $\Eone$-algebras.
\end{remark}


\section{Explicit Yangian bar complex computations}
\label{sec:yangian-bar-explicit}
\index{Yangian!bar complex!explicit}

\subsection{\texorpdfstring{RTT bar complex for $Y(\mathfrak{sl}_2)$}{RTT bar complex for Y(sl2)}}

\begin{computation}[Yangian \texorpdfstring{$Y(\mathfrak{sl}_2)$}{Y(sl_2)}: bar complex
through degree 3]
\label{comp:yangian-sl2-bar}
\index{Yangian!sl2@$\mathfrak{sl}_2$!bar complex}

The Yangian $Y(\mathfrak{sl}_2)$ in Drinfeld's first realization
has generators $\{e_0, f_0, h_0, e_1, f_1, h_1\}$ with
relations encoding the $\mathfrak{sl}_2$ Lie algebra in Drinfeld
mode~$0$ and the Yangian deformation in mode~$1$.

\emph{RTT presentation.}
In the RTT formalism (Definition~\ref{def:yangian-rtt}), the
algebra is generated by the matrix entries $t_{ij}^{(r)}$
($i,j \in \{1,2\}$, $r \geq 0$) of the generating matrix
$T(u)=\mathbb{1}+\hbar\sum_{r \geq 1}T^{(r)}u^{-r}$ subject to
$R_{12}(u-v) T_1(u) T_2(v) = T_2(v) T_1(u) R_{12}(u-v)$,
where
\[
R_\hbar(u)=\mathbb{1}-\frac{\hbar P}{u}
\]
is the rational $R$-matrix
($P$ is the permutation operator).

\emph{Degree 1.}
$\bar{B}^1(Y(\mathfrak{sl}_2))$ is spanned by the
generators $t_{ij}^{(r)}$ for all $r \geq 1$
(the augmentation ideal excludes $t_{ij}^{(0)} = \delta_{ij}$).
At filtration degree $\leq n$: $\dim \bar{B}^1_{\leq n}
= 4n$ (four matrix entries at each RTT mode $r = 1, \ldots, n$).

\emph{Degree 2.}
Elements $[t_{ij}^{(r)} | t_{kl}^{(s)}] \otimes \eta_{12}$
with $r, s \geq 1$. The coefficient form of the RTT relation is
\begin{align}
[t_{ij}^{(r)},t_{kl}^{(s)}]
&=
\hbar\sum_{a=1}^{\min(r,s)}
\bigl(
t_{kj}^{(a-1)}t_{il}^{(r+s-a)}
-t_{kj}^{(r+s-a)}t_{il}^{(a-1)}
\bigr),
\qquad t_{ab}^{(0)}=\delta_{ab}.
\label{eq:yangian-bar-d2}
\end{align}
The residue convention is
\[
\operatorname*{Res}_{u=0}R_\hbar(u)=-\hbar P,\qquad
\operatorname*{Res}_{u=0}R_{-\hbar}(u)=+\hbar P.
\]
With normalized bar generators dual to the coefficients
$T_{ij}^{(r)}$, the Yangian deformation parameter~$\hbar$ is absorbed
in the pairing. The desuspended degree-$2$ bar corestriction is the
negative of the linear part of~\eqref{eq:yangian-bar-d2}; the
remaining summands are quadratic in positive RTT modes and enter the
higher syzygy differential:
\[
\frac{R_\hbar(u)_{ij,kl} - \delta_{ij}\delta_{kl}}{u}
= -\hbar\,\frac{\delta_{il}\delta_{kj}}{u^2}
\quad (\text{the permutation term})
\]
At mode window $(r,s)$ this gives:
\begin{align}
d([t_{ij}^{(r)} | t_{kl}^{(s)}]) &=
-t_{il}^{(r+s-1)} \delta_{kj} + t_{kj}^{(r+s-1)} \delta_{il}
\label{eq:yangian-d2-explicit}
\end{align}
(the linear RTT corestriction in the bar sign convention).

\emph{Explicit at lowest levels.}
At $r = s = 1$:
\begin{align*}
d([t_{11}^{(1)} | t_{22}^{(1)}])
&= -t_{12}^{(1)}\,\delta_{21} + t_{21}^{(1)}\,\delta_{12} = 0
\quad\text{(diagonal Cartan elements have regular OPE)} \\
d([t_{12}^{(1)} | t_{21}^{(1)}])
&= -t_{11}^{(1)}\,\delta_{22} + t_{22}^{(1)}\,\delta_{11}
= -t_{11}^{(1)} + t_{22}^{(1)} = -h_0^{(1)}
\end{align*}
where $h_0^{(1)} = t_{11}^{(1)} - t_{22}^{(1)}$ is the Cartan
generator in mode~$1$. The vanishing of the first line reflects
the $\mathfrak{gl}_2$ commutation relation
$[E_{11}, E_{22}] = 0$: the bar differential sees only the
singular part of the OPE, which vanishes for commuting generators.

\emph{Degree 3: Serre-type relations.}
At degree~3, the bar differential encodes the Yangian Serre
relations. The key element:
\begin{align*}
d([e_1 | e_1 | f_1] \otimes \eta_{12} \wedge \eta_{13})
&= [e_1 | h_1] \otimes \eta_{13}
+ \text{(other terms)}
\end{align*}
The Serre relation
$[e_1, [e_1, f_0]] - [e_0, [e_1, f_1]] = \frac{1}{4}\{h_0, e_0\}$
appears as a cohomological relation in $H^3(\bar{B}(Y))$.

\emph{Dimension table.}
\begin{center}
\renewcommand{\arraystretch}{1.2}
\begin{tabular}{c|cccc}
Filtration degree & $\dim \bar{B}^1$ & $\dim \bar{B}^2$ &
$\dim \bar{B}^3$ & $\dim H^2$ \\
\hline
$\leq 1$ & 4 & 16 & 64 & 6 \\
$\leq 2$ & 8 & 64 & 512 & 22 \\
$\leq 3$ & 12 & 144 & 1728 & 42
\end{tabular}
\end{center}
Note: the column ``$\dim H^2$'' reports the Koszul dual dimension
$(A^!)_2$ of the truncated RTT \emph{quadratic algebra}
$Y(\mathfrak{gl}_2)_{\leq L}$, which grows with the filtration
window~$L$. This differs from the \emph{chiral} bar cohomology
$\dim H^2(\barBgeom(Y(\mathfrak{sl}_2)^{\mathrm{ch}})) = 10$
(Table~\ref{tab:bar-dimensions}), which uses
$4$ independent generators only (Remark~\ref{rem:yangian-gl2-kunneth}).
\end{computation}


\subsection{Yangian Koszul dual: explicit verification}

\begin{computation}[Yangian duality: \texorpdfstring{$R_\hbar(u) \mapsto R_\hbar(u)^{-1}$}{R_hbar(u) R_hbar(u)^-1}]
\label{comp:yangian-dual-explicit}
\index{Yangian!Koszul dual!explicit}

The $\Eone$-chiral Koszul duality of
Theorem~\ref{thm:yangian-koszul-dual}
sends $R_\hbar(u)$ to~$R_\hbar(u)^{-1}$.
For the rational~$R$-matrix:
\[
R_\hbar(u)=\mathbb{1}-\frac{\hbar P}{u}
=\frac{u-\hbar P}{u},\qquad
R_\hbar(u)^{-1}
=\frac{u}{u-\hbar P}
=\frac{u^2}{u^2-\hbar^2}
\left(\mathbb{1}+\frac{\hbar P}{u}\right).
\]

\emph{Interpretation.}
The dual $R$-matrix $R_\hbar(u)^{-1}$ defines the dual Yangian
$Y(\mathfrak{sl}_2)^!$ via the RTT relation
$R_{\hbar,12}^{-1}(u-v) T_1^!(u) T_2^!(v)
= T_2^!(v) T_1^!(u) R_{\hbar,12}^{-1}(u-v)$.
The scalar factor $(1-\hbar^2/u^2)^{-1}$ is central and invertible
in the completed Laurent ring. Multiplying an RTT relation by such a
scalar gauge does not change its relation span. Hence the dual RTT
algebra is the sign-reversed algebra:
\[
Y_{R_\hbar^{-1}}(\mathfrak{sl}_2)
\;\cong\;
Y_{R_{-\hbar}}(\mathfrak{sl}_2).
\]

\emph{No pole shift.}
The apparent poles of the exact inverse at $u=\pm\hbar$ belong to the
scalar gauge, not to a new collision divisor of the RTT algebra. The
ordered bar complex still has its collision at $u=0$, and the
degree-$2$ corestriction keeps the same mode window:
\[
d^!([t_{ij}^{(r)} | t_{kl}^{(s)}])
= t_{il}^{(r+s-1)} \delta_{kj}
 - t_{kj}^{(r+s-1)} \delta_{il},
\]
the sign-reversed adjoint of~\eqref{eq:yangian-d2-explicit}. This is
the finite-stage bar adjointness of the normalized RTT kernel.
\end{computation}

\begin{proposition}[Scalar-gauge inverse and sign reversal;
\ClaimStatusProvedHere]\label{conj:yangian-spectral-selfdual}
For every finite RTT stage and after weight completion, Koszul
duality sends the normalized rational kernel
$R_\hbar(u)=1-\hbar P/u$ to the scalar-gauge equivalent kernel
$R_{-\hbar}(u)=1+\hbar P/u$. It does not shift the collision
parameter. Equivalently,
\[
A_N(R_\hbar^{-1})\cong A_N(R_{-\hbar}),\qquad
\widehat{\bar B}_{\mathrm{wt}}(Y_\hbar)^!
\simeq
\widehat{\bar B}_{\mathrm{wt}}(Y_{-\hbar})
\]
with the degree-$2$ mode condition $m=r+s-1$ unchanged.
\end{proposition}

\begin{proof}
The identity
\[
R_\hbar(u)^{-1}
=\left(1-\frac{\hbar^2}{u^2}\right)^{-1}R_{-\hbar}(u)
\]
follows from $P^2=1$. The scalar factor is central and invertible in
the completed Laurent ring, hence it multiplies the RTT relation by a
unit and preserves the relation span at every finite mode cutoff. The
bar differential is dual to that relation span, so the only change is
the sign of the permutation residue
$\operatorname*{Res}_{u=0}R_{\pm\hbar}(u)=\mp\hbar P$; the summation
condition in the coefficient extraction remains $m=r+s-1$.
\end{proof}

\begin{remark}[K\"unneth decomposition for \texorpdfstring{$\widehat{\mathfrak{gl}}_2$}{gl_2}]
\label{rem:yangian-gl2-kunneth}
\index{Yangian!Kunneth decomposition@K\"unneth decomposition}
In the chiral bar complex, the mode-$1$ generators
$T_{ij}^{(1)}$ ($1 \leq i,j \leq 2$) span
$\mathfrak{gl}_2 = \mathfrak{sl}_2 \oplus \mathfrak{gl}_1$, and
the OPE among them reproduces the $\widehat{\mathfrak{gl}}_2$
current algebra (double pole from the level, simple pole from
the Lie bracket). Higher RTT modes $T_{ij}^{(r)}$
($r \geq 2$) appear as composites
(normal-ordered products of mode-$1$ fields).
By the K\"unneth formula for bar complexes
of direct products:
\[
H^n(\barBgeom(\widehat{\mathfrak{gl}}_2))
= \sum_{i+j = n}
H^i(\barBgeom(\widehat{\mathfrak{sl}}_2))
\otimes
H^j(\barBgeom(\widehat{\mathfrak{gl}}_1))
\]
with $H^i(\widehat{\mathfrak{sl}}_2) = 3, 5, 7, 9, 11, \ldots = 2i{+}1$
(Table~\ref{tab:bar-dimensions}) and
$H^j(\widehat{\mathfrak{gl}}_1) = p(j{-}2)$ (partition function,
$p({-}1) = p(0) = 1$). This gives
$H^n(\widehat{\mathfrak{gl}}_2) = 1, 4, 9, 16, 30, \ldots$
At degree~$1$, the Yangian bar cohomology matches
$\widehat{\mathfrak{gl}}_2$ exactly ($H^1 = 4$).
At degree~$2$, the Yangian exceeds the K\"unneth prediction:
$H^2(Y) = 10 > H^2(\widehat{\mathfrak{gl}}_2) = 9$.
At degree~$3$, the gap widens to $28 - 24 = 4$.
\end{remark}


\subsection{Coulomb branch bar complex}

\begin{computation}[Coulomb branch and Yangian bar]
\label{comp:coulomb-bar}
\index{Coulomb branch!bar complex}

For the Jordan quiver with dimension vector $\mathbf{d} = (n)$
and framing vector $\mathbf{w} = (1)$, the Coulomb branch
$\mathcal{M}_C(n, 1)$ is:
\[
\mathcal{M}_C(n, 1)
= \operatorname{Spec}\bigl(
\bC[z_1, \ldots, z_n]^{S_n}[\hbar]\bigr)
\cong \bC^n/S_n \times \bC_\hbar
\]
(the symmetric product of the affine line, quantized by $\hbar$).

The quantized Coulomb branch algebra
$A_\hbar(n, 1)$ is generated by
$e_k = \sum_{i=1}^n z_i^k$ ($k \geq 1$) with
the shifted Yangian relation:
\[
[e_k, e_l] = \hbar \cdot k \delta_{k+l, 0} \cdot n
+ \text{(lower terms)}
\]

\emph{Bar complex.}
The bar complex $\bar{B}(A_\hbar(n,1))$ has degree-2
elements $[e_k | e_l] \otimes \eta_{12}$ with differential:
\[
d([e_k | e_l]) = [e_k, e_l]
= \hbar k n \delta_{k+l} + \sum_{j=1}^{k-1}
\binom{k}{j} e_{k-j} e_{l+j} - (k \leftrightarrow l)
\]
At $n = 1$ (the Heisenberg case): $A_\hbar(1,1) \cong
\mathcal{H}_\hbar$ (the Heisenberg algebra at level $\hbar$),
and the bar complex reduces to the Heisenberg bar complex
of Computation~\ref{comp:heisenberg-deg3-full}.

At $n = 2$: the algebra $A_\hbar(2,1)$ contains both the
Heisenberg subalgebra (generated by $e_1 = z_1 + z_2$) and
a relative subalgebra (generated by $e_2 - e_1^2/2
= (z_1-z_2)^2/2$). The bar complex detects the interaction
between these two sectors via the residue at the diagonal
$z_1 = z_2$ in the symmetric product.
\end{computation}


\section{Higher-rank Yangians and quantum groups}
\label{sec:higher-rank-yangians}
\index{Yangian!higher rank}

\subsection{\texorpdfstring{Yangian $Y(\mathfrak{sl}_3)$: bar complex}{Yangian Y(sl3): bar complex}}

\begin{computation}[\texorpdfstring{$Y(\mathfrak{sl}_3)$}{Y(sl_3)} bar complex through degree 2]
\label{comp:yangian-sl3-bar}
\index{Yangian!sl3@$\mathfrak{sl}_3$!bar complex}

The fundamental RTT calculation is first a calculation for
$Y_\hbar(\mathfrak{gl}_3)$: the $R$-matrix acts on
$V=\mathbb C^3$, and the generators are the entries of the
$3\times 3$ matrix
$T(u) = (t_{ij}(u))_{1 \leq i,j \leq 3}$ with
$t_{ij}(u) = \delta_{ij} + \sum_{r=1}^\infty t_{ij}^{(r)} u^{-r}$,
subject to the RTT relation with $R$-matrix
$R_{12}(u) = 1 - P_{12}/u$.
The $\mathfrak{sl}_3$ Yangian is obtained from this RTT window after
the usual quantum-determinant/trace reduction. The table below is
therefore a finite RTT matrix-window computation, not a closed formula
for the trace-reduced chiral bar cohomology.

\emph{Degree 1.}
The generators $\{t_{ij}^{(r)} : 1 \leq i,j \leq 3, \; r \geq 1\}$
form the desuspended bar generators $s^{-1}t_{ij}^{(r)}$.
At filtration degree $\leq 1$: $\dim \bar{B}^1 = 9$
(the nine entries of the first-order matrix $T^{(1)}$).
At filtration degree $\leq n$: $\dim \bar{B}^1 = 9n$.

\emph{Degree 2.}
The bar differential encodes the RTT relation.
For generators at levels $r, s \geq 1$:
\[
[t_{ij}^{(r)},t_{kl}^{(s)}]
=
\sum_{a=1}^{\min(r,s)}
\bigl(
t_{kj}^{(a-1)}t_{il}^{(r+s-a)}
-t_{kj}^{(r+s-a)}t_{il}^{(a-1)}
\bigr),
\qquad t_{ab}^{(0)}=\delta_{ab}.
\]
The degree-$2$ bar corestriction records the desuspended linear part:
\[
d([t_{ij}^{(r)} | t_{kl}^{(s)}])
=-\delta_{kj}t_{il}^{(r+s-1)}
 +\delta_{il}t_{kj}^{(r+s-1)}.
\]
The quadratic positive-mode summands in the RTT coefficient relation
enter the degree-$3$ syzygy differential.

\emph{Serre relations.}
For $\mathfrak{sl}_3$, the Serre relations are:
\[
[t_{12}^{(1)}, [t_{12}^{(1)}, t_{23}^{(1)}]] = 0, \qquad
[t_{23}^{(1)}, [t_{23}^{(1)}, t_{12}^{(1)}]] = 0.
\]
These appear as degree-3 bar cocycles: the element
$[t_{12}^{(1)} | t_{12}^{(1)} | t_{23}^{(1)}]
- 2[t_{12}^{(1)} | t_{23}^{(1)} | t_{12}^{(1)}]
+ [t_{23}^{(1)} | t_{12}^{(1)} | t_{12}^{(1)}]$
is exact in $\bar{B}^3$, reflecting the Serre relation.

\emph{Dimension table.}
\begin{center}
\renewcommand{\arraystretch}{1.2}
\begin{tabular}{c|ccc}
Filtration degree & $\dim \bar{B}^1$ & $\dim \bar{B}^2$ &
$\dim H^2$ \\
\hline
$\leq 1$ & 9 & 81 & 27 \\
$\leq 2$ & 18 & 324 & 99
\end{tabular}
\end{center}
These numbers are not obtained by multiplying the
$\mathfrak{sl}_2$ table. Rank enters through the full matrix-index
set and through the rank of the RTT relation map; trace reduction and
higher RTT relations must be imposed before reading the result as a
statement about $Y_\hbar(\mathfrak{sl}_3)$ itself.
\end{computation}

\begin{proposition}[Finite-window rank dependence of the Yangian bar complex;
\ClaimStatusProvedHere]
\label{prop:yangian-rank-dependence}
\index{Yangian!bar complex!rank dependence}

Let $V_N$ be the trace-reduced first RTT window for
$Y_\hbar(\mathfrak{sl}_N)$, so
$\dim V_N=N^2-1$. Then:
\begin{enumerate}
\item $\dim \bar{B}^1_{\leq 1}=N^2-1$.
\item The raw ordered degree-$2$ chain group has
 $\dim \bar{B}^2_{\leq 1}=(N^2-1)^2$.
\item The associated-graded commutative RTT window has quadratic-dual
 degree-$2$ dimension $\binom{N^2-1}{2}$. Before trace reduction,
 the corresponding $\mathfrak{gl}_N$ number is $\binom{N^2}{2}$.
\item The actual degree-$2$ bar cohomology of a finite RTT window is
 $\ker(d_2)/\operatorname{im}(d_3)$ and depends on the imposed
 quantum-determinant, Serre, and higher RTT relations. It is not
 determined by $N$ and the raw chain dimensions alone.
\item For fixed bar degree~$n$ and fixed filtration window, the raw
 chain dimension is $O(N^{2n})$.
\end{enumerate}
\end{proposition}

\begin{proof}
(i) At filtration degree~$1$, the generators are
$\{t_{ij}^{(1)} : 1 \leq i,j \leq N\}$ modulo the trace
condition, giving $N^2 - 1$ generators (matching $\dim \mathfrak{sl}_N$).

(ii) The degree-2 bar space $\bar{B}^2$ is spanned by
$\{[t_{ij}^{(1)} | t_{kl}^{(1)}]\}$, giving $(N^2-1)^2$ elements.

(iii) Passing to the associated graded kills the lower bracket term
and leaves the symmetric algebra on $V_N$; its quadratic dual has
degree-$2$ part $\wedge^2 V_N^*$, of dimension $\binom{N^2-1}{2}$.
The unreduced RTT matrix window uses $V=\mathfrak{gl}_N$, giving
$\binom{N^2}{2}$.

(iv) In the genuine finite RTT window the differential refines
the associated-graded commutative differential. It is dual to the
RTT relation span, and $d_3$ sees the cubic/Serre consequences.
Thus the cohomology is computed by the displayed kernel modulo image;
no dimension formula follows from the raw counts alone.

(v) A bar word of length~$n$ chooses $n$ generators from a vector
space of dimension $O(N^2)$, hence has raw dimension $O(N^{2n})$.
\end{proof}

\begin{computation}[The
\texorpdfstring{$\mathfrak{sl}_3$}{sl3} KZ residue and fundamental
Yang seed; \ClaimStatusProvedHere]
\label{comp:sl3-yangian-from-ordered-bar}
\index{Yangian!sl3 KZ residue and Yang seed@$\mathfrak{sl}_3$ KZ residue and Yang seed|textbf}
\index{sl3@$\mathfrak{sl}_3$!Yangian rank-2 verification}
For the KZ/DK comparison starting from
$\widehat{\mathfrak{sl}}_3$ at affine level~$k$, the ordered bar
computes the classical Casimir residue. The non-scalar fundamental
Yang/FRT seed on
$V = \mathbb{C}^3$ takes Yang's additive form
\begin{equation}
\label{eq:sl3-yang-r-matrix}
R^{\mathfrak{sl}_3}(u) \;=\; u\cdot I + P,
\end{equation}
where $P\in\End(V\otimes V)$ is the permutation operator. The Casimir
$\Omega = \tfrac12\sum_a \lambda_a\otimes\lambda_a$
\textup{(}Gell-Mann basis, fundamental trace normalization\textup{)}
acts on $V\otimes V \cong \operatorname{Sym}^2 V \oplus \wedge^2 V$
with eigenvalues $+2/3$ on the symmetric six-dimensional summand and
$-4/3$ on the antisymmetric three-dimensional summand. In the
fundamental representation,
$\Omega^{\mathrm{fund}} = P - I/3$, so the additive Yang form
\eqref{eq:sl3-yang-r-matrix} is recovered up to normalization from the
Casimir $r$-matrix.

The collision residue from the bar complex on the ordered configuration
space gives the trace-form affine residue
\[
r^{\mathrm{coll}}_k(z) \;=\; \frac{k\,\Omega_{\mathrm{tr}}}{z}.
\]
The KZ connection uses the different coefficient
\[
r^{\mathrm{KZ}}_k(z) \;=\;
\frac{\Omega_{\mathrm{KZ}}}{(k+h^\vee)z}.
\]
With the trace-form convention used for the collision residue,
\(\Omega_{\mathrm{KZ}}=2h^\vee\Omega_{\mathrm{tr}}\); for
\(\mathfrak{sl}_3\), \(h^\vee=3\), so
\[
r^{\mathrm{KZ}}_k(z)
=\frac{6\,\Omega_{\mathrm{tr}}}{(k+3)z}.
\]
These two rational kernels are related only after the Casimir
normalisation and coupling convention have been fixed; they are not
equal as functions of the affine level. In the fundamental
representation the Casimir identity
\(\Omega^{\mathrm{fund}}=P-I/3\) identifies the non-scalar part of the
KZ kernel with the permutation seed \(P\). The scalar \(I/3\) term
does not affect the Yang--Baxter equation. The
spectral parameter Yang--Baxter equation
\[
R_{12}(u)\,R_{13}(u+v)\,R_{23}(v)
\;=\;
R_{23}(v)\,R_{13}(u+v)\,R_{12}(u)
\]
holds for the Yang/FRT operator \(R(u)=uI+P\) in the fundamental. It is
not inferred from the CYBE alone, and the naive formula
\(uI+\Omega_{\mathrm{adj}}\) does \emph{not} give the adjoint
specialization of the universal Yangian \(R\)-matrix. The adjoint Casimir on
$V_{\mathrm{adj}}\otimes V_{\mathrm{adj}}$ has eigenvalue spectrum
$\{\,2\;(\text{mult }27),\;\;
0\;(\text{mult }20),\;\;
-3\;(\text{mult }16),\;\;
-6\;(\text{mult }1)\,\}$,
matching the Clebsch--Gordan decomposition
$\mathbf{8}\otimes\mathbf{8}
= \mathbf{27}\oplus\mathbf{10}\oplus\mathbf{10}^\ast
\oplus\mathbf{8}\oplus\mathbf{8}\oplus\mathbf{1}$.
\end{computation}

\begin{remark}[Why Yang's additive form is fundamental]
\label{rem:sl3-yangian-fundamental-only}
\index{R-matrix!fundamental representation only}
For the fundamental representation of $\mathfrak{sl}_N$ one has
$\Omega = P - I/N$, so Yang's $R$-matrix \(uI+P\) is recovered up to
scalar gauge and spectral shift. For higher representations the universal $R$-matrix
of the Yangian $Y(\mathfrak{sl}_N)$ has additional structure beyond
$R(u) = uI + \Omega$, and the spectral parameter Yang--Baxter equation
holds only for the universal $R$-matrix specialized in pairs of
representations, not for the naive Yang form. The
$\mathfrak{sl}_3$ case is the first rank where this distinction is
visible: the adjoint $\mathbf{8}\otimes\mathbf{8}$ already requires
the genuine universal $R$-matrix because the Casimir spectrum is
nondegenerate on four nontrivial summands.
\end{remark}


\subsection{Quantum group connection}

\begin{remark}[Yangian and quantum-loop bar complexes]
\label{rem:yangian-quantum-group}
\index{quantum group!bar complex}

The Yangian $Y(\fg)$ and the quantum-loop algebra
$U_q(\widehat{\fg})$ are related by the rational/trigonometric
dichotomy:
\begin{center}
\renewcommand{\arraystretch}{1.2}
\begin{tabular}{l|cc}
& \textbf{Yangian $Y(\fg)$} & \textbf{$U_q(\widehat{\fg})$} \\
\hline
$R$-matrix type & Rational: $R_\hbar(u) = 1-\hbar P/u$ &
Trigonometric: $R(u) = q^{u} P + q^{-u}(1-P)$ \\
Spectral parameter & Additive: $u - v$ & Multiplicative: $u/v$ \\
Bar propagator & $\eta_{ij} = d\log(z_i - z_j)$ &
$\omega_{ij} = d\log(z_i/z_j)$ \\
FM space & $\overline{\mathrm{FM}}_n(\mathbb{C})$ &
$\overline{\mathrm{FM}}_n(\mathbb{C}^\times)$ \\
Koszul dual & $Y_{-\hbar}(\fg)$ & $U_{q^{-1}}(\widehat{\fg})$
\end{tabular}
\end{center}

The finite Drinfeld--Jimbo algebra $U_q(\fg)$ is the KZ monodromy
target after Kazhdan--Lusztig transport. It has no additive spectral
parameter by itself. The spectral parameter belongs to
$U_q(\widehat{\fg})$ and to Yangian evaluation modules, and enters
the finite side only after choosing the quantum-loop/evaluation
comparison.

At roots of unity this separation becomes load-bearing. The
semisimplified Kazhdan--Lusztig target, Lusztig's small quantum group
\(u_q(\fg)\), and the quantum-loop evaluation category are three
different surfaces. Specializing \(q\) to a root of unity does not
produce a Yangian, and a PBW count for \(u_q(\fg)\) is not a Yangian
evaluation-module dimension unless a separate comparison functor is
constructed.

The bar complex of $U_q(\widehat{\fg})$ uses the trigonometric
$R$-matrix as the degree-2 differential. The FM compactification
$\overline{\mathrm{FM}}_n(\mathbb{C}^\times)$ replaces the
rational propagator $d\log(z_i - z_j)$ with the trigonometric
propagator $d\log(z_i/z_j)$, and the boundary strata of the
FM space encode the coproduct structure of the quantum-loop algebra.

At $q = 1$, the quantum-loop algebra degenerates to the enveloping
algebra of the loop algebra and the trigonometric
$R$-matrix becomes the identity, reducing the bar complex to
the current-algebra Chevalley--Eilenberg complex. This is a
trigonometric degeneration, not an assertion that the finite
Drinfeld--Jimbo algebra is a Yangian.
\end{remark}


\subsection{Transfer matrix and Baxter Q-operator}

\begin{computation}[Transfer matrix in the bar complex]
\label{comp:transfer-matrix-bar}
\index{transfer matrix!bar complex}

The transfer matrix of the Yangian $Y(\mathfrak{sl}_2)$ is:
\[
\mathcal{T}(u) = \operatorname{tr}_{\mathbb{C}^2} T(u)
= t_{11}(u) + t_{22}(u).
\]
This is a generating function for commuting elements:
$[\mathcal{T}(u), \mathcal{T}(v)] = 0$ for all $u, v$.

\emph{Bar complex interpretation.}
The commutativity $[\mathcal{T}(u), \mathcal{T}(v)] = 0$ means that
the degree-2 bar element
\[
[\mathcal{T}^{(r)} | \mathcal{T}^{(s)}] \in \bar{B}^2(Y(\mathfrak{sl}_2))
\]
is exact for all $r, s$:
$d[\mathcal{T}^{(r)} | \mathcal{T}^{(s)}] = [\mathcal{T}^{(r)},
\mathcal{T}^{(s)}] = 0$.
Thus the elements $[\mathcal{T}^{(r)} | \mathcal{T}^{(s)}]$
are \emph{bar cocycles}, defining classes in $H^2(\bar{B})$.
These classes are trivial (exact) because
$\mathcal{T}^{(r)} = t_{11}^{(r)} + t_{22}^{(r)}$ is a sum
of generators, and:
\begin{align*}
d([t_{11}^{(r)} | t_{22}^{(s)}] + [t_{22}^{(s)} | t_{11}^{(r)}])
&= [t_{11}^{(r)}, t_{22}^{(s)}] + [t_{22}^{(s)}, t_{11}^{(r)}] = 0.
\end{align*}

The Baxter Q-operator $Q(u)$ satisfies a TQ-relation. In the
additive algebraic convention for the Yang kernel the shifts are
\(\pm1\); after the physical XXX rescaling \(u\mapsto iu\) used
below they are \(\pm i\). In the bar complex, the relation appears
as a degree-2 cocycle condition with a shift in the spectral
parameter filtration.
\end{computation}


\subsection{Classical residue and quantized \texorpdfstring{$R$}{R}-matrix}
\label{subsec:quantum-rmatrix-bar}
\index{R-matrix!quantum!from bar complex}

Two rational kernels enter the affine-to-Yangian comparison. The
ordered affine collision residue is
\(r^{\mathrm{coll}}_k(z)=k\,\Omega_{\mathrm{tr}}/z\). The KZ
connection coefficient is
\(r^{\mathrm{KZ}}_k(z)=\Omega_{\mathrm{KZ}}/((k+h^\vee)z)\). The
Drinfeld--Kohno transport quantizes the KZ coefficient with
\(\hbar_{\mathrm{KZ}}=\pi i/(k+h^\vee)\); the Yangian/FRT kernel is
the rational evaluation form obtained after this transport, not the
affine collision residue with the level prefix erased.

\begin{proposition}[Classical residue from the bar complex and
quantized $R$-matrix after KZ/DJ input; \ClaimStatusConditional]
\label{prop:rmatrix-from-bar}
\index{R-matrix!from bar complex|textbf}
\index{bar complex!R-matrix extraction}
\index{collision residue!classical and quantum $R$-matrix}
For a modular Koszul chiral algebra $\cA = \hat{\fg}_k$ with
$\fg$ simple:
\begin{enumerate}[label=\textup{(\roman*)}]
\item \emph{Classical $r$-matrix surfaces.}
 The genus-$0$ binary collision residue of the bar-complex MC
 element is
 \begin{equation}\label{eq:classical-r-from-bar}
 r^{\mathrm{coll}}_k(z)
 = \operatorname{Res}^{\mathrm{coll}}_{0,2}(\Theta_\cA)
 = \frac{k\,\Omega_{\mathrm{tr}}}{z}
 \;\in\; \fg \otimes \fg \otimes z^{-1}\bC[z^{-1}],
 \end{equation}
 where \(\Omega_{\mathrm{tr}}\) is the Casimir for the normalized
 trace form. The KZ transport surface is instead
 \begin{equation}\label{eq:kz-r-from-bar}
 r^{\mathrm{KZ}}_k(z)
 =\frac{\Omega_{\mathrm{KZ}}}{(k+h^\vee)z},
 \end{equation}
 where \(\Omega_{\mathrm{KZ}}\) denotes the KZ/Sugawara-normalized
 Casimir tensor used in the landscape convention,
 \(\Omega_{\mathrm{KZ}}=2h^\vee\Omega_{\mathrm{tr}}\). This is a
 convention for comparing the trace-form collision tensor with the
 KZ coefficient, not the dual-basis Casimir for the Killing form.
 The denominator \(k+h^\vee\) is the Sugawara denominator from
 dualizing the zeroth product via the level-shifted invariant form
 \textup{(}Computation~\textup{\ref{comp:sl2-collision-residue-kz}}\textup{)}.
 The two expressions are convention-related representatives, not a
 generic equality in \(k\).
 In the fundamental representation of $\mathfrak{sl}_N$:
 $\Omega = P - \mathbb{1}/N$, where $P$ is the permutation
 operator. The classical Yang--Baxter equation
 \begin{equation}\label{eq:cybe-bar}
 [r_{12}(u), r_{13}(u{+}v)] + [r_{12}(u), r_{23}(v)]
 + [r_{13}(u{+}v), r_{23}(v)] = 0
 \end{equation}
 follows from the infinitesimal braid relations for~$\Omega$.

\item \emph{Quantized $R$-matrix after KZ/DJ input.}
 The KZ-transported deformation, with
 $\hbar_{\mathrm{KZ}} = \log q$ and
 $q = e^{\pi i / (k + h^\vee)}$, produces
 the quantum $R$-matrix:
 \begin{equation}\label{eq:quantum-R-from-bar}
 R_q = \mathbb{1} + \hbar_{\mathrm{KZ}}\, r_{\mathrm{DJ}}
 + O(\hbar_{\mathrm{KZ}}^2),
 \end{equation}
 where \(r_{\mathrm{DJ}}\) is the Drinfeld--Jimbo classical
 \(r\)-matrix whose symmetrization is the Casimir appearing in
 \(r^{\mathrm{KZ}}_k\). The higher-order terms are determined
 recursively by the quantum Yang--Baxter equation. The parameter
 \(\hbar_{\mathrm{KZ}}\) is neither the Yangian spectral parameter
 \(u\) nor the bare FRT deformation parameter \(\hbar\); identifying
 deformation parameters is part of the chosen DK/FRT normalization.
 For
 $\mathfrak{sl}_2$ in the fundamental:
 \begin{equation}\label{eq:quantum-R-sl2-explicit}
 R_q = q^{H \otimes H / 4}\,
 \sum_{n \geq 0}
 \frac{q^{n(n-1)/2}\,(q - q^{-1})^n}{[n]_q!}\;
 E^n \otimes F^n,
 \end{equation}
 where $[n]_q = (q^n - q^{-n})/(q - q^{-1})$ and only the
 $n = 0, 1$ terms contribute in the fundamental.

\item \emph{Classical bar MC and quantized QYBE.}
 The classical Yang--Baxter equation for the collision residue is a
 formal consequence of the coassociativity of the bar coproduct
 $\Delta_{\barB}\colon \barB(\cA) \to
 \barB(\cA) \otimes \barB(\cA)$.
 In the semiclassical degree-$3$ projection, coassociativity
 $(\Delta \otimes \mathrm{id}) \circ \Delta
 = (\mathrm{id} \otimes \Delta) \circ \Delta$
 projects to the CYBE~\eqref{eq:cybe-bar}.
 The full quantum Yang--Baxter equation is obtained only after a
 Drinfeld--Jimbo or KZ quantization of this classical tensor has been
 chosen; it is not forced by the classical bar coproduct alone.
\end{enumerate}
\end{proposition}

\begin{proof}
Part~(i): The bar propagator
$\eta_{12} = d\log(z_1 - z_2)$ absorbs one pole order from the
affine OPE. The double pole $k\,g^{ab}/(z{-}w)^2$
shifts to $k\,g^{ab}/(z{-}w)$ after $d\log$ absorption; the simple pole
$f^{ab}{}_c J^c/(z{-}w)$ becomes regular after $d\log$
absorption and drops. The degree-$2$ projection of the MC element
 therefore extracts only the Casimir part:
 \(r^{\mathrm{coll}}_k(z)=k\,\Omega_{\mathrm{tr}}/z\)
 (Proposition~\ref{prop:affine-propagator-matching}). The CYBE
reduces to the infinitesimal braid relations
$[\Omega_{12}, \Omega_{23}] + [\Omega_{13}, \Omega_{23}] = 0$,
which follow from the Jacobi identity in~$\fg$ (the Casimir
$\Omega$ satisfies $[\Omega_{12} + \Omega_{13}, \Omega_{23}] = 0$
because $\Omega$ is $\fg$-invariant: $[\Delta(x), \Omega] = 0$
for all $x \in \fg$). The KZ coefficient
\(\Omega_{\mathrm{KZ}}/((k+h^\vee)z)\) is obtained after the
Sugawara/KZ convention is chosen; Lemma~\ref{lem:g1sf-casimir-bridge}
keeps this convention separate from the trace-form collision
residue.

Part~(ii): The universal $R$-matrix~\eqref{eq:quantum-R-sl2-explicit}
is Drinfeld's~\cite{Drinfeld85} and Jimbo's~\cite{Jimbo85}. In the
spin-$\frac{1}{2}$ representation, $E = e_{12}$, $F = e_{21}$,
$H = e_{11} - e_{22}$, and the sum truncates at $n = 1$ because
$E^2 = F^2 = 0$ on $\bC^2$. The diagonal factor $q^{H \otimes H/4}$
evaluates to $\operatorname{diag}(q^{1/4}, q^{-1/4}, q^{-1/4}, q^{1/4})$,
and the $n = 1$ term contributes $(q - q^{-1}) e_{12} \otimes e_{21}$.
After incorporating the normalization convention, this yields the
matrix of Theorem~\ref{thm:quantum-rmatrix-shadow}. At
\(\hbar_{\mathrm{KZ}} = 0\)
$(q = 1)$: $R_q \to \mathbb{1}$, and the derivative
$dR_q/d\hbar_{\mathrm{KZ}}|_{\hbar_{\mathrm{KZ}}=0}$ recovers
the Drinfeld--Jimbo triangular \(r\)-matrix. Its symmetric part is the
quadratic Casimir used by the KZ connection.

Part~(iii): The bar coalgebra $\barB(\cA)$ is a coassociative
factorization coalgebra
(Theorem~\ref{thm:bar-modular-operad}). The MC element
$\Theta_\cA = D_\cA - d_0$ satisfies
$D_\cA^2 = 0$ (Theorem~\ref{thm:convolution-d-squared-zero}),
so the MC equation $D\Theta + \frac{1}{2}[\Theta, \Theta] = 0$
holds at all degrees. The degree-$3$ projection of this equation is
the CYBE for the degree-$2$ collision residue
(Theorem~\ref{thm:collision-depth-2-ybe}). A quantum
\(R_q\) satisfying QYBE requires the additional quantization input
specified in part~(ii).
\end{proof}

\begin{theorem}[Fundamental quantum \texorpdfstring{$R$}{R}-matrix and classical residue;
\ClaimStatusProvedHere]
\label{thm:quantum-rmatrix-shadow}
\index{R-matrix!quantum!shadow tower origin}
\index{Drinfeld--Kohno!R-matrix quantization}
For the KZ level~$k$ specialization
\(q = e^{\pi i/(k+h^\vee)}\), the Drinfeld--Jimbo quantum group
$U_q(\mathfrak{sl}_2)$ has universal $R$-matrix whose
fundamental representation is
\begin{equation}\label{eq:quantum-R-sl2}
R_q =
\begin{pmatrix}
q & 0 & 0 & 0 \\
0 & 1 & q - q^{-1} & 0 \\
0 & 0 & 1 & 0 \\
0 & 0 & 0 & q
\end{pmatrix}
\end{equation}
in the ordered basis $|{+}{+}\rangle, |{+}{-}\rangle,
|{-}{+}\rangle, |{-}{-}\rangle$, where $|{\pm}\rangle$ is
the spin-$\frac{1}{2}$ basis.

The $\hbar_{\mathrm{DJ}}$-expansion
$R_q = \mathbb{1} + \hbar_{\mathrm{DJ}}\, r + O(\hbar_{\mathrm{DJ}}^2)$
with $\hbar_{\mathrm{DJ}} = \log q$ recovers the FRT-gauge classical matrix
\[
r_{\mathrm{FRT}}
= e_{11}\otimes e_{11}+e_{22}\otimes e_{22}
  +2\,e_{12}\otimes e_{21}.
\]
Its symmetrization is \(2P\), hence
\[
\frac12\bigl(r_{\mathrm{FRT}}+r_{\mathrm{FRT}}^{21}\bigr)
-\frac12\,\mathbb{1}
=P-\frac12\,\mathbb{1}
=\Omega_{\mathfrak{sl}_2}^{\mathrm{fund}}.
\]
The KZ monodromy coefficient is therefore the Casimir surface
\(\Omega_{\mathrm{KZ}}/((k+2)z)\), whereas the ordered affine
collision residue remains \(k\,\Omega_{\mathrm{tr}}/z\).
\end{theorem}

\begin{proof}
The universal $R$-matrix of $U_q(\mathfrak{sl}_2)$ is
$\mathcal{R} = q^{H \otimes H/4}
\sum_{n \geq 0} \frac{q^{n(n-1)/2}(q - q^{-1})^n}{[n]_q!}
E^n \otimes F^n$.
In the spin-$\frac{1}{2}$ representation ($E = e_{12}$,
$F = e_{21}$, $H = e_{11} - e_{22}$), only the $n = 0$ and
$n = 1$ terms contribute. The $n = 0$ term gives the diagonal
$q^{H \otimes H/4}$:
\[
q^{H \otimes H / 4}
= \operatorname{diag}(q^{1/4}, q^{-1/4}, q^{-1/4}, q^{1/4})
\cdot q^{3/4} \cdot (\ldots)
\]
After multiplication by the standard scalar and diagonal FRT gauge,
this gives~\eqref{eq:quantum-R-sl2}.

The $\hbar_{\mathrm{DJ}}$-expansion:
\[
R_q=\mathbb{1}
+\hbar_{\mathrm{DJ}}\,(e_{11}\otimes e_{11}+e_{22}\otimes e_{22}
       +2e_{12}\otimes e_{21})
+O(\hbar_{\mathrm{DJ}}^2),
\]
where \(e_{ij}\) are the \(2\times2\) matrix units. This is the
FRT-gauge classical \(r\)-matrix. Passing to the unitarized
Drinfeld--Kohno gauge replaces it by the Casimir representative in
the KZ connection; it does not identify the affine collision
coefficient \(k\Omega_{\mathrm{tr}}\) with
\(\Omega_{\mathrm{KZ}}/(k+2)\).
\end{proof}

\begin{proposition}[Colored $R$-matrices and Casimir eigenvalues;
\ClaimStatusProvedElsewhere]
\label{prop:colored-rmatrix}
\index{R-matrix!colored}
For irreducible representations $V_\lambda$, $V_\mu$ of
$U_q(\mathfrak{sl}_2)$, with \(q=e^{\pi i/(k+2)}\) when the
representation is obtained from KZ level~\(k\), the $R$-matrix
$\check R_{V_\lambda, V_\mu}=P\circ R_{V_\lambda,V_\mu}$ acts on the
irreducible component
$V_\nu \subset V_\lambda \otimes V_\mu$ with eigenvalue
\begin{equation}\label{eq:colored-eigenvalue}
\check R\big|_{V_\nu}
= \epsilon_{\lambda,\mu}^{\nu}\,
q^{(c_\nu-c_\lambda-c_\mu)/2}
\end{equation}
in the ribbon normalization, where
\(c_\lambda=(\lambda,\lambda+2\rho)\) and
\(\epsilon_{\lambda,\mu}^{\nu}\in\{\pm1\}\) is the symmetry sign of
the summand.
For $V = V_1$ \textup{(}fundamental\textup{)}, the decomposition
$V_1 \otimes V_1 = V_0 \oplus V_2$ gives, with
\((\alpha,\alpha)=2\), the ribbon-gauge eigenvalues
\(-q^{-3/2}\) on the antisymmetric singlet and \(q^{1/2}\) on the
symmetric triplet, up to an overall scalar gauge. The FRT matrix
\eqref{eq:quantum-R-sl2} is one such gauge representative.
\end{proposition}

\begin{proof}
The eigenvalue formula~\eqref{eq:colored-eigenvalue} is the standard
ribbon-category normalization of the Drinfeld--Jimbo braiding. It
follows from the fact that the universal $R$-matrix satisfies
$(\Delta \otimes \operatorname{id})(\mathcal{R})
= \mathcal{R}_{13}\mathcal{R}_{23}$
and the Casimir element $C = \sum_a T^a T_a$ is a central element
of $U_q(\fg)$ with eigenvalue \(c_\lambda\) on $V_\lambda$; the
categorical symmetry supplies the sign on the antisymmetric summand.
\end{proof}


\subsection{Bethe ansatz from the shadow connection}
\label{subsec:bethe-ansatz-shadow}
\index{Bethe ansatz!from shadow tower}

The shadow connection $\nabla^{\mathrm{sh}}$ on the modular
convolution algebra, restricted to genus~$0$ and degree~$n$, gives the
KZ connection on $\operatorname{Conf}_n(\bP^1)$ through its binary
collision residue. Bethe equations are not a formal consequence of the
KZ residue alone: they enter after one chooses the corresponding
Yangian or quantum-loop transfer matrix and applies the algebraic
Bethe ansatz.

\begin{theorem}[Bethe ansatz from the MC equation;
\ClaimStatusConditional]
\label{thm:bethe-from-mc}
\index{Bethe ansatz!XXX chain}
\index{Heisenberg chain!from bar complex}
For the antiferromagnetic Heisenberg XXX chain of length~$L$ with
$M$ down spins, the Bethe ansatz equations
\begin{equation}\label{eq:xxx-bae}
\prod_{j \neq i} \frac{u_i - u_j + i}{u_i - u_j - i}
= \left(\frac{u_i + \frac{i}{2}}{u_i - \frac{i}{2}}\right)^{\!L},
\qquad i = 1, \ldots, M
\end{equation}
arise from the shadow connection after this transfer-matrix choice as
follows. The rational
$R$-matrix \(R_{\mathrm{XXX}}(u)=u\,\mathbb{1}+iP\), equivalently
Yang's additive \(u\,\mathbb{1}+P\) after \(u\mapsto iu\),
\textup{(}Theorem~\textup{\ref{thm:rtt-all-classical-types}(i))}
defines the transfer matrix
$\mathcal{T}(u) = \operatorname{tr}_{\mathrm{aux}}
\prod_{j=1}^{L} R_{0j}(u)$,
whose eigenvalue conditions are the
BAE~\eqref{eq:xxx-bae}.

The energy eigenvalues are
\begin{equation}\label{eq:xxx-energy}
E = \frac{L}{4} - \frac{1}{2} \sum_{a=1}^{M}
 \frac{1}{u_a^2 + \frac{1}{4}},
\end{equation}
where $\{u_a\}$ are the Bethe roots solving~\eqref{eq:xxx-bae}.
\end{theorem}

\begin{proof}
The transfer matrix $\mathcal{T}(u)$ commutes for distinct
spectral parameters: $[\mathcal{T}(u), \mathcal{T}(v)] = 0$
(a consequence of the YBE, verified in
Computation~\ref{comp:transfer-matrix-bar}). The algebraic Bethe
ansatz \textup{(}Faddeev~\cite{Faddeev96}\textup{)} constructs
simultaneous eigenvectors by applying $B(u_1) \cdots B(u_M)$ to
the pseudovacuum $|\Omega\rangle = |{\uparrow \cdots \uparrow}\rangle$,
where $B(u)$ is the off-diagonal entry of the monodromy matrix
$T(u) = \bigl[\begin{smallmatrix} A & B \\ C & D \end{smallmatrix}\bigr]$.
The eigenvector condition $\mathcal{T}(u)\,|u_1, \ldots, u_M\rangle
= \Lambda(u)\,|u_1, \ldots, u_M\rangle$ is satisfied if and only if
the Bethe roots $\{u_a\}$ solve~\eqref{eq:xxx-bae}.

The energy formula follows from
$H_{\mathrm{XXX}} = \frac{J}{2}
\frac{d}{du}\log \mathcal{T}(u)\big|_{u=\frac{i}{2}}$
(Sutherland's formula), giving~\eqref{eq:xxx-energy}.

Numerical verification for $L = 4, 6, 8$: the Bethe-ansatz energies
agree with exact diagonalization to machine precision ($< 10^{-10}$);
see \texttt{bethe\_ansatz\_shadow.py}.
\end{proof}

\begin{computation}[Bethe ansatz vs exact diagonalization]
\label{comp:bethe-exact}
\index{Bethe ansatz!numerical verification}
For the XXX chain at $L = 6$, half-filling ($M = 3$), the
ground-state Bethe roots are $u_1 \approx -0.6557$,
$u_2 = 0$, $u_3 \approx 0.6557$, yielding ground-state energy
$E_0 \approx -1.2019$ \textup{(}in units with $J = 1$\textup{)},
matching exact diagonalization of the $64 \times 64$
Hamiltonian matrix in the $S^z = 0$ sector.
The Bethe roots satisfy the string
hypothesis at large~$L$: the ground-state roots distribute
symmetrically about $u = 0$ with spacing $\sim 1/L$.

The TQ relation
\begin{equation}\label{eq:tq-relation}
\mathcal{T}(u)\, Q(u)
= a(u)\, Q(u - i) + d(u)\, Q(u + i),
\end{equation}
where $a(u) = (u + \frac{i}{2})^L$, $d(u) = (u - \frac{i}{2})^L$,
and $Q(u) = \prod_{a=1}^M (u - u_a)$ is the Baxter $Q$-operator,
encodes the BAE as the regularity condition for the eigenvalue
$\Lambda(u) = a(u) Q(u-i)/Q(u) + d(u) Q(u+i)/Q(u)$.
The poles of the right-hand side at $u = u_a$ cancel if and only
if~\eqref{eq:xxx-bae} holds.
\end{computation}

\begin{remark}[Ordered amplitudes and the symmetric Bethe shadow]
\label{rem:yangcomp-ordered-bethe-shadow}
\index{Bethe ansatz!ordered amplitudes}
\index{averaging map!Bethe shadow}
For a length-$L$ spin chain on $V=\mathbb{C}^2$, the ordered
genus-$0$ chiral datum carries $V^{\otimes L}$ with its
pure-braid KZ monodromy. The averaging map
\[
\av_L\colon V^{\otimes L}\longrightarrow \operatorname{Sym}^L(V)
\]
forgets ordered site amplitudes and retains the symmetric spectral
data: the transfer eigenvalue $\Lambda(u)$, the Baxter polynomial
$Q(u)=\prod_a(u-u_a)$, and the quantum determinant. Thus
\[
\dim\ker(\av_L)=2^L-(L+1),
\]
but $\ker(\av_L)$ is not a Bethe eigenspace. Bethe vectors
$B(u_1)\cdots B(u_M)|\Omega\rangle$ live in the ordered Hilbert space
$V^{\otimes L}$; under the usual completeness hypotheses they span
the appropriate magnon-sector eigenspaces. Only differences of
ordered vectors with the same symmetric spectral data are killed by
averaging. This is the ordered/symmetric distinction behind
Remark~\ref{rem:ordered-center-bethe}.
\end{remark}

\begin{remark}[Nested Bethe ansatz for higher rank]
\label{rem:nested-bethe}
\index{Bethe ansatz!nested!sl3@$\mathfrak{sl}_3$}
For $\mathfrak{sl}_3$, the Bethe ansatz involves two species of
Bethe roots $(u_a^{(1)}, u_b^{(2)})$ with coupled equations:
\begin{align*}
\prod_{j \neq i} \frac{u_i^{(1)} - u_j^{(1)} + i}
{u_i^{(1)} - u_j^{(1)} - i}
&= \left(\frac{u_i^{(1)} + \frac{i}{2}}{u_i^{(1)} - \frac{i}{2}}\right)^{\!L}
\prod_{b} \frac{u_i^{(1)} - u_b^{(2)} - \frac{i}{2}}
{u_i^{(1)} - u_b^{(2)} + \frac{i}{2}}, \\
\prod_{j \neq i} \frac{u_i^{(2)} - u_j^{(2)} + i}
{u_i^{(2)} - u_j^{(2)} - i}
&= \prod_{a} \frac{u_i^{(2)} - u_a^{(1)} - \frac{i}{2}}
{u_i^{(2)} - u_a^{(1)} + \frac{i}{2}}.
\end{align*}
The nesting reflects the rank-$2$ structure of $\mathfrak{sl}_3$:
each simple root contributes one species of Bethe roots, and the
Cartan matrix appears in the coupling. The classical-type
$R$-matrices of Theorem~\ref{thm:rtt-all-classical-types} produce
analogous nested systems for types $B$, $C$, $D$, with the Cartan
matrix of the corresponding Lie algebra governing the coupling
structure
\textup{(}\texttt{bethe\_ansatz\_shadow.py}\textup{)}.
\end{remark}

\begin{theorem}[Bethe ansatz equations from the shadow potential;
\ClaimStatusConditional]
\label{thm:bae-from-mc}
\index{Bethe ansatz!from MC element|textbf}
\index{shadow potential!Bethe ansatz}
\index{Nekrasov--Shatashvili!gauge/Bethe correspondence}
The Bethe ansatz equations for the standard rational, trigonometric,
elliptic, and nested examples below arise as stationarity conditions
for the shadow potential
$\cW(\{u_i\})$, which is the genus-$0$ shadow amplitude
evaluated on the spin chain configuration space.
\begin{enumerate}[label=\textup{(\roman*)}]
\item \emph{XXX chain \textup{(}rational, from $\hat{\mathfrak{sl}}_2$\textup{)}.}
	 The shadow potential is
	 \begin{equation}\label{eq:xxx-shadow-potential}
	 \cW_{\mathrm{XXX}}(\{u_i\})
	 = L\sum_{i=1}^{M}
	 \bigl(\Phi(u_i+\tfrac{i}{2})-\Phi(u_i-\tfrac{i}{2})\bigr)
	 -\frac12\sum_{i\ne j}
	 \bigl(\Phi(u_i-u_j+i)-\Phi(u_i-u_j-i)\bigr),
	 \end{equation}
	 where $\Phi(u) = u\log u - u$ and a branch of \(\log\) has been
	 chosen. The BAE~\eqref{eq:xxx-bae} are the exponential form of
	 \(\partial \cW/\partial u_i=0\). The shadow connection
 $\nabla^{\mathrm{sh}}$
 \textup{(}Theorem~\textup{\ref{thm:shadow-connection})}
 restricted to genus~$0$, degree~$L$, produces the KZ connection
 on the configuration space $\operatorname{Conf}_L(\bP^1)$;
 the algebraic Bethe ansatz is the transfer-matrix quantization of
 this KZ Hamiltonian system, not a consequence of the collision
 residue alone.

\item \emph{XXZ chain \textup{(}trigonometric, from $U_q(\hat{\mathfrak{sl}}_2)$\textup{)}.}
 Setting $q = e^\eta$ with anisotropy $\Delta = \cosh \eta$,
 the XXZ Bethe equations are
 \begin{equation}\label{eq:xxz-bae}
 \prod_{j \neq i}
 \frac{\sinh(u_i - u_j + \eta)}{\sinh(u_i - u_j - \eta)}
 = \left(
 \frac{\sinh(u_i + \eta/2)}{\sinh(u_i - \eta/2)}
 \right)^{\!L},
 \quad i = 1, \ldots, M.
 \end{equation}
 These arise from the genus-$0$ shadow obstruction tower on
 the cylinder $\bC^* = \bC / \bZ$: the bar propagator becomes
 $d\log(e^{z_1} - e^{z_2})$, and the collision residue extracts
 the trigonometric $r$-matrix
 $r^{\mathrm{trig}}(u) = \Omega\,\coth(u)$.

\item \emph{XYZ chain \textup{(}elliptic, from elliptic quantum group\textup{)}.}
 On an elliptic curve $E_\tau = \bC / (\bZ + \bZ\tau)$, the bar
 propagator $d\log E(z,w) = \zeta_\tau(z - w)\,d(z-w)$ produces
 the elliptic $r$-matrix
 $r^{\mathrm{ell}}(u, \tau)$
 \textup{(}Proposition~\textup{\ref{prop:elliptic-rmatrix-shadow}}\textup{)},
 and the BAE involve Jacobi theta functions. The XYZ chain
 has Hamiltonian
 $H = \sum_j (J_x \sigma_j^x \sigma_{j+1}^x
 + J_y \sigma_j^y \sigma_{j+1}^y
 + J_z \sigma_j^z \sigma_{j+1}^z)$
 with $(J_x, J_y, J_z)$ parametrized by $(k', k, 1)$
 in terms of the elliptic modulus.

\item \emph{Higher rank.}
 For $\mathfrak{sl}_N$, the nested Bethe ansatz involves
 $N - 1$ species of Bethe roots $\{u_a^{(\alpha)}\}$
 ($\alpha = 1, \ldots, N{-}1$) with coupled equations
 governed by the Cartan matrix of $\mathfrak{sl}_N$:
 \begin{equation}\label{eq:nested-bae}
 \prod_{j \neq i}
 \frac{u_i^{(\alpha)} - u_j^{(\alpha)} + i}
 {u_i^{(\alpha)} - u_j^{(\alpha)} - i}
 = \prod_{\beta \sim \alpha}\,
 \prod_{b}
 \frac{u_i^{(\alpha)} - u_b^{(\beta)} - iA_{\alpha\beta}/2}
 {u_i^{(\alpha)} - u_b^{(\beta)} + iA_{\alpha\beta}/2},
 \end{equation}
 where \(A_{\alpha\beta}=-C_{\alpha\beta}>0\) for adjacent nodes
 and $\beta \sim \alpha$ denotes adjacency in the Dynkin diagram.
\end{enumerate}
\end{theorem}

\begin{proof}
Part~(i) follows from Theorem~\ref{thm:bethe-from-mc} by
rewriting the BAE in logarithmic form:
$L\log((u_i + i/2)/(u_i - i/2))
= \sum_{j \neq i} \log((u_i - u_j + i)/(u_i - u_j - i))$.
Define $\cW = L \sum_i \mathrm{Li}(u_i) - \sum_{i<j} \mathrm{Li}(u_i - u_j)$
where $\mathrm{Li}$ is the appropriate logarithmic primitive;
then $e^{\partial \cW/\partial u_i} = 1$ is exactly the BAE.
The identification with the shadow connection uses the fact that
the genus-$0$ degree-$n$ shadow amplitude
$\mathrm{Sh}_{0,n}(\Theta_\cA)$ on $\operatorname{Conf}_n(\bP^1)$
reduces to the KZ Hamiltonian
$\sum_{i \neq j} \Omega_{ij}/(z_i - z_j)$
by Proposition~\ref{prop:kz-from-shadow} below;
the Bethe ansatz diagonalizes this Hamiltonian via the
algebraic Bethe ansatz of Faddeev~\cite{Faddeev96}.

Part~(ii): On the cylinder, the lattice $\bZ$ periodicity converts
$1/z$ to $\cot(z)$ (or $\coth(z)$ after Wick rotation), giving the
trigonometric $r$-matrix. The XXZ BAE follow from the transfer matrix
with the trigonometric $R$-matrix $R(u) = \sinh(u + \eta)\,\mathbb{1}
+ \sinh(\eta)\,P/\sinh(u)$ by the same algebraic Bethe ansatz procedure.

Part~(iii): On the elliptic curve $E_\tau$, the bar propagator
$d\log E(z,w)$ extracts the Weierstrass $\zeta$-function instead of
$1/z$, giving $r^{\mathrm{ell}}(z,\tau)$ with the elliptic
deformation of the pole structure. The XYZ chain's
integrability via the elliptic $R$-matrix is
Baxter's~\cite{Baxter82}.

Part~(iv): The nested Bethe ansatz for $\mathfrak{sl}_N$ follows
from the classical-type RTT presentation
(Theorem~\ref{thm:rtt-all-classical-types}): the $R$-matrix for
$\mathfrak{sl}_N$ in the fundamental defines a transfer matrix,
and the successive ``auxiliary spaces'' for each simple root yield
the $N - 1$ species of Bethe roots. The Cartan-matrix coupling
arises from the overlap of adjacent root systems in the nesting.
Numerical verification for $\mathfrak{sl}_3$ at $L = 4$ confirms
agreement between nested Bethe ansatz energies and exact
diagonalization to precision $< 10^{-10}$
(\texttt{bethe\_ansatz\_shadow.py}).
\end{proof}


\section{Yangian representation theory and bar modules}
\label{sec:yangian-rep-bar}
\index{Yangian!representations}

\begin{definition}[Evaluation module]\label{def:evaluation-module-yangian}
\index{evaluation module|textbf}
For $a \in \mathbb{C}$ and $V$ a finite-dimensional
$\mathfrak{sl}_N$-representation, the \emph{evaluation module}
$V(a)$ of $Y_\hbar(\mathfrak{sl}_N)$ is the pullback along the
evaluation homomorphism. In the type-$A$ RTT window this homomorphism is
\[
\operatorname{ev}_a(T_{ij}(u))
 =
\delta_{ij}+\frac{\hbar\,\rho(E_{ij})}{u-a},
\qquad
T_{ij}^{(r)}\longmapsto a^{r-1}\rho(E_{ij})
\quad (r\ge1),
\]
after the usual trace/quantum-determinant reduction from
$\mathfrak{gl}_N$ to $\mathfrak{sl}_N$. The parameter $a$ is the
additive spectral point. It is not an affine Kac--Moody level.
\end{definition}

\begin{proposition}[Evaluation quotient bar complex; \ClaimStatusProvedHere]
\label{prop:eval-module-bar}

Fix \(\hbar\ne0\), and let $I_a=\ker(\operatorname{ev}_a)$. The
evaluation quotient $Y_\hbar(\mathfrak{sl}_N)/I_a$ is
$U(\mathfrak{sl}_N)$ on the finite evaluation lane. Hence the
relative bar complex of the evaluation quotient acting on $V(a)$ is
the Chevalley--Eilenberg complex of $\mathfrak{sl}_N$ with
coefficients in~$V$:
\[
\bar{B}\bigl(Y_\hbar(\mathfrak{sl}_N)/I_a;\,V(a)\bigr)
\simeq C^*(\mathfrak{sl}_N, V)
\]
as a chain complex. The full $Y_\hbar$-module bar complex maps to this
quotient complex; the kernel is generated by the evaluation ideal
$I_a$ and belongs to the completed Yangian bar surface, not to the
finite evaluation quotient.
\end{proposition}

\begin{proof}
Expanding
\[
\frac{1}{u-a}=\sum_{r\ge1}a^{r-1}u^{-r}
\]
shows that $\operatorname{ev}_a$ identifies every RTT mode with the
same $\mathfrak{sl}_N$ action multiplied by the scalar $a^{r-1}$.
Modulo $I_a$ no independent higher-mode generator remains, and the
quotient algebra is $U(\mathfrak{sl}_N)$ with its usual action on~$V$,
after the harmless rescaling by the non-zero deformation parameter
\(\hbar\). The bar complex of this quotient module is therefore the
standard Chevalley--Eilenberg complex \(C^*(\mathfrak{sl}_N,V)\), with
differential determined by the $\mathfrak{sl}_N$-action and the Lie
bracket. The evaluation point remains visible only in the embedding of
the quotient into the RTT mode tower; it does not create an affine
level.
\end{proof}

\begin{remark}[Tensor product of evaluation modules]
\label{rem:tensor-eval}
\index{evaluation module!tensor product}

For two evaluation modules $V(a)$ and $W(b)$, the coproduct makes
$V(a)\otimes W(b)$ a Yangian module for all $a,b$. When $a\neq b$, the
universal rational $R$-intertwiner is regular, and the bar complex of
the tensor product is related to the bar complexes of the factors by a
K\"unneth-type formula on the evaluation quotient:
\[
\bar{B}(Y; V(a) \otimes W(b)) \simeq
\bar{B}(Y; V(a)) \otimes_{\bar{B}(Y)} \bar{B}(Y; W(b))
\]
where the tensor product is taken over the bar construction.

When $a=b$ (coincident evaluation points), the tensor product remains
defined, but the normalized $R$-intertwiner has the collision pole
$R_\hbar(a-b)=1-\hbar P/(a-b)$. The pole belongs to the ordered
configuration divisor. The FM blow-up at the diagonal records the
collision direction and produces the propagator used by the Yangian
bar complex.
\end{remark}

\subsection{\texorpdfstring{Yangian category $\mathcal{O}$ from bar-cobar duality}{Yangian category O from bar-cobar duality}}
\label{sec:yangian-category-O}
\index{Yangian!category O}

We first fix the ambient category.

\begin{definition}[Yangian category \texorpdfstring{$\mathcal{O}$}{O}]
\label{def:yangian-category-O}
\index{Yangian!category O!definition}
For $\fg = \mathfrak{sl}_N$, the \emph{Yangian category~$\mathcal{O}$}
$\mathcal{O}_{Y_\hbar(\mathfrak{sl}_N)}$ is the full subcategory of
$Y_\hbar(\mathfrak{sl}_N)$-modules~$M$ satisfying:
\begin{enumerate}[label=\textup{(\roman*)}]
\item $M$ is finitely generated over~$Y_\hbar(\mathfrak{sl}_N)$;
\item $M$ has a weight decomposition
 $M = \bigoplus_{\mu \in \mathfrak{h}^*} M_\mu$ with
 finite-dimensional weight spaces;
\item the set of weights of~$M$ is contained in a finite union
 of cones $\{\mu \leq \lambda_i\}$ for the dominance order.
\end{enumerate}
The \emph{polynomial subcategory}
$\mathcal{O}_{\mathrm{poly}} \subset \mathcal{O}_{Y_\hbar}$ is the
full subcategory of modules whose composition factors all have
polynomial highest $\ell$-weights in the sense of
Drinfeld~\cite{Drinfeld88}. For $N = 2$, every simple in
$\mathcal{O}$ is finite-dimensional, so
$\mathcal{O}_{\mathrm{poly}} = \mathcal{O}_{Y_\hbar}$. For
$N \ge 3$, it is an open question whether
$\mathcal{O}_{\mathrm{poly}}$ exhausts $\mathcal{O}_{Y_\hbar}$
\textup{(}Remark~\textup{\ref{rem:catO-thick-generation-verification}}\textup{)}.
\end{definition}

The evaluation module bar complex
\textup{(}Proposition~\ref{prop:eval-module-bar}\textup{)} yields a
BGG-type resolution theory on the intrinsic bar-comodule surface. On
lanes where the corresponding finite-type dualized Yangian-module
interpretation is available, the same calculation may be read on the
dual Yangian side.

\begin{theorem}[Evaluation-quotient BGG resolution; \ClaimStatusConditional]
\label{thm:yangian-bgg}
\index{BGG resolution!Yangian}
For $\mathfrak{g} = \mathfrak{sl}_N$ and a finite-dimensional
irreducible $\mathfrak{g}$-module $L(\lambda)$, the evaluation
module $L(\lambda)(a)$ admits the following relative bar resolution
over the evaluation quotient \(Y_\hbar(\mathfrak{sl}_N)/I_a\):
\begin{equation}\label{eq:yangian-bgg}
\cdots \to \bar{B}_2(Y/I_a; L(\lambda)(a))
 \to \bar{B}_1(Y/I_a; L(\lambda)(a))
 \to L(\lambda)(a) \to 0
\end{equation}
where $\bar{B}_n(Y/I_a; L(\lambda)(a))
\simeq C^n(\mathfrak{sl}_N, L(\lambda))$
is the Chevalley--Eilenberg complex. The absolute
\(Y_\hbar\)-module bar complex maps to this quotient complex; no
identification with the Chevalley--Eilenberg complex is asserted
before quotienting by \(I_a\). In particular:
\begin{enumerate}[label=\textup{(\roman*)}]
\item The resolution length is $\dim \mathfrak{sl}_N
 = N^2 - 1$.
\item The terms are
 $\bar{B}_n \cong L(\lambda) \otimes
 \bigwedge^n \mathfrak{sl}_N^*$,
 with differentials given by the dual of the
 $\mathfrak{sl}_N$-action.
\item Under the module bar functor~$\Phi$, the resolution maps to the
 corresponding bar-comodule resolution on the dual bar side. On any
 finite-type dualized evaluation lane where this bar image is
 identified with a \(Y_{-\hbar}(\mathfrak{sl}_N)\)-module, this may
 be read as the dual BGG resolution with inverse RTT kernel. A
 simultaneous reflection \(a\mapsto -a\) is an extra choice of
 spectral coordinate, not part of the quotient calculation itself.
\end{enumerate}
\end{theorem}

\begin{proof}
Parts~(i) and~(ii) follow immediately from
Proposition~\ref{prop:eval-module-bar}: the bar complex
is the CE complex $C^*(\mathfrak{sl}_N, L(\lambda))$,
which has terms $L(\lambda) \otimes \bigwedge^n \mathfrak{sl}_N^*$
and length $\dim \mathfrak{sl}_N$.

For part~(iii): the module bar functor sends the relative
evaluation-quotient resolution~\eqref{eq:yangian-bgg} termwise to the
corresponding bar-comodule resolution on the dual bar side, and
Proposition~\ref{prop:yangian-module-koszul} identifies the
underlying evaluation object with the dual inverse-kernel datum on the
lanes where that Yangian-module realization is available. The
Chevalley--Eilenberg differential depends only on the
$\mathfrak{sl}_N$-action, so the quotient-level bar-side resolution has
the same chain-complex shape. The higher RTT modes remain in the
kernel of the map from the absolute Yangian bar complex to this
relative quotient complex.
\end{proof}

\begin{corollary}[Bar-comodule Ext comparison for Yangian evaluation modules;
\ClaimStatusConditional]
\label{cor:yangian-ext-exchange}
\index{Ext!Yangian modules}
For evaluation modules $L(\lambda)(a)$ and $L(\mu)(b)$
of $Y(\mathfrak{sl}_N)$,
\[
\operatorname{Ext}^n_{Y_\hbar}(L(\lambda)(a), L(\mu)(b))
\;\cong\;
\operatorname{Hom}_{D(\operatorname{CoMod}^{\mathrm{conil}}_{\barB(Y_\hbar)})}
\bigl(
 \Phi(L(\lambda)(a)),\,
 \Phi(L(\mu)(b))[n]
\bigr)
\]
via the module bar functor~$\Phi$. On any finite-type dualized
evaluation lane where the target bar-comodules are identified with
$Y_{-\hbar}(\mathfrak{sl}_N)$-modules, this yields the corresponding
Ext comparison on the dual Yangian module side.
\end{corollary}

\begin{proof}
By Proposition~\ref{prop:ext-tor-exchange},
\[
\operatorname{Ext}^n_{\cA}(L, M)
\cong
\operatorname{Hom}_{D(\operatorname{CoMod}^{\mathrm{conil}}_{\barB(\cA)})}
(\Phi(L), \Phi(M)[n]).
\]
This is exactly the Yangian evaluation-module instance of
Proposition~\ref{prop:ext-tor-exchange}. Proposition~\ref{prop:yangian-module-koszul}
identifies the resulting bar image with the $(-\hbar,-a)$ dual datum
on the evaluation lane, which is why the last sentence may be read on
the dual Yangian module side when the required finite-type dualized
realization is available.
\end{proof}

\begin{remark}[Conditional extension to full Yangian
category \texorpdfstring{$\mathcal{O}$}{O}]
\label{rem:yangian-cat-O-conditional}
The results above hold unconditionally for
evaluation modules (which reduce to finite-dimensional
$\mathfrak{sl}_N$-representations). This is the endpoint of DK-1 in
the Yangian ladder. Extending to
the \emph{full} Yangian category~$\mathcal{O}$ (which
includes infinite-dimensional highest-weight modules,
Verma modules $M(\lambda, a)$, and their subquotients) requires
the Koszulness conjecture
\textup{(}Remark~\ref{rem:yangian-collapse-conj}\textup{)}.
Even if $Y(\mathfrak{g})$ is Koszul, however, one still needs a
separate ambient extension theorem realizing the larger bar-comodule
comparison on the $Y_{-\hbar}$ module surface and respecting the
highest-weight structures of category~$\mathcal{O}$. The remaining
step is therefore not a new local computation on evaluation modules,
but it is more than generation alone.
\end{remark}

\begin{proposition}[DK-2 reduction to thick generation, conditional on an ambient exact extension;
\ClaimStatusConditional]
\label{prop:yangian-dk2-thick-generation}
Let
\[
\mathcal{D}_{\hbar}^{\mathrm{eval}}
\subset D^b(\mathcal{O}_{Y_\hbar}),
\qquad
\mathcal{D}_{-\hbar}^{\mathrm{eval}}
\subset D^b(\mathcal{O}_{Y_{-\hbar}})
\]
be the smallest thick triangulated subcategories containing all
evaluation modules. Assume there exist exact functors
\[
F\colon D^b(\mathcal{O}_{Y_\hbar}) \to D^b(\mathcal{O}_{Y_{-\hbar}}),
\qquad
G\colon D^b(\mathcal{O}_{Y_{-\hbar}}) \to D^b(\mathcal{O}_{Y_\hbar})
\]
whose restrictions on evaluation modules agree with the evaluation
correspondence of Proposition~\ref{prop:yangian-module-koszul} and are
mutually quasi-inverse there. Then:
\begin{enumerate}[label=\textup{(\roman*)}]
\item the functors $F$ and $G$ restrict to quasi-inverse exact
 equivalences
 \[
 F_{\mathrm{eval,thick}}\colon
 \mathcal{D}_{\hbar}^{\mathrm{eval}}
 \xrightarrow{\sim}
 \mathcal{D}_{-\hbar}^{\mathrm{eval}},
 \qquad
 G_{\mathrm{eval,thick}}\colon
 \mathcal{D}_{-\hbar}^{\mathrm{eval}}
 \xrightarrow{\sim}
 \mathcal{D}_{\hbar}^{\mathrm{eval}};
 \]
\item if
 \[
 \mathcal{D}_{\hbar}^{\mathrm{eval}}
 = D^b(\mathcal{O}_{Y_\hbar}),
 \qquad
 \mathcal{D}_{-\hbar}^{\mathrm{eval}}
 = D^b(\mathcal{O}_{Y_{-\hbar}}),
 \]
 then the DK-2 extension from evaluation modules to the full
 completed Yangian category~$\mathcal{O}$ holds.
\end{enumerate}
In particular, once an ambient exact extension of the evaluation-locus
duality is granted, DK-2 reduces to thick generation of evaluation
modules, not to a new local bar computation.
\end{proposition}

\begin{proof}
Exact functors preserve shifts, cones, and direct summands, so both
$F$ and $G$ carry thick subcategories generated by evaluation modules
to the corresponding thick subcategories on the opposite side.
Because $F$ and $G$ are mutually quasi-inverse on evaluation modules,
the full subcategories on which the unit
$\mathrm{id}\to G\!\circ F$ and counit $F\!\circ G\to\mathrm{id}$ are
isomorphisms are thick and contain the evaluation generators.
Therefore these unit and counit maps are isomorphisms on
$\mathcal{D}_{\hbar}^{\mathrm{eval}}$ and
$\mathcal{D}_{-\hbar}^{\mathrm{eval}}$, proving~\textup{(i)}.
Statement~\textup{(ii)} is then immediate.
\end{proof}

\begin{proposition}[Thick generation by evaluation modules
in type~\texorpdfstring{$A$}{A}; \ClaimStatusProvedHere]
\label{prop:dk2-thick-generation-typeA}
\index{Yangian!thick generation}
For $\fg = \mathfrak{sl}_N$ with $N \ge 2$, every finite-dimensional
irreducible $Y(\mathfrak{sl}_N)$-module is a subquotient of a tensor
product of fundamental evaluation modules. Consequently, the
thick triangulated subcategory
\[
\mathcal{D}_{\hbar}^{\mathrm{eval}}
\;\subset\;
D^b\bigl(\mathrm{Rep}_{\mathrm{fd}}(Y_\hbar(\mathfrak{sl}_N))\bigr)
\]
generated by evaluation modules equals
$D^b(\mathrm{Rep}_{\mathrm{fd}}(Y_\hbar(\mathfrak{sl}_N)))$.
\end{proposition}

\begin{proof}
We argue in three steps.

\emph{Step~1: Every simple is a subquotient of evaluation tensor
products.}
Let $L$ be a finite-dimensional irreducible
$Y(\mathfrak{sl}_N)$-module. By the Drinfeld classification
\cite{Drinfeld88}, $L$ is determined by an $N$-tuple of monic
polynomials $(P_1(u), \ldots, P_{N-1}(u))$ (the Drinfeld
polynomials). Write $P_i(u) = \prod_{j=1}^{d_i}(u - a_{ij})$.
The tensor product
\[
\bigotimes_{i=1}^{N-1}\; \bigotimes_{j=1}^{d_i}\;
V_{\omega_i}(a_{ij})
\]
of fundamental evaluation modules
$V_{\omega_i}(a) = \bigwedge^i \mathbb{C}^N(a)$ is a finite-dimensional
$Y(\mathfrak{sl}_N)$-module whose unique irreducible quotient (in a suitable
ordering of factors) has Drinfeld polynomials $(P_1, \ldots, P_{N-1})$
\cite{Drinfeld88}. Hence $L$ is a quotient, and in particular a
subquotient, of a tensor product of fundamental evaluation modules.

\emph{Step~2: Every simple lies in the thick closure.}
Let $M = \bigotimes V_{\omega_i}(a_{ij})$ be the tensor product
from Step~1. Since $M$ is finite-dimensional, it has a
finite composition series $0 = M_0 \subset M_1 \subset \cdots
\subset M_r = M$ with simple graded pieces
$M_k/M_{k-1} \cong L_k$. Each short exact sequence
$0 \to M_{k-1} \to M_k \to L_k \to 0$ yields an exact triangle
$M_{k-1} \to M_k \to L_k \to M_{k-1}[1]$ in
$D^b(\mathrm{Rep}_{\mathrm{fd}})$. The tensor product~$M$ is itself
in the thick closure of evaluation modules (each factor
$V_{\omega_i}(a_{ij})$ is an evaluation module). By descending
induction on the composition series, every composition factor~$L_k$
lies in the thick closure.

\emph{Step~3: All objects lie in the thick closure.}
An arbitrary object $X \in D^b(\mathrm{Rep}_{\mathrm{fd}})$ is a
bounded complex of finite-dimensional modules, each of finite length
with simple factors. By Step~2, every simple is in the thick closure
$\mathcal{D}_{\hbar}^{\mathrm{eval}}$. A thick subcategory of a
triangulated category is closed under extensions (via the octahedral
axiom) and bounded complexes (via iterated cones), so all of
$D^b(\mathrm{Rep}_{\mathrm{fd}})$ is contained in
$\mathcal{D}_{\hbar}^{\mathrm{eval}}$.
\end{proof}

\begin{remark}[Thick generation and DK-2]
\label{rem:dk2-thick-generation-consequence}
Proposition~\ref{prop:dk2-thick-generation-typeA} verifies
hypothesis~(ii) of Proposition~\ref{prop:yangian-dk2-thick-generation}
for the finite-dimensional regime. The remaining DK-2 gap for
type~$A$ is therefore:
\begin{enumerate}[label=\textup{(\alph*)}]
\item \emph{Ambient extension}: one needs exact functors on the chosen
 version of category~$\mathcal{O}$ extending the evaluation-locus
 Yangian correspondence to the ambient derived category; and
\item \emph{Infinite-dimensional extension}: thick generation of
 the chosen finite-length part of
 $D^b(\mathcal{O}_{Y_\hbar})$ by the same seed, with every
 standard, Verma-type, or prefundamental simple placed in that
 thick closure by an actual finite filtration.
\end{enumerate}
For~(a), the key input is a theorem that the relevant
bar-comodule comparison admits a Yangian-module realization on
category~$\mathcal{O}$ and respects its highest-weight structure.
For~(b), BGG-type resolutions or parabolic induction statements are
not enough by themselves: they must identify the actual composition
factors with objects already in the evaluation-generated thick
subcategory and must do so inside the completed ambient category used
for DK.

The type~$A$ thick generation result also
verifies the finite-dimensional part of the full-generation
hypothesis needed for
Proposition~\ref{prop:yangian-dk3-generated-core} \textup{(}DK-3\textup{)}.
\end{remark}

\begin{corollary}[Thick generation for all simple types;
\ClaimStatusProvedHere]
\label{cor:dk2-thick-generation-all-types}%
\index{thick generation!all simple types|textbf}%
For every simple Lie algebra~$\mathfrak{g}$, the evaluation
modules $\{V_{\omega_i}(a)\}$ thickly generate\/
$D^b(\mathrm{Rep}_{\mathrm{fd}}(Y_\hbar(\mathfrak{g})))$.
\end{corollary}

\begin{proof}
The three-step argument of
Proposition~\ref{prop:dk2-thick-generation-typeA} is
type-independent. Step~1 uses the Drinfeld classification
\textup{(}valid for all simple~$\mathfrak{g}$\textup{)}: every
finite-dimensional irreducible $Y_\hbar(\mathfrak{g})$-module is
the unique irreducible quotient of an ordered tensor product of
fundamental evaluation modules $V_{\omega_{i_1}}(a_1) \otimes
\cdots \otimes V_{\omega_{i_k}}(a_k)$, hence is a subquotient of
this tensor product.
Steps~2 and~3 are purely categorical
\textup{(}Lemma~\ref{lem:composition-thick-generation}
and closure of thick subcategories under bounded complexes\textup{)}
and carry over unchanged.
\end{proof}

\begin{lemma}[Thick generation from finite composition series;
\ClaimStatusProvedHere]
\label{lem:composition-thick-generation}
\index{thick generation!from composition series}
Let $\cA$ be an Artinian abelian category and
$\cC \subset D^b(\cA)$ a thick triangulated subcategory containing
every simple object of~$\cA$ (viewed in degree~$0$). Then
$\cC = D^b(\cA)$.
\end{lemma}

\begin{proof}
\emph{Step~1: Objects of $\cA$ lie in~$\cC$.}
Every $M \in \cA$ has a finite composition series
$0 = M_0 \subset M_1 \subset \cdots \subset M_r = M$
with each $M_i/M_{i-1}$ simple, hence in~$\cC$ by hypothesis.
We argue by induction on~$r$. For $r = 0$, $M = 0 \in \cC$.
For $r \ge 1$, the short exact sequence
$0 \to M_{r-1} \to M \to M_r/M_{r-1} \to 0$
yields a distinguished triangle
$M_{r-1} \to M \to M_r/M_{r-1} \xrightarrow{+1}$
in $D^b(\cA)$.
The outer terms lie in~$\cC$ (by induction and hypothesis),
so $M \in \cC$ since thick subcategories are closed under cones.

\emph{Step~2: Bounded complexes.}
For a bounded complex $C = [C^a \to \cdots \to C^b]$ with
$C^i \in \cA$, the stupid truncation triangle
$\sigma_{\ge a+1}C \to C \to C^a[-a] \xrightarrow{+1}$
and induction on the amplitude $b - a$ give $C \in \cC$.
\end{proof}

\begin{theorem}[Thick generation of category~\texorpdfstring{$\mathcal{O}$}{O} by
evaluation modules, type~\texorpdfstring{$A$}{A}; \ClaimStatusConditional]
\label{thm:catO-thick-generation}
\index{thick generation!Category O@Category $\mathcal{O}$|textbf}
\index{Drinfeld--Kohno!derived!Category O@Category $\mathcal{O}$}
For $\fg = \mathfrak{sl}_N$ with $N \ge 2$, let
$\mathcal{O}_{Y_\hbar}$ be a full subcategory of
$\mathcal{O}_{Y_\hbar(\mathfrak{sl}_N)}$ satisfying:
\begin{enumerate}[label=\textup{(\alph*)}]
\item $\mathcal{O}_{Y_\hbar}$ is Artinian
 \textup{(}every object has finite composition length\textup{)};
\item every simple object of $\mathcal{O}_{Y_\hbar}$ is
 finite-dimensional.
\end{enumerate}
Then the thick triangulated subcategory
$\mathcal{D}_\hbar^{\mathrm{eval}} \subset
D^b(\mathcal{O}_{Y_\hbar})$
generated by evaluation modules equals
$D^b(\mathcal{O}_{Y_\hbar})$.
Consequently, under the ambient exact-extension hypothesis of
Proposition~\textup{\ref{prop:yangian-dk2-thick-generation}}, this
verifies the generation input needed for DK-2 on
$\mathcal{O}_{Y_\hbar}$.
\end{theorem}

\begin{proof}
By Proposition~\ref{prop:dk2-thick-generation-typeA}, every
finite-dimensional simple $Y_\hbar(\mathfrak{sl}_N)$-module lies in
$\mathcal{D}_\hbar^{\mathrm{eval}}$.
Hypothesis~(b) then gives that every simple of
$\mathcal{O}_{Y_\hbar}$ lies in
$\mathcal{D}_\hbar^{\mathrm{eval}}$.
Hypothesis~(a) and Lemma~\ref{lem:composition-thick-generation}
yield $D^b(\mathcal{O}_{Y_\hbar}) =
\mathcal{D}_\hbar^{\mathrm{eval}}$.
\end{proof}

\begin{proposition}[Asymptotic prefundamentals as a conditional
generation criterion; \ClaimStatusConditional]
\label{prop:mc3-asymptotic-prefundamentals}
\index{thick generation!prefundamentals}
\index{QQ-system!Frenkel--Hernandez}
Let $\fg$ be simple of simply-laced type and let
$Y_\hbar(\fg)$ act on the Hernandez--Jimbo asymptotic category inside
a chosen separated completion. For each fundamental weight~$\omega_i$
and spectral parameter~$a \in \C$, let $L^{\pm}_{\omega_i, a}$ denote
the positive and negative prefundamental modules constructed
in~\cite{HernandezJimbo12,FrenkelHernandez15}. Assume:
\begin{enumerate}[label=\textup{(\roman*)},nosep]
\item these prefundamentals are objects of the chosen completed
 category and the relevant subcategory is finite length;
\item the Frenkel--Hernandez QQ-system identities lift from
 $K_0$-identities to finite exact filtrations there; and
\item tensoring with the finite-dimensional modules used in those
 filtrations preserves the completed subcategory.
\end{enumerate}
Then the objects whose simples occur in those lifted filtrations lie
in the thick subcategory generated by the
$L^{\pm}_{\omega_i,a}$. The statement is independent of the rank-$1$
Baxter hyperplane $b=a-\tfrac12$, but it is a criterion inside the
chosen completion, not a proof that the full category
$D^b(\mathcal O)$ is generated.
\end{proposition}

\begin{proof}
The Frenkel--Hernandez QQ-system~\cite{FrenkelHernandez15}
gives Pluecker-type identities for prefundamental characters in the
Grothendieck ring. A $K_0$ identity alone does not put an object in a
thick subcategory. Under hypotheses \textup{(i)--(iii)}, however, the
identity is represented by finite exact filtrations whose factors are
tensor products of the named prefundamentals and finite-dimensional
modules; closure under cones and tensoring gives the stated
membership. No type-$A$ Baxter balance is used. Without the filtration
and completion hypotheses this remains only a Grothendieck-group
identity.
\end{proof}

\begin{proposition}[Reflection-equation singular vectors for twisted
Yangians; \ClaimStatusConditional]
\label{prop:mc3-reflection-equation}
\index{twisted Yangian!reflection equation}
\index{Sklyanin determinant}
Let $Y^{\mathrm{tw}}_\hbar(\mathfrak{so}_n)$ or
$Y^{\mathrm{tw}}_\hbar(\mathfrak{sp}_n)$ denote the twisted
Yangians of Olshanski, with reflection-equation generator
$S(u)$ satisfying $R(u-v) S_1(u) R(u+v) S_2(v) = S_2(v)
R(u+v) S_1(u) R(u-v)$. Fix a finite-dimensional block. Assume the
Molev--Ragoucy classification realizes every simple in that block as
a subquotient of a generic reflection-equation fusion product of
evaluation modules, and assume that the Sklyanin determinant separates
the blocks under consideration. Then those evaluation modules thickly
generate the bounded derived category of the fixed block.
\end{proposition}

\begin{proof}
Molev--Ragoucy~\cite{MolevRagoucy02} classify finite-dimensional
irreducibles via Drinfeld-type tuples constrained by the
Sklyanin determinant $\mathrm{sdet}\, S(u)$, a central series
whose values label the block. The stated fusion hypothesis puts every
simple in the block in the thick closure of the evaluation seed.
Lemma~\ref{lem:composition-thick-generation} then gives the bounded
derived block. No Baxter constraint enters: the reflection equation
and the Sklyanin determinant are the block data in this twisted
setting. The conclusion is not a statement about ordinary Yangian
category~$\mathcal O$.
\end{proof}

\begin{theorem}[Finite-dimensional evaluation core from Drinfeld
polynomials;
\ClaimStatusProvedHere]
\label{thm:mc3-universal-multiplicity-free}
\index{Bethe subalgebra!joint spectrum}
\index{multiplicity-free!Chari--Moura}
For $\fg$ simple of any type and any
finite-dimensional irreducible $Y_\hbar(\fg)$-module~$L(\psi)$,
$L(\psi)$ is a subquotient of an ordered tensor product of fundamental
evaluation modules determined by its Drinfeld polynomials. Hence the
fundamental evaluation modules thickly generate
$D^b(\operatorname{Rep}_{\mathrm{fd}}(Y_\hbar(\fg)))$.
Generic Bethe-algebra diagonalisation and multiplicity-free
$q$-character results control generic finite-dimensional blocks; they
do not enlarge this theorem to full category~$\mathcal O$ or to a
completed prefundamental category.
\end{theorem}

\begin{proof}
The Drinfeld classification writes the highest $\ell$-weight of
$L(\psi)$ by Drinfeld polynomials. The roots of those polynomials
choose fundamental evaluation factors, and the ordered tensor product
has $L(\psi)$ as its irreducible quotient after the usual generic
ordering and limiting argument. Thus every simple finite-dimensional
module lies in the thick closure of the fundamental evaluation seed.
Lemma~\ref{lem:composition-thick-generation} gives the bounded
derived finite-dimensional category. The Bethe subalgebra is not used
as a substitute for this classification, and its simple-spectrum
claims are restricted to the generic finite-dimensional locus where
the cited Bethe results apply.
\end{proof}

\begin{proposition}[Elliptic thick generation on the theta-divisor complement;
\ClaimStatusConditional]
\label{prop:mc3-elliptic-theta-divisor}
\index{elliptic quantum group!theta-divisor}
\index{Felder algebra}
Let $E_{\rho, \eta}(\mathfrak{sl}_N)$ be the Felder elliptic
quantum group at modulus~$\rho$ and step~$\eta$, acting on
finite-dimensional regular blocks over the complement of the theta
divisor $\Theta \subset E_\rho$. Assume the elliptic Bethe
eigenvectors are complete, the joint spectrum is simple, and fusion
products have finite composition series off~$\Theta$. Then the
elliptic evaluation modules at points of $E_\rho \setminus \Theta$
thickly generate the corresponding finite-dimensional regular block.
\end{proposition}

\begin{proof}
Felder--Varchenko~\cite{FelderVarchenko96} construct the
evaluation representations; the dynamical R-matrix
$R(z, \lambda)$ has poles precisely on $\Theta$, and
fusion is well-defined on the complement. Under the stated
completeness and finite-length hypotheses, the Bethe eigenbasis
separates the simples in the regular block; the composition-length
argument of Lemma~\ref{lem:composition-thick-generation} then gives
thick generation of that block. The theta-divisor excision is a
regularity condition for the elliptic block, not a proof of generation
for an unrestricted elliptic category~$\mathcal O$.
\end{proof}

\begin{proposition}[Super-Yangian parity-balanced Baxter;
\ClaimStatusConditional]
\label{prop:mc3-super-parity-balance}
\index{super-Yangian!parity balance}
\index{Baxter polynomials!graded}
For the super-Yangian $Y_\hbar(\mathfrak{gl}_{m|n})$, let
$Q^\pm(u) = \prod_{k} (u - a^\pm_k)^{b^\pm_k}$ denote the
$\Z/2$-graded Baxter polynomials tracking even and odd
simple roots. The thick triangulated subcategory of
finite-dimensional $Y_\hbar(\mathfrak{gl}_{m|n})$-modules
generated by evaluation modules satisfying the parity-balance
constraint
\[
 b^+ - b^- \;=\; a - m + n
\]
equals the bounded derived category of the parity-balanced
finite-dimensional block, provided the graded Drinfeld classification
and the super-Shapovalov determinant give finite-length fusion
filtrations in that block.
\end{proposition}

\begin{proof}
Gow~\cite{Gow07} and Tsymbaliuk~\cite{Tsymbaliuk20super}
classify finite-dimensional irreducibles of
$Y_\hbar(\mathfrak{gl}_{m|n})$ by pairs $(Q^+, Q^-)$ of
polynomials of degrees $b^\pm$ satisfying the parity-balance
relation above; at $n = 0$ this recovers the ordinary
type-$A$ Baxter $b = a$ with the shift~$-1/2$ absorbed into
conventions. The $\Z/2$-graded Drinfeld presentation gives
an evaluation homomorphism $Y_\hbar(\mathfrak{gl}_{m|n}) \to
U(\mathfrak{gl}_{m|n})$. Under the stated finite-length fusion
hypotheses, every simple in the parity-balanced block is a
subquotient of tensor products of such evaluation modules, and
Lemma~\ref{lem:composition-thick-generation} applies. The shift
$-m+n$ is block data in the super setting; it is not a universal
Baxter hyperplane for ordinary Yangians.
\end{proof}

\begin{remark}[The unified mechanism: Baxter is a type-$A$ rational chart]
\label{rem:mc3-unified-mechanism}
\index{Baxter constraint!chart-local}
Propositions~\ref{prop:mc3-asymptotic-prefundamentals},
\ref{prop:mc3-reflection-equation},
\ref{prop:mc3-elliptic-theta-divisor},
\ref{prop:mc3-super-parity-balance} and
Theorem~\ref{thm:mc3-universal-multiplicity-free}
replace the type-$A$ Baxter constraint
$b = a - 1/2$ of
Theorem~\ref{thm:catO-thick-generation} by structural data
intrinsic to each family: the Frenkel--Hernandez QQ-system
(asymptotic), the Sklyanin determinant (twisted), the
Chari--Moura Bethe spectrum (uniform), the theta-divisor
complement (elliptic), and the super-trace parity balance
(graded). Read together, they identify the Baxter balance as a
rational type-$A$ chart: spectral-parameter families are governed in
each setting by the central series or dynamical locus natural to that
setting. The only unconditional theorem in this list is the
finite-dimensional evaluation core. The asymptotic, twisted, elliptic,
and super statements are conditional generation criteria in their
named ambients and do not close the compact/completed DK extension.
\end{remark}

\begin{remark}[Verification of hypotheses and scope]
\label{rem:catO-thick-generation-verification}
\index{Category O@Category $\mathcal{O}$!finite-dimensionality of simples}
Hypotheses (a) and~(b) of
Theorem~\ref{thm:catO-thick-generation} are hypotheses on the chosen
finite-length ambient category. They hold in the following
finite-dimensional or polynomial subcategories; they are not automatic
for the full Hernandez--Jimbo category~$\mathcal O$.
\begin{enumerate}[label=\textup{(\roman*)}]
\item \emph{$N = 2$, finite-dimensional block.}
 The finite-dimensional representation category for
 $Y_\hbar(\mathfrak{sl}_2)$ is controlled by
 Kirillov--Reshetikhin modules \cite{ChariPressley94}.
 Thick generation is therefore
 \emph{unconditional} for the finite-dimensional
 $\mathfrak{sl}_2$ block.
\item \emph{Polynomial category~$\mathcal{O}$, any~$N$.}
 The Hernandez--Leclerc subcategory
 $\mathcal{O}_{\mathrm{poly}} \subset
 \mathcal{O}_{Y_\hbar(\mathfrak{sl}_N)}$ (modules
 with polynomial highest $\ell$-weights
 \cite{HernandezLeclerc10}) has finite-dimensional simple objects
 when it is restricted to the Drinfeld-polynomial locus
 \cite{Drinfeld88}. Hence thick generation holds for that
 finite-dimensional polynomial subcategory at all~$N$.
\item \emph{Full category~$\mathcal{O}$ at $N \ge 3$.}
 The full category contains asymptotic and prefundamental objects
 outside the finite-dimensional Drinfeld-polynomial locus. The
 remaining question is therefore not settled by the polynomial
 subcategory: one must specify the completed finite-length ambient
 category and prove that its simples or standards lie in the
 evaluation-generated thick closure.
\end{enumerate}

The obstruction in
Remark~\ref{rem:cat-o-generation-obstruction} is therefore not the
use of composition series.
Lemma~\ref{lem:composition-thick-generation} shows that a
finite composition series with simple factors in a thick
subcategory \emph{does} prove membership: every short exact
sequence in an abelian category yields a distinguished triangle
in the derived category, and thick subcategories are closed under
cones. The generation obstruction for category~$\mathcal{O}$
therefore has two separate parts: the representation-theoretic
finite-length/seed-containment statement, and the ambient exact
extension/completion statement required by the DK comparison.
\end{remark}

\begin{remark}[Bar-comodule route to a factorization Kazhdan candidate]
\label{rem:bar-cobar-kazhdan-candidate}
\index{Kazhdan functor!factorization}
On the evaluation-generated factorization core
$\operatorname{Fact}_{\Eone}^{\mathrm{eval}}(\cA)$ of an $\Eone$-chiral Yangian-type algebra~$\cA$,
Theorem~\ref{thm:e1-module-koszul-duality} gives the intrinsic module
bar functor
\[
\barB_{\mathrm{mod}}\colon
\operatorname{Mod}(\cA)
\to
\operatorname{CoMod}^{\mathrm{conil}}_{\barB(\cA)}.
\]
On evaluation generators, the proved DK square
\textup{(}Theorem~\ref{thm:derived-dk-yangian}\textup{)} and the
evaluation-locus factorization theorem
\textup{(}Theorem~\ref{thm:factorization-dk-eval}\textup{)} show that
any additional comparison functor from the resulting bar-comodule image
to the quantum-group evaluation core must send
$V(a)_\hbar \mapsto V(-a)_{-\hbar}$ and reverse the braiding
$R \mapsto R^{-1}$.

Thus bar-cobar remains a natural route to the factorization Kazhdan
transport required by
Proposition~\ref{prop:yangian-dk3-generated-core}, but the actual
bar-comodule-to-quantum-group comparison is extra input rather than a
formal consequence of the module bar comparison itself.
\end{remark}

\begin{remark}[Remaining DK-3 input after the bar-cobar candidate]
\label{rem:dk3-remaining-after-barcobar}
The previous remark isolates the real remaining DK-3 inputs on the
evaluation-generated core:
\begin{enumerate}[label=\textup{(\roman*)}]
\item a comparison from the relevant bar-comodule image to the
 quantum-group evaluation core, together with a quasi-inverse on that
 core; and
\item compatibility of the bar-cobar unit and counit with
 factorization products after passing through that comparison.
\end{enumerate}

Concretely, the bar-cobar adjunction itself provides the intrinsic
bar-side unit and counit
$\eta\colon \mathrm{id} \to \Omega \circ \barB$ and counit
$\varepsilon\colon \barB \circ \Omega \to \mathrm{id}$ are
quasi-isomorphisms on Koszul objects
\textup{(}Theorem~\textup{\ref{thm:e1-chiral-koszul-duality})}.
On evaluation modules, the proved DK-1 theorem shows that these
adjunction maps match the expected evaluation-level identifications.
Beyond that locus, however, neither the bar-comodule/quantum-group
comparison nor its factorization-monoidal/quasi-inverse enhancement is
automatic from bar-cobar alone.
\end{remark}

\begin{lemma}[Monoidal extension to thick closures;
\ClaimStatusProvedHere]
\label{lem:monoidal-thick-extension}
Let $F\colon \cC \to \cD$ be an exact functor between tensor
triangulated categories equipped with a lax monoidal structure
$\mu_{M,N}\colon F(M \otimes N) \to F(M) \otimes F(N)$.
Let $\mathcal{S} \subset \cC$ generate~$\cC$ as a thick subcategory.
If $\mu_{M,N}$ is a quasi-isomorphism for all $M, N \in \mathcal{S}$,
then $\mu_{M,N}$ is a quasi-isomorphism for all $M, N \in \cC$.
In particular, $F$ is strong monoidal.
\end{lemma}

\begin{proof}
For fixed $M \in \mathcal{S}$, the collection
$\{N \in \cC : \mu_{M,N} \text{ is a quasi-isomorphism}\}$
contains~$\mathcal{S}$ and is closed under shifts (naturality of~$\mu$),
cones (five lemma, since $F(M \otimes -)$ and
$F(M) \otimes F(-)$ are both exact), and direct summands
(retract argument). Hence it equals~$\cC$.
Exchanging the two variables completes the proof.
\end{proof}

\begin{conjecture}[Finite-dimensional extension criterion from evaluation generators,
type~\texorpdfstring{$A$}{A}; \ClaimStatusConjectured]
\label{conj:dk-fd-typeA}
\index{Drinfeld--Kohno!derived!finite-dimensional|textbf}
For $\fg = \mathfrak{sl}_N$ with $N \ge 2$, assume:
\begin{enumerate}[label=\textup{(\alph*)}]
\item there is an exact ambient extension
 \[
 \widetilde{\Phi}\colon
 D^b(\operatorname{Rep}_{\mathrm{fd}}(Y_\hbar(\mathfrak{sl}_N)))
 \to
 D^b(\operatorname{Rep}_{\mathrm{fd}}(Y_{-\hbar}(\mathfrak{sl}_N)))
 \]
 together with a quasi-inverse, whose restriction to the
 evaluation-generated subcategory agrees with the proved
 evaluation-locus equivalence of
 Theorems~\textup{\ref{thm:factorization-dk-eval}}
 and~\textup{\ref{thm:derived-dk-yangian}\textup{;}};
\item $\widetilde{\Phi}$ carries a lax monoidal structure whose
 structure maps are quasi-isomorphisms on pairs of evaluation
 modules; and
\item there is a compatible braided monoidal comparison from the image
 of~$\widetilde{\Phi}$ to
 $\operatorname{Fact}_{\Eone}^{\mathrm{fd}}(U_q(\mathfrak{sl}_N))^{\mathrm{op}}$
 agreeing with the evaluation-level DK/Kazhdan identifications on
 generators.
\end{enumerate}
Then thick generation by evaluation modules promotes the
evaluation-locus DK equivalence to a braided monoidal equivalence on
all finite-dimensional representations; in particular the composite
$KZ_{\barB}$ gives the finite-dimensional factorization DK bridge
under these ambient-extension and comparison hypotheses.
\end{conjecture}

\begin{remark}[Finite-dimensional extension mechanism]
Proposition~\ref{prop:dk2-thick-generation-typeA} shows that
evaluation modules thickly generate
$D^b(\operatorname{Rep}_{\mathrm{fd}}(Y_{\pm\hbar}(\mathfrak{sl}_N)))$.
Applying Proposition~\ref{prop:yangian-dk2-thick-generation} to the
ambient extension of hypothesis~\textup{(a)} extends the
evaluation-locus equivalence to all finite-dimensional
representations.

By hypothesis~\textup{(b)}, the ambient extension carries a lax
monoidal structure that is an isomorphism on pairs of evaluation
modules. Lemma~\ref{lem:monoidal-thick-extension} and thick generation
therefore make it a strong monoidal structure on the whole
finite-dimensional category.

Hypothesis~\textup{(c)} identifies the resulting finite-dimensional
Yangian-side equivalence with the corresponding quantum-group
factorization category on the same thick closure, yielding the stated
finite-dimensional DK bridge.
\end{remark}

\begin{corollary}[Type-$A$ evaluation-generated extension principle;
\ClaimStatusConditional]\label{cor:dk-partial-conj}
For $\mathfrak{g}=\mathfrak{sl}_N$, the proof of
Conjecture~\ref{conj:dk-fd-typeA} extends verbatim to any thick
subcategory of $\operatorname{Fact}_{\Eone}$ generated by evaluation
modules whenever the same ambient exact/lax monoidal extension and
comparison hypotheses are available there. In particular, on the
evaluation-generated core of
$D^b(\mathcal{O}_{Y_\hbar(\mathfrak{sl}_N)})$, the only additional
obstacles are the ambient extension/comparison input and the question
of whether evaluation modules thickly generate the full category.
\end{corollary}

\begin{proof}
This is exactly the same thick-closure argument as in
Conjecture~\ref{conj:dk-fd-typeA}, applied on the chosen
evaluation-generated subcategory.
\end{proof}

\begin{remark}[Application to the full DK-3 conjecture]
\label{rem:dk-partial-conj-application}
Corollary~\ref{cor:dk-partial-conj} isolates the exact remaining
input in Conjecture~\ref{conj:full-derived-dk}: one must prove that
evaluation modules thickly generate
$D^b(\mathcal{O}_{Y_\hbar(\mathfrak{sl}_N)})$ and that the required
ambient extension/comparison package exists on that category.
\end{remark}

\begin{corollary}[Factorization DK for the finite-dimensional
polynomial locus, type~\texorpdfstring{$A$}{A}; \ClaimStatusConditional]
\label{cor:dk-poly-catO}
\index{Drinfeld--Kohno!derived!polynomial Category O@polynomial Category $\mathcal{O}$}
For $\fg = \mathfrak{sl}_N$ with $N \ge 2$, let
$\mathcal{O}_{\mathrm{poly}}^{\mathrm{fd}}$ be the full
finite-dimensional Drinfeld-polynomial subcategory of
$\mathcal{O}_{\mathrm{poly}}$. Assume the ambient
extension/comparison hypotheses of
Conjecture~\ref{conj:dk-fd-typeA} hold on this subcategory. Then the
induced factorization Drinfeld--Kohno functor is a braided monoidal
equivalence
\[
KZ_{\barB}\colon
\operatorname{Fact}_{\Eone}^{\mathcal{O}_{\mathrm{poly}}^{\mathrm{fd}}}
 (Y_\hbar(\mathfrak{sl}_N))
\;\xrightarrow{\;\sim\;}\;
\operatorname{Fact}_{\Eone}^{\mathrm{fd}}
 (U_q(\mathfrak{sl}_N))^{\mathrm{op}}.
\]
\end{corollary}

\begin{proof}
Remark~\ref{rem:catO-thick-generation-verification}(ii) and
Theorem~\ref{thm:catO-thick-generation} show that
$D^b(\mathcal{O}_{\mathrm{poly}}^{\mathrm{fd}})$ is thickly generated
by evaluation modules at all~$N$. Under the stated ambient-extension
and comparison hypotheses, Conjecture~\ref{conj:dk-fd-typeA}
therefore applies verbatim on this thick closure.
\end{proof}

\begin{remark}[Finite-dimensional polynomial locus is the natural DK domain]
\label{rem:poly-O-natural-domain}
Corollary~\ref{cor:dk-poly-catO} is the natural statement of
factorization DK for the Yangian, for the following reason.
The quantum group side of the equivalence is
$\operatorname{Rep}_{\mathrm{fd}}(U_q(\mathfrak{sl}_N))$: finite-dimensional
representations classified by Drinfeld polynomials. These correspond,
under the KL identification, to
$\mathcal{O}_{\mathrm{poly}}^{\mathrm{fd}}$, not to the full
category~$\mathcal{O}$.
The ``full category~$\mathcal{O}$ generation problem'' of
Corollary~\ref{cor:dk-partial-conj}
is therefore a question about whether
$\mathcal{O}$ and $\mathcal{O}_{\mathrm{poly}}^{\mathrm{fd}}$
coincide as DK-relevant categories, not about whether the DK
equivalence extends beyond its natural finite-dimensional domain.

If $\mathcal{O} \supsetneq
\mathcal{O}_{\mathrm{poly}}^{\mathrm{fd}}$ (i.e., if there exist
simple modules of $\mathcal{O}$ that are \emph{not}
finite-dimensional), then these additional modules have no quantum
group counterpart under the present formulation and lie outside the
finite-dimensional DK statement. The present theorem surface proves
only the thick-generation extension criterion on the finite-dimensional
polynomial locus, not the ambient extension/comparison package itself.

Whether $\mathcal{O} =
\mathcal{O}_{\mathrm{poly}}^{\mathrm{fd}}$ is an open question; see
Corollary~\ref{cor:dk-partial-conj}. If the inclusion is strict, the
additional modules are either irrelevant to finite-dimensional DK
or indicate missing structure on the quantum group side; the
present framework does not distinguish these possibilities.
\end{remark}

\begin{remark}[DK ladder after the finite-dimensional theorem]
\label{rem:dk-fd-synthesis}
DK-0/1: proved for all simple~$\fg$. DK-2/3: the thick-generation
statements for finite-dimensional Yangian representations are proved
\textup{(}Proposition~\ref{prop:dk2-thick-generation-typeA} and
Corollary~\ref{cor:dk2-thick-generation-all-types}\textup{)}. For
finite-length category~$\mathcal O$ subcategories with
finite-dimensional simples, Theorem~\ref{thm:catO-thick-generation}
gives the corresponding conditional generation criterion. Promotion
of evaluation-locus DK beyond the generators is still conditional on
the ambient extension/comparison package of
Conjecture~\ref{conj:dk-fd-typeA}; a separate H-level sectorwise route
exists at all simple types
\textup{(}Corollary~\ref{cor:dk23-all-types}\textup{)}. Extension
beyond $\mathcal{O}_{\mathrm{poly}}^{\mathrm{fd}}$ requires
completed/coderived techniques. DK-4: structural framework proved
(MC4 closed via
$\mathcal{W}_N$ rigidity, Theorem~\ref{thm:completed-bar-cobar-strong});
algebraic identification open (Conjecture~\ref{conj:dk4-formal-moduli}).
DK-5: still conjectural; even after the all-types generated-core DK
comparison, it remains downstream of the post-core
extension/completion packet and the later categorical bridge inputs.
\end{remark}

\begin{lemma}[Finite-dimensional thick-closure constraint;
\ClaimStatusProvedHere]
\label{lem:fd-thick-closure}
\index{thick generation!finite-dimensional constraint}
Let $\cA$ be an abelian category with enough injectives, and let
$\cA^{\mathrm{fd}} \subset \cA$ denote the full subcategory of
finite-length objects. Then
\[
\operatorname{thick}_{D^b(\cA)}\!\bigl(\cA^{\mathrm{fd}}\bigr)
\;=\; D^b(\cA^{\mathrm{fd}}).
\]
In particular, if $\cA$ contains objects of infinite length, the
thick closure of finite-length objects in $D^b(\cA)$ is strictly
smaller than $D^b(\cA)$.
\end{lemma}

\begin{proof}
Every complex in $\operatorname{thick}(\cA^{\mathrm{fd}})$ is built
from finite-length objects by finitely many shifts, cones, and
retracts. Since the full subcategory
$D^b(\cA^{\mathrm{fd}}) \subset D^b(\cA)$ is a thick triangulated
subcategory containing $\cA^{\mathrm{fd}}$, one inclusion is
immediate. Conversely, every bounded complex with finite-length
cohomology lies in $\operatorname{thick}(\cA^{\mathrm{fd}})$ by
d\'{e}vissage (induction on the amplitude and composition length of
cohomology, using the stupid truncation triangles).
\end{proof}

\begin{remark}[Implications for the generation problem]
\label{rem:fd-thick-closure-implications}
\index{thick generation!scope limitation}
Lemma~\ref{lem:fd-thick-closure} shows that the literal generation
question
$\mathcal{C}_{\mathrm{eval}} \stackrel{?}{=}
D^b(\mathcal{O}_{Y_\hbar(\mathfrak{sl}_N)})$
has a \emph{positive} answer when restricted to
$\mathcal{O}_{\mathrm{poly}}$ or to $N = 2$ (where all simples are
finite-dimensional), and a \emph{negative} answer if the full
category~$\mathcal{O}$ contains genuine infinite-length objects.
The correct target for extending beyond
$\mathcal{O}_{\mathrm{poly}}$ is therefore not literal thick
generation in $D^b(\mathcal{O})$, but one of:
\begin{enumerate}[label=\textup{(\alph*)}]
\item the statement
 $\operatorname{thick}(\text{evaluation modules}) =
 D^b(\mathcal{O}_{\mathrm{poly}})$
 (matching the proved scope);
\item \emph{localizing} generation in an Ind- or
 pro-completed enhancement
 $\widehat{D}(\mathcal{O})$,
 where evaluation modules generate a localizing
 subcategory whose compact part recovers the evaluation-generated
 factorization core;
\item enlargement of the generator set to include
 prefundamental or asymptotic modules.
\end{enumerate}
This explains the role of the polynomial subcategory as a
\emph{natural domain} for the DK equivalence
(Remark~\ref{rem:catO-thick-generation-verification}).
\end{remark}

\subsubsection*{Strategies for the category~$\mathcal{O}$
generation problem}
\label{sec:cat-O-strategies}
\index{Category O@Category $\mathcal{O}$!generation strategies}

Corollary~\ref{cor:dk-partial-conj} reduces the full
Conjecture~\ref{conj:full-derived-dk} to a single problem: whether
the DK-relevant subcategory of
$\mathcal{O}_{Y_\hbar(\mathfrak{sl}_N)}$ is thickly generated
by evaluation modules.
By Lemma~\ref{lem:fd-thick-closure}, this is equivalent to asking
whether every simple in the DK-relevant subcategory is
finite-dimensional (i.e., whether the DK-relevant subcategory
is contained in $\mathcal{O}_{\mathrm{poly}}$).
We develop four approaches, in decreasing
order of structural ambition. Each reduces the generation problem to
a different (but related) input. See
Remark~\ref{rem:cat-o-generation-obstruction} in the concordance
for the summary assessment.

\medskip

\noindent\textbf{Strategy~IV\@: Sectorwise finiteness
generalization.}

\begin{proposition}[Loop-weight filtration of the Yangian bar
complex; \ClaimStatusProvedHere]
\label{prop:yangian-bar-loop-weight}
\index{Yangian!bar complex!loop-weight filtration}
For any simple Lie algebra~$\fg$, the factorization bar complex
$\barB^{\mathrm{fact}}(Y_\hbar(\fg))$ admits:
\begin{enumerate}[label=\textup{(\alph*)}]
\item an exact direct sum decomposition by root lattice weight:
 $\barB^{\mathrm{fact}} = \bigoplus_{\gamma \in
 Q(\fg)} \barB^{\mathrm{fact}}_\gamma$;
\item a bounded-below descending filtration by loop degree:
 $F^0 \barB^{\mathrm{fact}}_\gamma \supset
 F^1 \barB^{\mathrm{fact}}_\gamma \supset \cdots$,
 where $F^p$ is spanned by bar monomials with total loop index
 $\ge p$;
\item in each bidegree $(\gamma, p)$, the graded piece
 $\operatorname{gr}^p \barB^{\mathrm{fact}}_\gamma =
 F^p / F^{p+1}$ is finite-dimensional over each stratum of
 $\operatorname{Ran}(X)$.
\end{enumerate}
The associated spectral sequence has $E_0$ differential given by
the bar differential of
$\operatorname{gr}^F Y_\hbar(\fg) \cong
U(\fg[t])$, the universal enveloping algebra of the
current algebra.
\end{proposition}

\begin{proof}
(a)~The root lattice grading of $Y_\hbar(\fg)$ (the generators $J^{(r)}_a$ carry weight
$\alpha(a) \in Q(\fg)$ determined by the root
system) is preserved by all defining relations. The bar
differential, being constructed from the product, preserves root
weight. Hence the bar complex decomposes as a direct sum over
$Q(\fg)$.

(b)~Define $\deg_{\mathrm{loop}} J^{(r)}_a = r$. The PBW
theorem for $Y_\hbar(\fg)$ \cite{Drinfeld85,molev-yangians} gives
a $k$-basis of ordered monomials in $\{J^{(r)}_a\}_{r \ge 0}$.
The filtration $F^p$ on the bar complex is spanned by monomials
$s^{-1}m_1 | \cdots | s^{-1}m_n$ with
$\sum_i \deg_{\mathrm{loop}}(m_i) \ge p$. The Yangian defining
relations express $[J^{(r)}_a, J^{(s)}_b]$ in terms of generators
of loop degree $\le r + s - 1$, so the bar differential satisfies
$d(F^p) \subset F^{p-1}$ and the filtration is compatible.

(c)~At fixed root weight $\gamma$ and loop degree $p$, the number
of PBW monomials is the number of partitions of $p$ into
non-negative integers assigned to basis elements whose root weights
sum to $\gamma$. This is finite.
\end{proof}

\begin{remark}[Comparison with the lattice bypass]
\label{rem:sectorwise-vs-lattice}
\index{sectorwise finiteness!Yangian vs.\ lattice}
Theorem~\ref{thm:lattice:factorization-koszul}
(DK-$1\frac{1}{2}$) proves the lattice factorization bar-cobar
equivalence via sectorwise finiteness: the $\Lambda$-grading gives
an exact decomposition, and within each sector, the weight
filtration degenerates at $E_1$ by Heisenberg acyclicity
(Theorem~\ref{thm:heisenberg-koszul-dual-early}).
Proposition~\ref{prop:yangian-bar-loop-weight} provides the
analogous structure with two differences:
\begin{enumerate}[label=\textup{(\alph*)}]
\item The loop direction is a \emph{filtration}, not a grading:
 the Yangian relations mix loop degrees
 ($[J^{(1)}_a, J^{(0)}_b]$ involves terms at loop degree~$0$),
 producing a spectral sequence rather than a direct sum
 decomposition.
\item The $E_1$ page is the factorization bar cohomology of the
 current algebra $U(\fg[t])$, which is nontrivial:
 for $\mathfrak{g}$ semisimple,
 $H^\bullet(\mathfrak{g}[t], k)$ is strictly larger than
 $H^\bullet(\mathfrak{g}, k)$
 (it involves the continuous cohomology of the current algebra).
\end{enumerate}
Consequently, the spectral sequence does \emph{not} degenerate at
$E_1$, and sectorwise finiteness alone does not immediately resolve
the enlargement problem beyond the evaluation-generated core.

Despite this, the structure provides a complexity reduction: any
completed/coderived enlargement controlled sectorwise reduces to a
countable family of finite-dimensional spectral sequence
computations at each $(\gamma, p)$ bidegree. If the same finiteness
and convergence package survives the prefundamental/asymptotic
extension, then the sectorwise method may enlarge the proved
evaluation-generated DK core without thick generation.

A successful extension carries the proved evaluation-generated
DK-2/3 equivalence to the
completed/coderived factorization domain relevant to the remaining
MC3 packets for \emph{all types} simultaneously, since the root
lattice grading and loop filtration exist for any
simple~$\mathfrak{g}$. The bottleneck is no longer convergence on
the generated core, but the extension of that control to the enlarged
completed/coderived domain. The $E_1$ computation is the first
well-posed input, and it is related to the cohomology of current
algebras via the Loday--Quillen--Tsygan theorem.
The abstract criterion subsuming both the lattice and Yangian
instances is
Conjecture~\ref{conj:factorization-finiteness-criterion}.
\end{remark}

\begin{computation}[The \texorpdfstring{$E_1$}{E_1} page for \texorpdfstring{$\mathfrak{sl}_2{[t]}$}{sl_2[t]}]
\label{comp:current-algebra-E1}
\index{current algebra!cohomology}
\index{Loday--Quillen--Tsygan theorem}
We compute $H^\bullet(\mathfrak{sl}_2[t], k)$ (the $E_1$
page of the Strategy~IV spectral sequence when
$\mathfrak{g} = \mathfrak{sl}_2$) by the Hochschild--Serre
spectral sequence for the extension
$0 \to \mathfrak{sl}_2 \otimes t\,k[t] \to \mathfrak{sl}_2[t]
\to \mathfrak{sl}_2 \to 0$.

\emph{Step~1: The abelian ideal.}
Let $\mathfrak{n} = \mathfrak{sl}_2 \otimes t\,k[t]$, a
pro-nilpotent Lie algebra. As an $\mathfrak{sl}_2$-module,
$\mathfrak{n} \cong \mathfrak{sl}_2^{\oplus \infty}$
(copies indexed by $t, t^2, t^3, \ldots$), each carrying the
adjoint representation $V_2$ (the $3$-dimensional irreducible
of~$\mathfrak{sl}_2$).

\emph{Step~2: Hochschild--Serre $E_2$ page.}
The HS spectral sequence has
\[
E_2^{p,q} = H^p(\mathfrak{sl}_2,\, H^q(\mathfrak{n}, k))
\;\Longrightarrow\; H^{p+q}(\mathfrak{sl}_2[t], k).
\]
The $\mathfrak{n}$-cohomology is the continuous dual of the
Chevalley--Eilenberg complex. Since~$\mathfrak{n}$ is
abelian as a graded Lie algebra (the Lie bracket
$[t^r x, t^s y] = t^{r+s}[x,y]$ maps to higher loop degree, so
it is zero on the associated graded), we have
\[
H^q(\mathfrak{n}, k) = \Lambda^q(\mathfrak{n}^*)
= \Lambda^q\!\Bigl(\bigoplus_{r \ge 1} V_2^*\Bigr)
\]
as an $\mathfrak{sl}_2$-representation, where
$V_2^* \cong V_2$ (the adjoint is self-dual).

\emph{Step~3: Invariants.}
For the $E_2$ page we need
$H^p(\mathfrak{sl}_2, \Lambda^q(\bigoplus_{r \ge 1} V_2))$.

By the K\"unneth formula and $\dim V_2 = 3$, the $q = 0$ line is
\[
H^p(\mathfrak{sl}_2, k)
= \begin{cases} k & p = 0, 3 \\ 0 & \text{else}. \end{cases}
\]

For $q = 1$ we get
\[
H^p(\mathfrak{sl}_2, \bigoplus_{r \ge 1} V_2)
= \bigoplus_{r \ge 1} H^p(\mathfrak{sl}_2, V_2)
= 0,
\]
because $V_2$ is irreducible and nontrivial.

For $q = 2$ we have
\[
\Lambda^2\Bigl(\bigoplus_r V_2\Bigr)
\supset
\bigoplus_{r < s} V_2 \otimes V_2
\oplus \bigoplus_r \Lambda^2(V_2).
\]

Here $\Lambda^2(V_2) \cong V_2$ (antisymmetric square of the
adjoint of~$\mathfrak{sl}_2$, isomorphic to the adjoint itself),
and $V_2 \otimes V_2 \cong V_4 \oplus V_2 \oplus V_0$
(Clebsch--Gordan for $3 \otimes 3 = 5 + 3 + 1$). The invariant
part $V_0$ appears once per pair $(r,s)$ with $r < s$, giving
\[
\dim H^0(\mathfrak{sl}_2, \Lambda^2(\textstyle\bigoplus_{r \ge 1}
V_2)) = \binom{\infty}{2} \;\text{(one copy per pair)}.
\]
In the graded setting (fixed total loop degree $p$), the number of
pairs $(r,s)$ with $r + s = p$ is $\lfloor (p-1)/2 \rfloor$ for
$p \ge 3$, giving a \emph{finite-dimensional} $E_2^{0,2}$
at each loop degree.

\emph{Step~4: Conclusion.}
The $E_1$ page of the Strategy~IV spectral sequence, at root
weight~$0$ and loop degree~$p$, has dimension bounded by the
$\mathfrak{sl}_2$-invariants in the exterior algebra of
$p - 1$ copies of $V_2$ (the copies of the adjoint at
loop degrees $1, \ldots, p-1$ contributing to total degree~$p$).
This is a finite computation at each~$p$. The first values are:
\begin{center}
\begin{tabular}{c|cccccc}
$p$ & 0 & 1 & 2 & 3 & 4 & 5 \\
\hline
$\dim E_1^{0,p}$ & 1 & 0 & 0 & 1 & 0 & 1 \\
\end{tabular}
\end{center}
where $p = 0$ is the constant, $p = 3$ is the CE $3$-cocycle of
$\mathfrak{sl}_2$, and $p = 5$ is the first contribution from the
$V_0$ summand at $q = 2$. The growth of
$\dim E_1^{0,p}$ is at most polynomial in~$p$, confirming the
complexity reduction claimed in
Remark~\ref{rem:sectorwise-vs-lattice}.
\end{computation}

\begin{computation}[The \texorpdfstring{$E_1$}{E_1} page for \texorpdfstring{$\mathfrak{sp}_4{[t]}$}{sp_4[t]} and
 all types via LQT]
\label{comp:current-algebra-E1-all-types}
\index{current algebra!cohomology!all types}
By the Feigin--Tsygan formulation of the
Loday--Quillen--Tsygan theorem, for any simple~$\mathfrak{g}$
with exponents $e_1, \ldots, e_r$:
\[
H^*(\mathfrak{g}[t], k)
\;\cong\;
\Lambda\bigl(\xi_{i,n} \mid
 1 \le i \le r,\; n \ge 0\bigr),
\qquad
\deg_{\mathrm{loop}}(\xi_{i,n}) = 2e_i + 1 + 2n.
\]
Thus $\dim E_1^{0,p}$ equals the number of subsets
of $\{2e_i + 1 + 2n : 1 \le i \le r,\, n \ge 0\}$
summing to~$p$, a restricted partition function computable
by dynamic programming.

For $\mathfrak{sp}_4 = C_2$ (exponents $1, 3$;
first non-simply-laced test case,
$h^\vee = 3 \neq h = 4$), generators appear at loop degrees
$\{3, 5, 7, 9, \ldots\} \cup \{7, 9, 11, \ldots\}$, giving:
\begin{center}
\begin{tabular}{c|cccccccccc}
$p$ & 0 & 1 & 2 & 3 & 4 & 5 & 6 & 7 & 8 & 9 \\
\hline
$\dim E_1^{0,p}(\mathfrak{sl}_2)$
 & 1 & 0 & 0 & 1 & 0 & 1 & 0 & 1 & 1 & 1 \\
$\dim E_1^{0,p}(\mathfrak{sl}_3)$
 & 1 & 0 & 0 & 1 & 0 & 2 & 0 & 2 & 2 & 2 \\
$\dim E_1^{0,p}(\mathfrak{sp}_4)$
 & 1 & 0 & 0 & 1 & 0 & 1 & 0 & 2 & 1 & 2 \\
\end{tabular}
\end{center}
The $\mathfrak{sp}_4$ sequence first departs from
$\mathfrak{sl}_2$ at $p = 7$ (two generators vs.\ one,
from the second exponent $e_2 = 3$). Growth is sub-exponential
($\sim \exp(\pi\sqrt{rp/12})$ by
Proposition~\ref{prop:lqt-e1-subexponential-growth})
for all types; this is the proved growth input for
Conjecture~\textup{\ref{conj:mc3-sectorwise-all-types}}.
\end{computation}

\medskip

\noindent\textbf{Strategy~III\@: Reduction to the
Kazhdan--Lusztig problem.}

\begin{proposition}[Thick generation via projective resolutions;
\ClaimStatusProvedHere]
\label{prop:thick-gen-projective}
\index{thick generation!projective resolution criterion}
Let $\mathcal{O}$ be a highest-weight category with finite global
dimension, and let $\mathcal{C}_{\mathrm{eval}} \subset
D^b(\mathcal{O})$ be the thick subcategory generated by evaluation
modules. Then $\mathcal{C}_{\mathrm{eval}} = D^b(\mathcal{O})$ if
and only if every projective indecomposable $P(\lambda)$ of
$\mathcal{O}$ admits a finite resolution
\begin{equation}\label{eq:proj-eval-resolution}
0 \to F_n \to \cdots \to F_1 \to F_0 \to P(\lambda) \to 0
\end{equation}
with each $F_i \in \mathcal{C}_{\mathrm{eval}}$.
\end{proposition}

\begin{proof}
($\Rightarrow$):
If $\mathcal{C}_{\mathrm{eval}} = D^b(\mathcal{O})$, then
$P(\lambda) \in D^b(\mathcal{O}) = \mathcal{C}_{\mathrm{eval}}$
and such a resolution exists by definition of thick closure.

($\Leftarrow$):
Suppose every $P(\lambda)$ admits a finite evaluation
resolution~\eqref{eq:proj-eval-resolution}. Then
$P(\lambda) \in \mathcal{C}_{\mathrm{eval}}$. Since $\mathcal{O}$
has finite global dimension, every object $M \in \mathcal{O}$ admits
a finite projective resolution. Each projective term lies in
$\mathcal{C}_{\mathrm{eval}}$, so $M \in
\mathcal{C}_{\mathrm{eval}}$ by iterated cones. Every object of
$D^b(\mathcal{O})$ is a bounded complex of objects of
$\mathcal{O}$, hence lies in $\mathcal{C}_{\mathrm{eval}}$.
\end{proof}

\begin{remark}[Kazhdan--Lusztig input for Yangian
category~\texorpdfstring{$\mathcal{O}$}{O}]
\label{rem:kl-input-yangian}
\index{Kazhdan--Lusztig!Yangian thick generation}
For $Y_\hbar(\mathfrak{sl}_N)$,
Proposition~\ref{prop:thick-gen-projective} reduces thick generation
to two inputs:
\begin{enumerate}[label=\textup{(\alph*)}]
\item \emph{Finite global dimension.}
 For truncated shifted Yangians, this follows from the
 Brundan--Kleshchev equivalence \cite{BrunKlesh06} with parabolic
 BGG category~$\mathcal{O}$ for $\mathfrak{gl}_M$, which has
 finite global dimension. For the full Yangian
 $Y(\mathfrak{sl}_N)$, the corresponding finiteness remains
 conjectural, but the analogous finiteness holds in Hernandez's
 category~$\mathcal{O}_\bZ$ \cite{Hernandez05}.
\item \emph{Projective modules in evaluation thick closure.}
 The projective covers $P(\lambda)$ in $\mathcal{O}$ have Verma
 flags with multiplicities $(P(\lambda) : M(\mu))$ given by
 Kazhdan--Lusztig polynomials (via BGG reciprocity:
 $(P(\lambda) : M(\mu)) = [M(\mu) : L(\lambda)]$).
 The combinatorial data is accessible from the KL theory for
 Yangians developed by Brundan--Kleshchev \cite{BrunKlesh06} and
 Kang--Kashiwara--Kim--Oh \cite{KKKO18}. The remaining step is
 to show that the Verma modules appearing in these flags lie in
 $\mathcal{C}_{\mathrm{eval}}$, connecting to Strategy~I\@.
\end{enumerate}
This is the most classical approach: the combinatorial structure is
provided by the KL theory, and the analytical input reduces to
resolvability of Verma modules. For type~$A$, the
Brundan--Kleshchev equivalence makes the finite global dimension
unconditional.
\end{remark}

\medskip

\noindent\textbf{Strategy~II\@: t-structure argument.}

This is now \emph{proved}: Lemma~\ref{lem:composition-thick-generation}
and Theorem~\ref{thm:catO-thick-generation} implement Strategy~II
completely. The point
(Remark~\ref{rem:catO-thick-generation-verification}) is that
thick generation reduces to the finite-dimensionality of simples,
which holds unconditionally for $\mathfrak{sl}_2$ and for
the polynomial subcategory $\mathcal{O}_{\mathrm{poly}}$ at all~$N$.

\begin{remark}[Residual obstruction beyond polynomial
category~\texorpdfstring{$\mathcal{O}$}{O}]
\label{rem:infinite-dim-simples}
\index{Category O@Category $\mathcal{O}$!infinite-dimensional simples}
The remaining question at $N \ge 3$
(Remark~\ref{rem:catO-thick-generation-verification}(iii)) is
whether the polynomial subcategory $\mathcal{O}_{\mathrm{poly}}$
exhausts the DK-relevant category~$\mathcal{O}$. If not, the
infinite-dimensional simples $L(\lambda, a)$ at non-polynomial
$\ell$-weights are the precise obstruction.

In the classical setting (without the Yangian spectral parameter),
infinite-dimensional simples at non-integral weight are orthogonal
to all finite-dimensional modules (different blocks of BGG
category~$\mathcal{O}$). The Yangian spectral parameter creates
inter-block extensions: the $R$-matrix poles at
$a - b \in \{0, \pm 1, \ldots, \pm (N-1)\}$ produce nontrivial
$\operatorname{Ext}^1_Y(V_n(b), L(\lambda, a))$ connecting
finite-dimensional and infinite-dimensional modules. Whether these
extensions suffice for thick generation beyond
$\mathcal{O}_{\mathrm{poly}}$ is the residual content of the
generation problem at $N \ge 3$.
\end{remark}

\begin{remark}[$\mathcal{O}_{\mathrm{poly}}$ does not thickly
generate~$\mathcal{O}^{\mathrm{sh}}$]
\label{rem:opoly-not-dense}
\index{prefundamental module!thick generation obstruction}
The negative prefundamental $L^-(b)$ has weight multiplicities
$\dim(L^-)_{-2k} = p(k)$, unbounded as $k \to \infty$.
By definition, every module in~$\mathcal{O}_{\mathrm{poly}}$ has
polynomially bounded weight multiplicities. Therefore
$L^- \notin \mathcal{O}_{\mathrm{poly}}$, and the Baxter
equivariance proved on~$\mathcal{O}_{\mathrm{poly}}$
(Proposition~\ref{prop:baxter-yangian-equivariance}) does
\emph{not} extend to~$L^-$ by thick generation from
$\mathcal{O}_{\mathrm{poly}}$. A separate argument
is needed: either the $\lambda = 0$ specialization
(which still requires the Yangian hyperplane
$b = a - 1/2$) or an
explicit pro-Weyl approximation of~$L^-$ by objects in
$\mathcal{O}_{\mathrm{poly}}$. Since $p(k)$ grows
superpolynomially (Hardy--Ramanujan), no finite-length
resolution of~$L^-$ by evaluation modules exists.
\end{remark}

\medskip

\noindent\textbf{Strategy~I\@: Highest-weight filtration and BGG
resolution.}

\begin{proposition}[BGG resolution criterion for thick generation;
\ClaimStatusProvedHere]
\label{prop:bgg-criterion}
\index{BGG resolution!thick generation}
Thick generation of
$D^b(\mathcal{O}_{Y_\hbar(\mathfrak{sl}_N)})$ by evaluation
modules holds if every Verma module
$M(\lambda, a) \in \mathcal{O}_{Y_\hbar}$ admits a finite
resolution
\begin{equation}\label{eq:verma-eval-resolution}
0 \to E_m \to \cdots \to E_1 \to E_0 \to M(\lambda, a) \to 0
\end{equation}
with each $E_i$ a finite direct sum of evaluation modules
$V_\mu(b)$ \textup{(}$V_\mu$ finite-dimensional\textup{)}.
\end{proposition}

\begin{proof}
If~\eqref{eq:verma-eval-resolution} holds, then
$M(\lambda, a) \in \mathcal{C}_{\mathrm{eval}}$ by iterated cones.
Every simple $L(\lambda, a)$ is a quotient of $M(\lambda, a)$;
the kernel is a module of $\mathcal{O}$ of strictly shorter
composition length. By induction on length, all simples lie in
$\mathcal{C}_{\mathrm{eval}}$.
Lemma~\ref{lem:composition-thick-generation} then gives
$\mathcal{C}_{\mathrm{eval}} = D^b(\mathcal{O})$.
\end{proof}

\begin{remark}[BGG resolution via spectral parameter extensions]
\label{rem:bgg-spectral-parameter}
\index{BGG resolution!spectral parameter}
For dominant integral $\lambda$, the Verma module $M(\lambda, a)$
surjects onto the finite-dimensional evaluation module
$L(\lambda, a)$ with kernel $M(\mu, a)$ at lower weight
$\mu = s_\alpha \cdot \lambda$. Iterating along the Weyl group
orbit produces the classical BGG resolution, all of whose terms are
Verma evaluation modules at the \emph{same} spectral parameter~$a$.

For non-dominant $\lambda$, the classical BGG resolution is
unavailable: $M(\lambda, a)$ is simple at fixed~$a$. The Yangian
structure opens a new channel: evaluation at \emph{different}
spectral parameters. The $R$-matrix intertwiner
\[
R_{V,W}(a - b)\colon V(a) \otimes W(b) \to W(b) \otimes V(a)
\]
has poles at $a - b \in \{0, \pm 1, \ldots, \pm (N-1)\}$
(for fundamental representations of $\mathfrak{sl}_N$).
At these poles, the residues create nontrivial extensions between
evaluation modules at distinct spectral parameters.

For $\mathfrak{sl}_2$: the $R$-matrix pole at $a - b = \pm 1$
produces extensions
\[
0 \to V_n(a) \to E \to V_m(a+1) \to 0
\]
for appropriate $n, m$. Iterating such extensions along a chain
of spectral parameters $a, a+1, a+2, \ldots$ builds complexes
whose totalization approaches Verma-type modules. Whether this
process terminates in a \emph{finite}
resolution~\eqref{eq:verma-eval-resolution} is the key analytical
question. For $\mathfrak{sl}_2$, the Ext groups
$\operatorname{Ext}^n_{\mathcal{O}}(V_m(b), M(\lambda, a))$ are
computable from the $q$-character theory of Hernandez
\cite{Hernandez05}, and finiteness of the resolution follows
from eventual vanishing of these Ext groups at large spectral
parameter separation $|a - b|$.
\end{remark}

\begin{computation}[R-matrix extensions for \texorpdfstring{$Y(\mathfrak{sl}_2)$}{Y(sl_2)}]
\label{comp:sl2-rmatrix-ext}
\index{R-matrix!extensions for sl2@extensions for $\mathfrak{sl}_2$}
\index{Yangian!sl2 R-matrix@$\mathfrak{sl}_2$ $R$-matrix}
We compute the extensions produced by R-matrix poles for
$Y(\mathfrak{sl}_2)$ acting on its fundamental
$2$-dimensional representation $V_1$.

\emph{Step~1: The R-matrix.}
For $V_1(a) \otimes V_1(b)$ the Yang $R$-matrix is
\begin{equation}\label{eq:sl2-rmatrix-fundamental}
R(u) = \frac{u \cdot \mathrm{id} + P}{u + 1}
\;\in\; \operatorname{End}(V_1 \otimes V_1),
\end{equation}
where $u = a - b$ and $P$ is the transposition.
The eigenvalues of~$P$ on $V_1 \otimes V_1 \cong V_2 \oplus V_0$
are $+1$ (on $\mathrm{Sym}^2(V_1) = V_2$) and $-1$
(on $\Lambda^2(V_1) = V_0$), giving
\[
R(u)\big|_{V_2} = \frac{u+1}{u+1} = 1, \qquad
R(u)\big|_{V_0} = \frac{u-1}{u+1}.
\]
The unique pole is at $u = -1$ (i.e.\ $a = b - 1$), with residue
projecting onto the $V_0$ summand.

\emph{Step~2: The extension at the pole.}
At $a - b = -1$, the intertwiner
$R_{V_1,V_1}(-1)\colon V_1(b-1) \otimes V_1(b) \to
V_1(b) \otimes V_1(b-1)$
has a pole. The residue
\[
\Res_{u = -1} R(u) = \frac{-\mathrm{id} + P}{1}
\bigg|_{\text{pole}} = \pi_{V_0}
\]
(the projector onto the antisymmetric line $V_0 \subset
V_1 \otimes V_1$) defines a nontrivial extension
\begin{equation}\label{eq:sl2-ext-fundamental}
0 \;\to\; V_0(b-\tfrac{1}{2})
\;\to\; E
\;\to\; V_2(b-\tfrac{1}{2})
\;\to\; 0
\end{equation}
in $\mathcal{O}_{Y(\mathfrak{sl}_2)}$, where $E$ is the image
of $V_1(b-1) \otimes V_1(b)$ at the specialization $a = b - 1$.
This gives
\begin{equation}\label{eq:sl2-ext1}
\dim \operatorname{Ext}^1_{\mathcal{O}}(V_2(c),\, V_0(c)) = 1
\end{equation}
for all~$c$: a single extension class, classified by the
R-matrix pole.

\emph{Step~3: Higher representations.}
For $V_n(a) \otimes V_m(b)$ with $n + m \le 2$, the
$R$-matrix poles occur at $a - b \in \{0, \pm 1\}$ (for
$\mathfrak{sl}_2$, the only fundamental is $V_1$, and the
general R-matrix factors through fundamental $R$-matrices).
The tensor product decomposition
$V_n(a) \otimes V_1(b)$ at the pole $a - b = -(n+1)/2$
produces
\[
0 \to V_{n-1}(c') \to V_n(a) \otimes V_1(b)\big|_{a = b - (n+1)/2}
\to V_{n+1}(c'') \to 0
\]
for appropriate shifted spectral parameters $c', c''$.
Iterating: starting from any finite-dimensional $V_n(a)$,
one can build a complex
\[
\cdots \to V_n(a) \otimes V_1(a+s_2) \otimes V_1(a+s_3)
\to V_n(a) \otimes V_1(a+s_1) \to V_n(a) \to M(\lambda, a) \to 0
\]
that resolves the Verma module $M(\lambda, a)$ by iterated tensor
products with evaluation modules at shifted spectral parameters
$s_1, s_2, \ldots$.

\emph{Step~4: Finiteness for $\mathfrak{sl}_2$.}
For $\mathfrak{sl}_2$, each step reduces the highest weight of the
kernel by exactly~$2$ (one $R$-matrix pole). A Verma module
$M(\lambda, a)$ at dominant integral weight~$\lambda$ therefore
admits a resolution of length exactly $\lambda + 1$ (the length of
the Weyl group orbit), recovering the classical BGG resolution
enriched with spectral parameters.

For non-integral $\lambda$, the chain of spectral parameter
extensions does not terminate after finitely many steps: there is no
lowest weight in the $\mathfrak{sl}_2$ Verma module when $\lambda$
is not a non-negative integer. This exhibits the precise
obstruction: Strategy~I succeeds exactly on the dominant integral
(= polynomial) subcategory, consistent with the proved
Theorem~\ref{thm:catO-thick-generation} for
$\mathcal{O}_{\mathrm{poly}}$.
\end{computation}

\begin{proposition}[Heart-capture criterion; \ClaimStatusProvedHere]
\label{prop:heart-capture-criterion}
\index{thick generation!heart-capture criterion}
Let $\cA$ be a finite-length abelian category and
$\mathcal{T} \subset D^b(\cA)$ a thick triangulated subcategory.
If every simple object of~$\cA$ lies in~$\mathcal{T}$, then
$\mathcal{T} = D^b(\cA)$.
\end{proposition}
\begin{proof}
By induction on Jordan--H\"older length, every object
of~$\cA$ lies in~$\mathcal{T}$: if $0 \to K \to M \to S \to 0$
is exact with $S$ simple and $\operatorname{length}(K) < \operatorname{length}(M)$,
then $K \in \mathcal{T}$ by the inductive hypothesis and $S \in \mathcal{T}$
by assumption, so the cone of $K \to M$ places $M$ in~$\mathcal{T}$.
For the full derived category, every bounded complex
$X^\bullet$ is built from its cohomology objects by canonical
truncation triangles: $\tau_{\leq n-1}X \to \tau_{\leq n}X \to H^n(X)[-n]$,
so by induction on cohomological amplitude, $X^\bullet \in \mathcal{T}$.
\end{proof}

\begin{proposition}[Standard-capture criterion; \ClaimStatusProvedHere]
\label{prop:standard-capture-criterion}
\index{thick generation!standard-capture criterion}
\index{highest-weight category!standard-capture}
Let $\cA$ be a finite-length highest-weight category with finite
weight poset~$\Lambda$ and standard objects
$\{\Delta(\lambda)\}_{\lambda \in \Lambda}$.
Let $\mathcal{T} \subset D^b(\cA)$ be a thick subcategory.
If every standard object $\Delta(\lambda) \in \mathcal{T}$, then
$\mathcal{T} = D^b(\cA)$.
\end{proposition}
\begin{proof}
Induct on the weight poset. For a maximal $\lambda$,
$\Delta(\lambda) = L(\lambda)$, so $L(\lambda) \in \mathcal{T}$.
For general~$\lambda$, the quotient
$\Delta(\lambda) \twoheadrightarrow L(\lambda)$ has kernel
filtered by simples $L(\mu)$ with $\mu > \lambda$, all of which
lie in~$\mathcal{T}$ by induction.
Since $\Delta(\lambda) \in \mathcal{T}$ and the kernel is in~$\mathcal{T}$,
the cone construction places $L(\lambda) \in \mathcal{T}$.
Heart-capture
(Proposition~\ref{prop:heart-capture-criterion}) then gives
$\mathcal{T} = D^b(\cA)$.
\end{proof}

\begin{corollary}[Sectorwise localizing generation; \ClaimStatusProvedHere]
\label{cor:sectorwise-localizing-generation}
\index{thick generation!sectorwise localizing}
\index{truncation sector!localizing generation}
Suppose a completed category
$\widehat{\mathcal{O}} = \bigcup_\sigma \mathcal{O}_\sigma^{\mathrm{trunc}}$
is a filtered union of finite truncation sectors, each a
finite-length highest-weight category. If in each sector the
evaluation objects $\mathrm{Ev}_\sigma$ generate all standard objects
under thick closure, then
$\operatorname{Loc}\langle \mathrm{Ev} \rangle = \widehat{\mathcal{O}}$.
\end{corollary}
\begin{proof}
In each finite sector~$\mathcal{O}_\sigma^{\mathrm{trunc}}$,
standards lie in~$\operatorname{thick}\langle \mathrm{Ev}_\sigma \rangle$
by hypothesis, so standard-capture
(Proposition~\ref{prop:standard-capture-criterion}) gives
$\operatorname{thick}\langle \mathrm{Ev}_\sigma \rangle =
D^b(\mathcal{O}_\sigma^{\mathrm{trunc}})$.
Since every object of $\widehat{\mathcal{O}}$ belongs to some sector
and localizing closure is closed under filtered colimits, the result
follows.
\end{proof}

\begin{remark}[Assessment and interaction of the four strategies]
\label{rem:cat-O-strategies-assessment}
\index{thick generation!strategy comparison}
\begin{center}
\renewcommand{\arraystretch}{1.2}
\begin{tabular}{lll}
\textbf{Strategy} & \textbf{Reduces to} & \textbf{Scope} \\
\hline
IV (sectorwise) & Current algebra $E_1$ computation
 & All types \\
III (KL reduction) & Finite global dimension +
 Verma $\in \mathcal{C}_{\mathrm{eval}}$ & Type $A$ \\
II (t-structure) & \textbf{Proved}: $\mathfrak{sl}_2$ + $\mathcal{O}_{\mathrm{poly}}$
 (Thm.~\ref{thm:catO-thick-generation}) & Type $A$ \\
I (BGG) & Finite evaluation resolution of Verma modules
 & All types \\
\end{tabular}
\end{center}
Strategy~II is proved (Theorem~\ref{thm:catO-thick-generation}); the residual question is whether $\mathcal{O}_{\mathrm{poly}}$ exhausts the DK-relevant category (Remark~\ref{rem:infinite-dim-simples}). Strategy~IV resolves DK-2/3 on the evaluation-generated core for all simple types via sectorwise convergence (Theorem~\ref{thm:h-level-factorization-kd}, Corollary~\ref{cor:dk23-all-types}), bypassing thick generation. The remaining bottleneck is the factorization bar cohomology of $U(\mathfrak{g}[t])$, which interacts with the Loday--Quillen--Tsygan theorem.
\end{remark}

\begin{conjecture}[Sectorwise DK extension to the completed/coderived MC3 domain for all types;
\ClaimStatusConjectured]\label{conj:mc3-sectorwise-all-types}
\index{thick generation!sectorwise finiteness!all types}
For every simple~$\mathfrak{g}$, the sectorwise root-weight
decomposition and loop filtration of
Proposition~\ref{prop:yangian-bar-loop-weight} extend from the
evaluation-generated core to the completed/coderived enlargement
used by the MC3 completion problem. Assume that,
on this enlarged domain, the associated spectral sequence
stabilizes at each tridegree, retains sectorwise sub-exponential
growth, and is compatible with factorization descent on
$\operatorname{Ran}(X)$. Then the sectorwise argument of
Theorem~\textup{\ref{thm:h-level-factorization-kd}} extends the
proved DK-2/3 equivalence from the evaluation-generated core to the
completed/coderived factorization domain relevant to the remaining
MC3 packets at all simple types, bypassing ordinary thick generation
as the route to that enlarged domain.

\emph{The proved part of Strategy~IV is already settled on the
evaluation-generated core by
Theorem~\textup{\ref{thm:h-level-factorization-kd}} and
Corollary~\textup{\ref{cor:dk23-all-types}}. The residual question
is whether the same sectorwise control survives the
completed/coderived enlargement.}
\end{conjecture}

\begin{proposition}[Sub-exponential growth of the \texorpdfstring{$E_1$}{E_1} page;
\ClaimStatusProvedHere]
\label{prop:lqt-e1-subexponential-growth}
\index{Loday--Quillen--Tsygan theorem!$E_1$ growth rate}
\index{partition asymptotics!current algebra}
For every simple Lie algebra~$\mathfrak{g}$ of rank~$r$,
the $E_1$ page of the Strategy~IV spectral sequence satisfies
\[
\dim E_1^{0,p}(\mathfrak{g}[t])
\;\sim\;
C(\mathfrak{g})\,p^{-3/4}\,
\exp\!\Bigl(\pi\sqrt{\tfrac{rp}{12}}\Bigr)
\qquad (p \to \infty),
\]
where $C(\mathfrak{g})$ depends on the exponents of~$\mathfrak{g}$.
In particular:
\begin{enumerate}[label=\textup{(\alph*)}]
\item the growth is sub-exponential: $\dim E_1^{0,p} < c^p$
 for every $c > 1$ and all sufficiently large~$p$;
\item the growth is \emph{not} polynomial: for every~$N$,
 $\dim E_1^{0,p} / p^N \to \infty$;
\item the leading growth constant $\pi\sqrt{r/12}$ depends
 only on the rank~$r$, not on the specific exponents.
\end{enumerate}
\end{proposition}

\begin{proof}
The Feigin--Tsygan formulation of the
Loday--Quillen--Tsygan theorem gives
\[
H^*(\mathfrak{g}[t], k)
\;\cong\;
\Lambda\bigl(\xi_{i,n} \mid
 1 \le i \le r,\; n \ge 0\bigr),
\qquad
\deg(\xi_{i,n}) = 2e_i + 1 + 2n.
\]
Hence $\dim E_1^{0,p}$ equals the number of subsets~$S$ of
$\{(i,n) : 1 \le i \le r,\, n \ge 0\}$ satisfying
$\sum_{(i,n) \in S} (2e_i + 1 + 2n) = p$.
The generating function is
\[
P_{\mathfrak{g}}(q)
= \sum_{p \ge 0} \dim E_1^{0,p}\, q^p
= \prod_{i=1}^{r} \prod_{n=0}^{\infty}
 \bigl(1 + q^{2e_i + 1 + 2n}\bigr).
\]
Each factor $\prod_{n=0}^{\infty}(1 + q^{a_i + 2n})$ with
$a_i = 2e_i + 1$ counts partitions into distinct parts of an
arithmetic progression with common difference~$2$. Applying
the saddle-point method, the logarithmic singularity at
$q = 1$ gives
\[
\log P_{\mathfrak{g}}(e^{-t})
\;\sim\; \frac{r\pi^2}{24\,t}
\qquad (t \to 0^+),
\]
since $\int_0^\infty \log(1 + e^{-x})\,dx = \pi^2/12$ and the $r$
independent progressions contribute additively (the starting
points~$a_i$ affect only subleading terms). The standard
Tauberian transfer yields the claimed asymptotic.
Part~(a) follows because the exponent $\pi\sqrt{rp/12}$ grows
as~$\sqrt{p}$, slower than any linear function~$cp$;
part~(b) from $\sqrt{p}$ growing faster than $N\log p$;
part~(c) from the leading term $r\pi^2/(24t)$ depending only
on~$r$.
\end{proof}

\begin{computation}[Finite-window sub-exponential growth for all types]
\label{comp:lqt-e1-growth-verification}
\index{Loday--Quillen--Tsygan theorem!finite-window growth}
The sub-exponential growth rate of
Proposition~\ref{prop:lqt-e1-subexponential-growth} is visible in
the finite enumeration for all simple types through $p=500$:
\begin{itemize}
\item Manuscript tables for $\mathfrak{sl}_2$, $\mathfrak{sl}_3$,
 $\mathfrak{sp}_4$ reproduced exactly through $p = 9$;
\item At $p = 200$: $\dim E_1^{0,200}(\mathfrak{sl}_2) = 159{,}733$,
 $\dim E_1^{0,200}(\mathfrak{sl}_3) = 98{,}874{,}604$
 (not polynomial);
\item Observed growth constants converge to $\pi\sqrt{r/12}$
 within $8\%$ at $p = 500$ for all types including the
 exceptional algebras $G_2$, $F_4$, $E_{6,7,8}$;
\item Same-rank algebras ($A_2$, $B_2$, $G_2$: all rank~$2$)
 confirm identical leading constant, different subleading;
\item Departure points from $A_1$: $p = 5$ for $A_2$,
 $p = 7$ for $C_2$, $p = 11$ for $G_2$ (matching manuscript).
\end{itemize}
Sectorwise finiteness for all simply-laced lattice
VOAs ($A_1$--$A_3$, $D_4$--$D_5$, $E_6$--$E_8$) holds at the
lattice-bar sector level in the same finite window, and the
sub-exponential versus exponential growth discrimination is the
comparison of $\log d_n$ with the two models $c\cdot n$ and
$C\sqrt n$.
\end{computation}

\subsubsection*{Resolution: the evaluation-generated core}

\begin{theorem}[Evaluation-generated core identification, type~\texorpdfstring{$A$}{A};
\ClaimStatusConditional]
\label{thm:eval-core-identification}
\index{thick generation!evaluation-generated core|textbf}
\index{Drinfeld--Kohno!derived!evaluation-generated core}
For $\fg = \mathfrak{sl}_N$ with $N \ge 2$, the thick triangulated
subcategory of $D^b(\mathcal{O}_{Y_\hbar(\mathfrak{sl}_N)})$
generated by evaluation modules satisfies
\begin{equation}\label{eq:eval-core}
\mathcal{D}_\hbar^{\mathrm{eval}}
\;=\; D^b\bigl(\operatorname{Rep}_{\mathrm{fd}}
 (Y_\hbar(\mathfrak{sl}_N))\bigr).
\end{equation}
In particular, the module Koszul duality functor restricts to a
braided monoidal equivalence
\begin{equation}\label{eq:dk-eval-core}
\Phi\colon
D^b\bigl(\operatorname{Rep}_{\mathrm{fd}}(Y_\hbar(\mathfrak{sl}_N))\bigr)
\;\xrightarrow{\;\sim\;}\;
D^b\bigl(\operatorname{Rep}_{\mathrm{fd}}
 (Y_{-\hbar}(\mathfrak{sl}_N))\bigr)^{\mathrm{op}},
\end{equation}
and DK-2/DK-3 are resolved for the evaluation-generated core
at all~$N$.
\end{theorem}

\begin{proof}
\emph{$(\supseteq)$:}
Every finite-dimensional simple $Y_\hbar(\mathfrak{sl}_N)$-module
lies in $\mathcal{D}_\hbar^{\mathrm{eval}}$ by
Proposition~\ref{prop:dk2-thick-generation-typeA}.
Since $\operatorname{Rep}_{\mathrm{fd}}$ is Artinian, every
object has a finite composition series with simple factors in
$\mathcal{D}_\hbar^{\mathrm{eval}}$.
Lemma~\ref{lem:composition-thick-generation} then gives
$D^b(\operatorname{Rep}_{\mathrm{fd}}) \subset
\mathcal{D}_\hbar^{\mathrm{eval}}$.

\emph{$(\subseteq)$:}
$\operatorname{Rep}_{\mathrm{fd}}$ contains all evaluation modules
and is a Serre subcategory of $\mathcal{O}_{Y_\hbar}$ (closed under
subobjects, quotients, and extensions, since a finite extension of
finite-dimensional modules is finite-dimensional). Therefore
$D^b(\operatorname{Rep}_{\mathrm{fd}})$ is a thick subcategory of
$D^b(\mathcal{O}_{Y_\hbar})$. Since it contains all evaluation
modules, $\mathcal{D}_\hbar^{\mathrm{eval}} \subset
D^b(\operatorname{Rep}_{\mathrm{fd}})$.

\emph{DK equivalence:}
The functor $\Phi$ sends evaluation modules at parameter
$\hbar$ to evaluation modules at parameter $-\hbar$
(Proposition~\ref{prop:yangian-module-koszul}). By the core
identification~\eqref{eq:eval-core}, $\Phi$ restricts to a functor
between the two sides of~\eqref{eq:dk-eval-core}.
Full faithfulness and essential surjectivity hold on the
evaluation-generated subcategory by
Theorem~\ref{thm:derived-dk-yangian}. Since
\eqref{eq:eval-core} identifies $D^b(\operatorname{Rep}_{\mathrm{fd}})$
with the thick closure of that subcategory, the same exactness and
thick-closure argument extends the equivalence to all of
$D^b(\operatorname{Rep}_{\mathrm{fd}})$. Strong monoidality then
extends from the evaluation-generated subcategory to the full thick
closure by Lemma~\ref{lem:monoidal-thick-extension}.
\end{proof}

\subsubsection*{Beyond finite dimensions: the shifted-prefundamental package}

\begin{theorem}[Baxter exact triangles on
\texorpdfstring{$\mathcal{O}_{\mathrm{poly}}$}{Opoly};
\ClaimStatusProvedHere]
\label{thm:baxter-exact-triangles-opoly}
\index{Baxter relation!exact triangle lift|textbf}
\index{shifted Yangian!Baxter exact triangles}
For $\fg = \mathfrak{sl}_2$ and generic spectral parameters,
the three-term TQ relations of Zhang~\cite{Zhang24} lift from
$K_0(\mathcal{O})$ to
$Y(\mathfrak{sl}_2)$-equivariant short exact sequences
on~$\mathcal{O}_{\mathrm{poly}}$.
Specifically, for each Verma module $M(\lambda, b)$ and
evaluation module~$V_1(a)$ with the spectral constraint
$b = a - (\lambda + 1)/2$
\textup{(}Proposition~\textup{\ref{prop:baxter-yangian-equivariance})},
the sequence
\[
0 \to M(\lambda - 1, b') \to V_1(a) \otimes M(\lambda, b)
 \to M(\lambda + 1, b'') \to 0
\]
is a natural short exact sequence of $Y(\mathfrak{sl}_2)$-modules,
functorial on~$\mathcal{O}_{\mathrm{poly}}$
\textup{(}Corollary~\textup{\ref{cor:baxter-naturality-opoly})}.
The $K_0$-image recovers Zhang's formula; the derived
octahedral compatibility is the derived form of the same functorial
short exact sequences under iterated TQ composition.
\end{theorem}

\begin{proof}
The Yangian-equivariant singular vector~$w_\lambda$
(Proposition~\ref{prop:baxter-yangian-equivariance}) generates
a sub-Verma $M(\lambda - 1) \hookrightarrow V_1(a) \otimes M(\lambda)$.
The quotient is necessarily $M(\lambda + 1)$ by weight comparison.
Naturality on $\mathcal{O}_{\mathrm{poly}}$ follows from
Corollary~\ref{cor:baxter-naturality-opoly}.
The octahedral axiom for composed maps $V_1 \otimes V_1 \otimes M$
follows from the Clebsch--Gordan decomposition
$V_1 \otimes V_1 \simeq V_0 \oplus V_2$ and the functoriality of
the SES with respect to the tensor product associator.
\end{proof}

\begin{theorem}[Baxter exact triangles on shifted envelope
\texorpdfstring{$\mathcal{O}^{\mathrm{sh}}_{\leq 0}$}{Osh};
\ClaimStatusProvedHere]
\label{thm:baxter-exact-triangles}
\index{Baxter relation!exact triangle lift!shifted envelope}
Theorem~\textup{\ref{thm:baxter-exact-triangles-opoly}} extends
from $\mathcal{O}_{\mathrm{poly}}$ to the full anti-dominant
shifted envelope
$\mathcal{O}^{\mathrm{sh}}_{\leq 0}
:= \bigoplus_{\mu \leq 0} \mathcal{O}_\mu$
for $\fg = \mathfrak{sl}_N$ at general~$N$.
Concretely, for each node~$i \in I$ and finite-dimensional companion
module $M_{k,x}^{(i)}$, there is a distinguished triangle
\[
A^-_{k,x,y} \to M_{k,x}^{(i)} \otimes L(\Psi_{i,x}/\Psi_{i,y})
\to A^+_{k,x,y} \to
\]
whose $K_0$-class recovers the formula of Zhang~\cite{Zhang24},
and which is functorial on~$\mathcal{O}^{\mathrm{sh}}_{\leq 0}$.
\end{theorem}

\begin{proof}
For type~$A$, Proposition~\ref{prop:categorical-cg-typeA} upgrades the
character identity of
Proposition~\ref{prop:prefundamental-clebsch-gordan} to an actual
module decomposition at generic spectral parameters:
\[
V_{\omega_i}(a) \otimes L^-_i(b)
\;\cong\;
\bigoplus_{\mu \in \operatorname{wt}(V_{\omega_i})}
L^-_i(\mathrm{shift} = \mu).
\]
The singular vector of
Proposition~\ref{prop:baxter-yangian-equivariance} identifies the
simple-root summand on the Baxter locus, and the block separation of
shifted prefundamental modules forces the filtration to split at the
module level. The resulting short exact sequences in
$\mathcal{O}^{\mathrm{sh}}_{\leq 0}$ yield the required distinguished
triangles. Their Grothendieck classes recover Zhang's TQ relation
because the same decomposition computes the $K_0$-identity, and
functoriality follows from the blockwise construction together with
Theorem~\ref{thm:baxter-exact-triangles-opoly}.
\end{proof}

\begin{remark}[Type-A resolution of Baxter exact triangles]
\label{rem:baxter-exact-triangles-typeA}
In type~$A$ ($\mathfrak{sl}_N$ for all $N \geq 2$),
Theorem~\ref{thm:baxter-exact-triangles} is proved at generic
spectral parameters by the categorical CG decomposition
(Proposition~\ref{prop:categorical-cg-typeA}): the module-level
splitting $V_{\omega_i}(a) \otimes L^-_i(b) \cong
\bigoplus L^-_i(b_\mu)$ provides the Baxter companion modules as
actual $Y(\mathfrak{sl}_N)$-submodules, and the associated short
exact sequences give the required derived exact triangles.
The proof combines three ingredients:
\textup{(1)}~the singular vector construction
(Proposition~\ref{prop:baxter-yangian-equivariance}),
\textup{(2)}~the character identity
(Proposition~\ref{prop:prefundamental-clebsch-gordan}),
and \textup{(3)}~the block separation of shifted prefundamental
modules in~$\mathcal{O}^{\mathrm{sh}}$, which forces the
character-level filtration to split as modules.
The generic-parameter condition suffices for the MC3 generation
arguments, since compact generation requires exact triangles only
at a cofinal family.
This extension to arbitrary simple $\mathfrak{g}$ is now complete only
for the categorical CG packet: the multiplicity-free $\ell$-weight
property (Chari--Moura) replaces the minuscule hypothesis, yielding
the categorical CG closure for all simple types
\textup{(}Theorem~\ref{thm:categorical-cg-all-types}\textup{)}.
\end{remark}

\begin{proposition}[Yangian equivariance of the Baxter singular vector;
\ClaimStatusProvedHere]
\label{prop:baxter-yangian-equivariance}
\index{Baxter relation!Yangian equivariance}
\index{spectral constraint!Baxter singular vector}
Let $V_1(a)$ be the $2$-dimensional evaluation module at spectral
parameter~$a$, and $M(\lambda, b)$ the Verma module of highest
weight~$\lambda$ at spectral parameter~$b$, both
for~$Y(\mathfrak{sl}_2)$. The singular vector
\[
 w_\lambda
 \;=\;
 \lambda \cdot (v_- \otimes v_\lambda)
 \;-\;
 v_+ \otimes f \cdot v_\lambda
 \;\in\; V_1(a) \otimes M(\lambda, b)
\]
(which generates a sub-Verma $M(\lambda-1)$ at the
$\mathfrak{sl}_2$ level) is annihilated by the Yangian coproduct
$\Delta(E)$ if and only if
\begin{equation}\label{eq:baxter-spectral-constraint}
 b \;=\; a - \frac{\lambda + 1}{2}.
\end{equation}
Under this constraint, $w_\lambda$ is a Yangian highest-weight
vector and the short exact sequence
\begin{equation}\label{eq:baxter-ses-yangian}
 0 \;\to\; M(\lambda - 1,\, b')
 \;\to\; V_1(a) \otimes M(\lambda,\, b)
 \;\to\; M(\lambda + 1,\, b'')
 \;\to\; 0
\end{equation}
is $Y(\mathfrak{sl}_2)$-equivariant, with spectral parameters
$b'$, $b''$ determined by the $\Delta(H)$-eigenvalues on $w_\lambda$
and on the quotient highest-weight vector respectively.
\end{proposition}

\begin{proof}
Direct computation using the Drinfeld coproduct
$\Delta(E) = a\,e \otimes 1 + 1 \otimes E
+ \tfrac{1}{2}(h \otimes e - e \otimes h)$
on $V_1(a) \otimes M(\lambda, b)$, where $E$ acts on $M(\lambda, b)$
via $E = b\,e$ (evaluation at parameter~$b$).

The weight-$(\lambda - 1)$ subspace of
$V_1(a) \otimes M(\lambda, b)$
is spanned by $u_1 = v_+ \otimes f \cdot v_\lambda$ and
$u_2 = v_- \otimes v_\lambda$. Computing:
\begin{align*}
 \Delta(E) \cdot u_1
 &= \bigl(b\lambda + \tfrac{\lambda}{2}\bigr)
 (v_+ \otimes v_\lambda),\\
 \Delta(E) \cdot u_2
 &= \bigl(a - \tfrac{\lambda}{2}\bigr)
 (v_+ \otimes v_\lambda),
\end{align*}
where the key simplification uses
$E \cdot v_\lambda = b \cdot e \cdot v_\lambda = 0$
and $E \cdot f \cdot v_\lambda = b\lambda \cdot v_\lambda$.
For $w_\lambda = \lambda \cdot u_2 - u_1$:
\[
 \Delta(E) \cdot w_\lambda
 \;=\;
 \lambda\bigl(a - \tfrac{\lambda}{2}\bigr)
 - b\lambda - \tfrac{\lambda}{2}
 \;=\;
 \lambda\Bigl(a - b - \frac{\lambda + 1}{2}\Bigr),
\]
which vanishes if and only if $b = a - (\lambda + 1)/2$.
An analogous computation for $\Delta(H)$ confirms that $w_\lambda$
is an $H$-eigenvector with eigenvalue determining the spectral
parameter~$b'$ of the sub-Verma.
\end{proof}

\begin{corollary}[Naturality on
$\mathcal{O}_{\mathrm{poly}}$; \ClaimStatusProvedHere]
\label{cor:baxter-naturality-opoly}
Under the spectral
constraint~\eqref{eq:baxter-spectral-constraint}, the
assignment
$M(\lambda) \mapsto
[\textup{SES}~\eqref{eq:baxter-ses-yangian}]$
defines a natural exact sequence of functors on
$\mathcal{O}_{\mathrm{poly}}(Y(\mathfrak{sl}_2))$.
\end{corollary}

\begin{proof}
On $\mathcal{O}_{\mathrm{poly}}$, Verma modules are the simple
objects and $\Hom(M(\lambda), M(\mu)) = 0$ for $\lambda \neq \mu$,
$\End(M(\lambda)) = \mathbb{C}$. Since all Hom spaces are at most
$1$-dimensional, naturality of the SES maps with respect to
morphisms in $\mathcal{O}_{\mathrm{poly}}$ is automatic from the
existence of the $Y(\mathfrak{sl}_2)$-equivariant SES at each
weight
(Proposition~\ref{prop:baxter-yangian-equivariance}).
\end{proof}

\begin{remark}[From R-matrix poles to Baxter equivariance]
\label{rem:rmatrix-baxter-connection}
\index{Baxter relation!R-matrix connection}
The spectral constraint $b = a - (\lambda+1)/2$ is the
\emph{Verma module analogue} of the R-matrix pole condition
$a - b = (n+1)/2$ for finite-dimensional tensor products
$V_1(a) \otimes V_n(b)$
(Computation~\ref{comp:sl2-rmatrix-ext}, Step~3). Setting
$\lambda = n$, the two constraints coincide, exhibiting the Baxter
SES as the analytic continuation of the R-matrix extension from
finite-dimensional to infinite-dimensional highest-weight modules.
This uniformity (the same operadic structure, R-matrix poles in
$\FM_2(\mathbb{C})$, governing both the evaluation-generated core
and the Verma-module Baxter project) is the local rank-$1$ model for
Theorem~\ref{thm:baxter-exact-triangles} and the first step
toward the non-type-$A$ categorical lift (H4 in the concordance): the
naturality established in
Corollary~\ref{cor:baxter-naturality-opoly} resolves the
H4~gap on $\mathcal{O}_{\mathrm{poly}}$, and
Theorem~\ref{thm:baxter-exact-triangles} closes the type-$A$
shifted-envelope step. The remaining open problem is the non-type-$A$
categorical lift together with the compact/completed MC3 comparison
packet.
\end{remark}

\begin{theorem}[Shifted-prefundamental generation on the shifted envelope;
\ClaimStatusConditional]
\label{thm:shifted-prefundamental-generation}
\index{shifted Yangian!prefundamental generation}
\index{shifted envelope!prefundamental generation}
\index{E1-chiral!thick generation}
For $\mathfrak{g} = \mathfrak{sl}_N$ ($N \geq 2$), in the separated
pro-completed shifted category
$\widehat{D}(\mathcal{O}^{\mathrm{sh}}_{\leq 0}(Y(\mathfrak{g})))$,
assume that the weight filtration is sectorwise complete and that the
standard-object spectral sequence stabilizes after finitely many pages in
each fixed weight sector. Then
every standard object of
$\mathcal{O}^{\mathrm{sh}}_{\leq 0}(Y(\mathfrak{g}))$ lies in the
thick idempotent-complete closure of
\[
\mathcal{G}_{\mathrm{Bax}}
\;:=\;
\{V_{\omega_i}(a) \mid i = 1, \ldots, N{-}1,\; a \in \mathbb{C}\}
\;\cup\;
\Bigl\{L^-_{i,a-\frac{1}{2}}
 \,\Bigm|\, i = 1, \ldots, N{-}1,\; a \in \mathbb{C}\Bigr\},
\]
where $V_{\omega_i}(a)$ are the fundamental evaluation modules
and $L^-_{i,a}$ are the negative prefundamental modules of
Hern'andez--Jimbo~\cite{HJZ25}, one per node of the
Dynkin diagram. Equivalently, the same statement holds for
\[
\mathcal{G}_{\mathrm{shift}}
\;:=\;
\{V_{\omega_i}(a) \mid i = 1, \ldots, N{-}1,\; a \in \mathbb{C}\}
\;\cup\;
\{L^-_{i,b} \mid i = 1, \ldots, N{-}1,\; b \in \mathbb{C}\},
\]
since $a \mapsto a - \tfrac{1}{2}$ is a bijection on~$\mathbb{C}$.
In particular, the image of
$D^b(\mathcal{O}^{\mathrm{sh}}_{\leq 0}(Y(\mathfrak{g})))$ inside
$\widehat{D}(\mathcal{O}^{\mathrm{sh}}_{\leq 0}(Y(\mathfrak{g})))$
is contained in the localizing subcategory generated by
$\mathcal{G}_{\mathrm{Bax}}$.
If $\mathcal{D}^{\mathrm{comp}}$ is a presentable stable
completed enhancement of
$\mathcal{O}^{\mathrm{sh}}_{\leq 0}(Y(\mathfrak{g}))$ whose compact
objects are the thick idempotent-complete closure of the image of
$D^b(\mathcal{O}^{\mathrm{sh}}_{\leq 0}(Y(\mathfrak{g})))$, then
\[
\mathcal{D}^{\mathrm{comp}}
\;=\;
\operatorname{Loc}\langle \mathcal{G}_{\mathrm{Bax}} \rangle.
\]
\end{theorem}

\begin{proof}
Four steps.

\smallskip\noindent\textit{Step~1 (Evaluation core).}
By the proved finite-dimensional thick-generation theorem
\textup{(}Proposition~\ref{prop:dk2-thick-generation-typeA}\textup{)},
the fundamental evaluation modules $\{V_{\omega_i}(a)\}$ compactly
generate $\operatorname{Rep}_{\mathrm{fd}}(Y(\mathfrak{g}))$ at
all~$N$.

\smallskip\noindent\textit{Step~2 (Prefundamental stability on the
Baxter hyperplane).}
Fix a node $i \in \{1, \ldots, N{-}1\}$ of the $A_{N-1}$~Dynkin diagram,
a spectral parameter $b \in \mathbb{C}$, and set
$a := b + \tfrac{1}{2}$. Let $L^-_{i}(b)$ denote the corresponding negative
prefundamental. The singular vector
\begin{equation}\label{eq:baxter-sv-prefundamental}
 w_0 \;=\; v_{\omega_i} \otimes f_i \cdot v_0
 \;\in\; V_{\omega_i}(a) \otimes L^-_i(b)
\end{equation}
(weight $\omega_i - \alpha_i$ in the first factor, weight
$-\alpha_i$ in the second) is annihilated by
$\Delta(E_j)$ for \emph{every} $j = 1, \ldots, N{-}1$ provided
\[
 b \;=\; a - \frac{1}{2}.
\]
The two cases:
\begin{itemize}[nosep]
\item $j = i$: Proposition~\ref{prop:baxter-yangian-equivariance}
 specialized to $\lambda = 0$ gives the actual spectral
 constraint $b = a - 1/2$, and on this hyperplane
 $\Delta(E_i)\cdot w_0 = 0$.
\item $j \neq i$: the commutator $[e_j, f_i] = 0$
 (Chevalley--Serre for distinct nodes in any simple
 Lie algebra), and $e_j \cdot v_{\omega_i} = 0$
 (highest-weight property of $V_{\omega_i}$), so every
 term vanishes.
\end{itemize}
\noindent
This is \textbf{rank-independent}: it uses only the
Chevalley--Serre presentation of~$\mathfrak{g}$,
not any $\mathfrak{sl}_2$-specific formula.
The Baxter constraint at $\lambda = 0$ is therefore exact and not
vacuous: the simple-root companion exists only on the codimension-$1$
locus $b = a - 1/2$. This does not shrink the prefundamental seed. The
projection
\[
\{(a,b) \in \mathbb{C}^2 \mid b = a - \tfrac{1}{2}\}
\longrightarrow \mathbb{C},
\qquad
(a,b) \longmapsto b
\]
is a bijection, so every negative prefundamental $L^-_{i}(b)$ occurs on
the Baxter locus, with unique companion evaluation parameter
$a = b + \tfrac{1}{2}$.

For generic parameters, Theorem~\ref{thm:baxter-exact-triangles} and
Proposition~\ref{prop:categorical-cg-typeA} give the required exact
triangles and shifted prefundamental decompositions on a cofinal family
of such pairs. Since thick closures are closed under cones, shifts,
and direct summands, every shifted prefundamental companion needed in
the anti-dominant shifted envelope lies in
$\operatorname{thick}\langle \mathcal{G}_{\mathrm{Bax}}\rangle$.

\smallskip\noindent\textit{Step~3 (Verma from pro-completion).}
By Theorem~\ref{thm:pro-weyl-recovery}, every ordinary standard module
$M(\Psi)$ of~$Y(\mathfrak{g})$ satisfies
$M(\Psi) \simeq R\varprojlim_m W_m$, where the Weyl
truncations~$W_m$ are finite-dimensional and the inverse system is
Mittag--Leffler ($R^1 \varprojlim = 0$). Each $W_m$ lies in
$\operatorname{thick}\langle\{V_{\omega_i}\}\rangle$ by Step~1.
Hence every ordinary standard lies in
$\operatorname{thick}\langle \mathcal{G}_{\mathrm{Bax}}\rangle$
inside the separated completion.

\smallskip\noindent\textit{Step~4 ($E_1$~compatibility and passage to the
shifted envelope).}
The $E_1$-braiding $R(u)$ is preserved throughout:
the $I$-adic filtration on~$Y(\mathfrak{g})$ is weight-graded
(PBW level filtration), hence $R$-compatible.
The sectorwise stabilization is part of the ambient hypothesis in the
statement, and the Mittag--Leffler vanishing in Step~3 depends only on the
weight filtration, not on the braiding data.
Under Koszul duality, $R(u) \mapsto R(u)^{-1}$
($E_1$~inversion,
Remark~\ref{rem:yangian-e1-inversion}).

\smallskip
Since every standard object of
$\mathcal{O}^{\mathrm{sh}}_{\leq 0}$ lies in
$\widehat{D}(\operatorname{thick}\langle
\mathcal{G}_{\mathrm{Bax}}\rangle)$ by Steps~2--3, and every
object of~$\mathcal{O}^{\mathrm{sh}}_{\leq 0}$ has a finite
filtration by standards (BGG), the image of
$D^b(\mathcal{O}^{\mathrm{sh}}_{\leq 0})$ in the separated
completion is contained in the localizing subcategory generated by
$\mathcal{G}_{\mathrm{Bax}}$.

\smallskip\noindent\textit{Conditional completion extension.}
Assume now that $\mathcal{D}^{\mathrm{comp}}$ is a presentable stable
completed enhancement whose compact objects are the thick
idempotent-complete closure of the image of
$D^b(\mathcal{O}^{\mathrm{sh}}_{\leq 0})$. Then the compact part of
$\mathcal{D}^{\mathrm{comp}}$ is already generated by
$\mathcal{G}_{\mathrm{Bax}}$, and the whole category is the localizing
subcategory generated by its compact objects. This is the classical
compact-generation mechanism of Neeman and Bondal--Van den Bergh; in
the stable $\infty$-categorical form used here it is the ind-completion
formalism of Lurie~\cite[Prop.~5.3.5.10, Prop.~5.5.7.8]{HTT}. Hence
$\mathcal{D}^{\mathrm{comp}} = \operatorname{Loc}\langle
\mathcal{G}_{\mathrm{Bax}}\rangle$.

This proves the unconditional Baxter-locus generation statement on the
shifted envelope. The remaining compact/completion packet is the
comparison between the compact shifted-prefundamental core and the
desired completed Yangian category. It is not used in the proved
shifted-envelope generation statement.
\end{proof}

\begin{conjecture}[Lift-and-lower \textup{(}Lemma~L\textup{)};
\ClaimStatusConjectured]
\label{conj:rank-independence-step2}
\index{MC3!Lemma L}
\index{lift-and-lower lemma}
Let~$\mathfrak{g}$ be any finite-dimensional simple Lie algebra and
let~$\mathfrak{g}^A \subseteq \mathfrak{g}$ denote any type-$A$
Levi subalgebra of maximal rank. Assume the categorical
Clebsch--Gordan closure for prefundamental modules of~$\mathfrak{g}$
\textup{(}Theorem~\textup{\ref{thm:categorical-cg-all-types}}\textup{)}.
Let $M(\lambda)$ be a standard in the shifted category
$\mathcal{O}^{\mathrm{sh}}_{\leq 0}(Y(\mathfrak{g}))$. Then
$M(\lambda)$ lies in the thick idempotent-complete closure of
$\mathcal{G}_{\mathrm{shift}} =
\{V_{\omega_i}(a)\} \cup \{L^-_{i,a}\}$ in
$\widehat{D}(\mathcal{O}^{\mathrm{sh}}_{\leq 0}(Y(\mathfrak{g})))$,
provided the following \emph{lift-and-lower} hypothesis holds:
the type-$A$ generation statement
\textup{(}Theorem~\textup{\ref{thm:shifted-prefundamental-generation}}\textup{)}
applied to~$\mathfrak{g}^A$ lifts to~$\mathfrak{g}$ under the
Hernandez block separation of~\cite{HJZ25}, and the resulting
filtration lowers back to the thick closure in the shifted
category of~$\mathfrak{g}$ via the prefundamental
Clebsch--Gordan decomposition
\textup{(}Proposition~\textup{\ref{prop:prefundamental-clebsch-gordan}}\textup{)}.
\end{conjecture}

\begin{remark}[Scope of Lemma~L]
\label{rem:lemma-L-status}
The lift-and-lower hypothesis is proved for type~$A$
\textup{(}tautologically, since $\mathfrak{g}^A = \mathfrak{g}$\textup{)}
and is the one remaining input required to promote
Theorem~\ref{thm:shifted-prefundamental-generation} and
Proposition~\ref{conj:mc3-automatic-generalization} from type~$A$
to arbitrary simple type. It is the categorical lift of the
$K_0$-level block-separation argument used in
Theorem~\ref{thm:categorical-cg-all-types}; the missing step is
the upgrade from the $\ell$-weight multiplicity-free decomposition
at the Grothendieck-group level to a filtration-level statement
in the derived shifted category. Pending Lemma~L, the layer~3
component of the MC3 split
\textup{(}Corollary~\textup{\ref{cor:mc3-all-types}}\textup{)}
for non-type-$A$ simple Lie algebras is conditional.
\end{remark}

\begin{remark}[Type-A resolution of shifted-prefundamental generation]
\label{rem:shifted-prefundamental-generation-typeA}
In type~$A$ ($\mathfrak{sl}_N$ for all $N \geq 2$),
Theorem~\ref{thm:shifted-prefundamental-generation} proves the
shifted-envelope generation statement on the Baxter locus
$b = a - \tfrac{1}{2}$: every standard object of
$\mathcal{O}^{\mathrm{sh}}_{\leq 0}$ lies in the thick
idempotent-complete closure of
$\{V_{\omega_i}(a)\} \cup \{L^-_{i,a-\frac{1}{2}}\}$ inside the
separated completion, and hence the image of
$D^b(\mathcal{O}^{\mathrm{sh}}_{\leq 0})$ is generated there by
this set. Since $a \mapsto a - \tfrac{1}{2}$ is bijective, this is
equivalent to generation by the full prefundamental family. If a chosen
completed enhancement is compactly generated with compact part given by
the shifted envelope, the same Baxter seed generates that enhancement
as a localizing subcategory. The compact-object identification for the
intended Yangian completion is the separate remaining packet
\textup{(}Conjecture~\textup{\ref{conj:dk-compacts-completion}}\textup{)}.
The proof uses the Baxter singular vector
(Proposition~\ref{prop:baxter-yangian-equivariance},
rank-independent) and the prefundamental Clebsch--Gordan closure
(Proposition~\ref{prop:prefundamental-clebsch-gordan}).
The prefundamental Clebsch--Gordan closure
(Proposition~\ref{prop:prefundamental-clebsch-gordan}) holds for
all simple~$\mathfrak{g}$ at the character level. Extension to
arbitrary type thus reduces to the categorical lift: upgrading the
$K_0$-level identity to a module decomposition in~$\mathcal{O}^{\mathrm{sh}}$.
\end{remark}

\begin{remark}[Removing the Baxter constraint from MC3]
\label{rem:mc3-baxter-constraint-removal}
\index{MC3!Baxter constraint!removal criterion}
\index{Baxter constraint!removal criterion}
\smallskip\noindent\emph{Proved core.}
Theorem~\ref{thm:shifted-prefundamental-generation} proves the
type-$A$ statement on the Baxter-locus seed
$\{V_{\omega_i}(a)\}\cup\{L^-_{i,a-\frac{1}{2}}\}$.
The issue is therefore not set-theoretic coverage of the
prefundamental family: because $a \mapsto a-\tfrac{1}{2}$ is a
bijection, every negative prefundamental $L^-_{i,b}$ already appears
on that seed. The issue is the proof mechanism. Step~2 uses the
rank-$1$ singular vector on the hyperplane
\eqref{eq:baxter-spectral-constraint}; away from that hyperplane, the
present manuscript has no module-level replacement for the Baxter
companion.

\smallskip\noindent\emph{Alternative generators and the deformation
route.}
The natural replacements are not new finite-dimensional modules, since
Corollary~\ref{cor:dk2-thick-generation-all-types} already shows that
fundamental evaluations generate only the evaluation core. The real
candidates are the asymptotic modules of the Kirillov--Reshetikhin
tower in the sense of \cite{HJZ25,Zhang24} and the Baxter--Rees family
$\mathcal{E}_i$ constructed later in this chapter.
Theorem~\ref{thm:yangian-generic-boundary-fibers} identifies its
generic fibers with Cartan-renormalized asymptotic modules and its
boundary fiber with the negative prefundamental module. A genuine
removal of the Baxter constraint proceeds by
deformation: first show that the generic asymptotic fibers already lie
in the localizing subcategory generated by the evaluation core inside
the chosen completion, then prove that this localizing subcategory is
stable under specialization along the flat family
$\mathcal{E}_i \to \mathbb{A}^1_q$. The boundary module $L^-_{i,b}$
then follows by degeneration, so the explicit Baxter singular
vector becomes only one rank-$1$ chart of the argument rather than the
engine of the proof.

\smallskip\noindent\emph{Why this is the correct deformation of the
classical theorem.}
In the non-shifted finite-dimensional theory, ordered tensor products
of fundamental evaluation modules generate without any spectral
constraint. The shifted problem is not addressed by perturbing the
singular vector off the hyperplane
\eqref{eq:baxter-spectral-constraint}. The sharper deformation is to
move from the finite-dimensional evaluation core to the asymptotic
chart of the Kirillov--Reshetikhin tower and then specialize to the
prefundamental boundary. The Baxter--Rees family supplies exactly that
interpolation.

\smallskip\noindent\emph{Completion boundary.}
The completion step remains separate. Theorem~\ref{thm:shifted-prefundamental-generation}
and Corollary~\ref{cor:sectorwise-localizing-generation} show that,
once a presentable stable completion is compactly generated by the
shifted envelope, the localizing extension is formal; see also
Remark~\ref{rem:francis-gaitsgory-mc3}. The shifted-Yangian inputs
imported from \cite{HZ24,GTL17,HJZ25} control tensor products, blocks,
and convention transfer, but they do not identify the intended
completed shifted Yangian category as such a compactly generated
localizing completion. That missing Yangian-side identification is
exactly Conjecture~\ref{conj:dk-compacts-completion}.

\smallskip\noindent\emph{Type \texorpdfstring{$B/C/D$}{B/C/D}.}
The honest all-type analogue of the type-$A$ Baxter hyperplane is not
a blind reuse of $b = a-\tfrac{1}{2}$. In the shifted-Yangian Baxter
package of \cite{HZ24}, the spectral data is organized by the
generalized simple roots
\[
 A_{i,a}
 \;=\;
 \prod_{j:\,c_{ij}\neq 0}
 \frac{\Psi_{j,a-d_{ij}}}{\Psi_{j,a+d_{ij}}},
 \qquad
 d_{ij} := \frac{(\alpha_i,\alpha_j)}{2},
\]
so the relevant shift is root-datum dependent. After converting
conventions, simply-laced type~$D$ has the same uniform local shift
pattern as type~$A$; the remaining obstruction there is still the
lift-and-lower categorical step of
Conjecture~\ref{conj:rank-independence-step2}. Types~$B$ and~$C$ are
different: the short/long-root asymmetry forces rootwise spectral
renormalization, producing mixed shifts rather than a single
universal half-step. The manuscript does not yet prove a
negative-prefundamental companion theorem with those renormalized
shifts. This is why the asymptotic/Baxter--Rees deformation route is
the sharper target: it attacks the spectral-normalization question and
the completion question at the same time.
\end{remark}

\begin{remark}[Genericity of prefundamental irreducibility at higher rank]
\label{rem:prefundamental-genericity-higher-rank}
\index{prefundamental module!genericity condition}
For $Y(\mathfrak{sl}_2)$, the negative prefundamental~$L^-(a)$
is irreducible for \emph{all} spectral parameters $a \in \mathbb{C}$
(Hern'andez--Jimbo~\cite{HJZ25}, Theorem~3.8).
For $Y(\mathfrak{sl}_N)$ with $N \ge 3$, irreducibility holds
for \emph{generic}~$a$, outside a countable, discrete subset
of~$\mathbb{C}$ where the Drinfeld polynomial of~$L^-_{i,a}$
satisfies special divisibility relations. The simple-root
$\lambda = 0$ Baxter calculation occurs on the hyperplane
$b = a - 1/2$, and the generic exact triangles used in the
type-$A$ generation argument are supplied on a cofinal family of
spectral parameters by
Remark~\ref{rem:baxter-exact-triangles-typeA}. The
Clebsch--Gordan decomposition holds for generic~$a, b$, regardless
of the genericity condition. Irreducibility enters
only in the identification of the graded pieces with simple
prefundamentals; at non-generic parameters, the graded pieces
may be indecomposable but reducible, and the thick-generation
conclusion still holds.
\end{remark}

\begin{computation}[Thick generation obstruction analysis for \texorpdfstring{$Y(\mathfrak{sl}_2)$}{Y(sl_2)}]
\label{comp:thick-generation-sl2}
\index{thick generation!obstruction analysis}
For $L^-_a$ the negative prefundamental module of~$Y(\mathfrak{sl}_2)$
with character $\operatorname{ch}(L^-) = \prod_{n \geq 1}(1-q^{2n})^{-1}$
(partition generating function), the completed character factorizes
formally as
\begin{equation}\label{eq:fock-verma-factorization}
 \operatorname{ch}(L^-)
 \;=\;
 \operatorname{ch}(M(0)) \cdot \prod_{n \geq 2}(1-q^{2n})^{-1},
\end{equation}
where the factor \(F_{\ge2}:=\prod_{n\ge2}(1-q^{2n})^{-1}\) has
coefficient \(p(k)-p(k-1)\) in weight~\(-2k\). This is an identity in
the completed character ring. It does not by itself exhibit \(M(0)\)
as a submodule, quotient, or direct summand of \(L^-\).

The \emph{obstruction sequence} $\delta(k) := p(k) - 1$ (the excess
of $L^-$ over $M(0)$ at weight~$-2k$) governs the resolution
complexity:
\begin{center}
\renewcommand{\arraystretch}{1.1}
\begin{tabular}{c|ccccccccccc}
$k$ & 0 & 1 & 2 & 3 & 4 & 5 & 6 & 7 & 8 & 9 & 10 \\
\hline
$p(k)$ & 1 & 1 & 2 & 3 & 5 & 7 & 11 & 15 & 22 & 30 & 42 \\
$\delta(k)$ & 0 & 0 & 1 & 2 & 4 & 6 & 10 & 14 & 21 & 29 & 41
\end{tabular}
\end{center}
The growth $\log\delta(k) \sim \pi\sqrt{2k/3}$ is
sub-exponential, matching the $E_1$-page growth from
Proposition~\ref{prop:lqt-e1-subexponential-growth}.

Sectorwise Hom bounds: for $V_n$ the $(n{+}1)$-dimensional
evaluation module,
\[
 \dim\operatorname{Hom}(L^-, V_n)
 \;\leq\;
 \begin{cases}
 \lfloor n/2\rfloor + 1 & \text{$n$ even,} \\
 0 & \text{$n$ odd}
 \end{cases}
\]
(weight parity obstruction for odd~$n$). The Euler characteristic
is $\chi(L^-, V_n) = \sum_{k=0}^{n/2} p(k)$ for even~$n$,
and character containment
$\operatorname{ch}(M(\lambda)) \leq
	 \operatorname{ch}(V_\lambda \otimes L^-)$
holds for \emph{all}~$\lambda$ as a coefficientwise inequality in the
completed character ring. By the prefundamental Clebsch--Gordan
character identity, \(V_\lambda\otimes L^-\) has the character of a
finite sum of shifted \(L^-\)'s, and these shifted characters dominate
the corresponding Verma characters coefficientwise.
The Hom bounds are \emph{exact} for even~$n$.

Three character-level consequences follow from the prefundamental
Clebsch--Gordan identity:
\begin{enumerate}[label=\textup{(R\arabic*)},nosep]
\item \(\operatorname{ch}(V_n \otimes L^-)
 = \sum_{j=0}^{n}\operatorname{ch}(L^-(n-2j))\):
 evaluation-stability of~$L^-$ in the completed character ring.
\item The kernel $K := \operatorname{ch}(L^-) - \operatorname{ch}(M(0))$
 decomposes with coefficients $p(k)-1$, the number of partitions of
 \(k\) with at least two parts.
\item The shifted coefficient sequence of \(K\) is recursively expressed
 in terms of the same partition numbers \(p(k)\).
\end{enumerate}
In~$K_0$: the ring generated by $\{[V_n], [L^-]\}$ is closed
under the tensor product (R1 gives the multiplication table).
The Grothendieck class $[M(\lambda)]$ lies in this ring
for all~$\lambda$ in the finite-window calculation, and the recursive
form of R1 gives the proposed all-\(\lambda\) extension.
\end{computation}

\begin{proposition}[Universal prefundamental Clebsch--Gordan; \ClaimStatusProvedHere]
\label{prop:prefundamental-clebsch-gordan}%
\index{prefundamental!Clebsch--Gordan}%
\index{Clebsch--Gordan!prefundamental}%
\index{Clebsch--Gordan!universal}%
Let\/ $\mathfrak{g}$ be any simple Lie algebra with Dynkin
node set~$I$, and let\/ $i \in I$. For any dominant
weight~$\lambda$ of~$\mathfrak{g}$, the Hern\'andez--Jimbo
negative prefundamental module $L^-_i$ satisfies
\begin{equation}\label{eq:prefundamental-cg}
 \operatorname{ch}\bigl(V_\lambda \otimes L^-_i\bigr)
 \;=\;
 \sum_{\mu \in \operatorname{wt}(V_\lambda)}
 \operatorname{mult}_\lambda(\mu)\,
 \operatorname{ch}\bigl(L^-_i(\mathrm{shift} = \mu)\bigr).
\end{equation}
In particular: \textup{(a)}~same-index closure holds at the character
level: the right-hand side is a sum of shifted $L^-_i$ characters only,
with no $L^-_j$ character for $j \neq i$;
\textup{(b)}~the total multiplicity equals $\dim V_\lambda$;
\textup{(c)}~for $\mathfrak{g} = \mathfrak{sl}_2$ and
$\lambda = n\omega$, this recovers the character identity
\(\operatorname{ch}(V_n\otimes L^-)
=\sum_{j=0}^{n}\operatorname{ch}(L^-(n{-}2j))\).
\end{proposition}

\begin{proof}
$\operatorname{ch}(V_\lambda) = \sum_\mu
\operatorname{mult}_\lambda(\mu)\, e^\mu$ is a finite Laurent
polynomial. $\operatorname{ch}(L^-_i)$ is a formal power series.
Multiplication distributes:
$\operatorname{ch}(V_\lambda) \cdot \operatorname{ch}(L^-_i)
= \sum_\mu \operatorname{mult}_\lambda(\mu)\,
 e^\mu \cdot \operatorname{ch}(L^-_i)
= \sum_\mu \operatorname{mult}_\lambda(\mu)\,
 \operatorname{ch}(L^-_i(\mathrm{shift} = \mu))$.
\end{proof}

\begin{corollary}[Universal character containment; \ClaimStatusProvedHere]
\label{cor:universal-character-containment}%
\index{character containment!universal}%
For any simple~$\mathfrak{g}$, node~$i \in I$, and dominant
weight~$\lambda$,
$\operatorname{ch}(M(\lambda))
\leq
\operatorname{ch}(V_\lambda \otimes L^-_i)$
\textup{(coefficient-by-coefficient)}.
The identity summand ($\mu = \lambda$) contributes with
multiplicity~$1$.
\end{corollary}

\begin{corollary}[$K_0$ generation for all simple types; \ClaimStatusProvedHere]
\label{cor:k0-generation-OY}%
\index{K0 generation@$K_0$ generation!category O@category $\mathcal{O}$}%
\index{thick generation!K0 level@$K_0$ level}%
\index{thick generation!all simple types}%
For any simple~$\mathfrak{g}$, every Verma class
$[M(\lambda)]$ in
$K_0(\mathcal{O}_{Y(\mathfrak{g})})$ lies in the
$\mathbb{Z}$-span of\/
$\{[V_\mu] \cdot [L^-_i] \mid \mu\text{ dominant},\; i \in I\}$.
\end{corollary}

\begin{proof}
By Proposition~\ref{prop:prefundamental-clebsch-gordan},
$[V_\lambda] \cdot [L^-_i] = \sum_\mu
\operatorname{mult}_\lambda(\mu)\, [L^-_i(\mathrm{shifted})]$.
The Fock--Verma factorization~\eqref{eq:fock-verma-factorization}
and Baxter TQ relations generate all $[M(\lambda)]$ by induction
on height.
Verified: $\mathfrak{sl}_2$ through $\mathfrak{sl}_6$.
\end{proof}

\begin{remark}[Evaluation-stability and the recursive kernel]%
\label{rem:eval-stability-recursive-kernel}%
\index{evaluation-stability}%
\index{kernel!recursive decomposition}%
\emph{(a) Evaluation-stability.}
$[V_1] \cdot [L^-] = [L^-(+1)] + [L^-(-1)]$ implies
$L^-(a) \in \operatorname{thick}\langle L^-(b), V_1\rangle$
for any $a, b$ at the $K_0$ level. The prefundamental family
forms a single orbit under evaluation-module tensoring.

\emph{(b) Recursive kernel.}
$K_1 := \operatorname{ch}(L^-) - \operatorname{ch}(M(0))$ has
coefficient $p(m){-}1$ at weight~$-2m$. Subtracting
$\operatorname{ch}(L^-(\mathrm{hw} = -4))$ yields $K_2$ with
coefficient $p(m){-}1{-}p(m{-}2)$ for $m \geq 3$: a self-similar
obstruction tower with partition-type growth at every level.

\emph{(c) Categorical upgrade.}
Lifting~\eqref{eq:prefundamental-cg} to an isomorphism
of $Y(\mathfrak{sl}_N)$-modules is proved in type~$A$ for generic
spectral parameters
(Proposition~\ref{prop:categorical-cg-typeA} below): the block
decomposition of~$\mathcal{O}^{\mathrm{sh}}$
forces the character-level filtration to split.
\end{remark}

\begin{proposition}[Categorical prefundamental CG decomposition,
 type~$A$; \ClaimStatusProvedHere]
\label{prop:categorical-cg-typeA}%
\index{prefundamental!categorical Clebsch--Gordan!type A|textbf}%
\index{Clebsch--Gordan!categorical lift}%
For\/ $\mathfrak{g} = \mathfrak{sl}_N$ with $N \geq 2$ and each
fundamental weight~$\omega_i$ ($1 \leq i \leq N{-}1$), the
character-level identity of
Proposition~\textup{\ref{prop:prefundamental-clebsch-gordan}}
lifts to a module decomposition: for generic spectral
parameters~$a,b$ \textup{(}$a - b$ outside a discrete resonance
set determined by the normalized $R$-matrix
denominator\textup{)},
\begin{equation}\label{eq:categorical-cg-typeA}
 V_{\omega_i}(a) \otimes L^-_i(b)
 \;\cong\;
 \bigoplus_{\mu \in \operatorname{wt}(V_{\omega_i})}
 L^-_i(\mathrm{shift} = \mu)
\end{equation}
as $Y(\mathfrak{sl}_N)$-modules in\/
$\mathcal{O}^{\mathrm{sh}}_{\leq 0}$, using the same shift
convention as
Proposition~\textup{\ref{prop:prefundamental-clebsch-gordan}}:
the $\ell$-highest weight of each summand encodes the full
weight~$\mu$, including its pairings with all simple coroots.
The generic-parameter condition suffices for all generation
arguments in the MC3 comparison, since compact generation requires
exact triangles only at a cofinal family of spectral parameters.
\end{proposition}

\begin{proof}
\emph{Step~1: Submodule construction.}
For each weight~$\mu \in \operatorname{wt}(V_{\omega_i})$,
the singular vector construction of
Proposition~\ref{prop:baxter-yangian-equivariance} (applied
to the $\mu$-weight line of~$V_{\omega_i}$) produces a
$Y(\mathfrak{sl}_N)$-equivariant embedding
\[
 \iota_\mu \colon
 L^-_i(\mathrm{shift} = \mu)
 \;\hookrightarrow\;
 V_{\omega_i}(a) \otimes L^-_i(b).
\]
The embedding sends the highest $\ell$-weight vector
of~$L^-_i(\mathrm{shift} = \mu)$ to the Yangian singular
vector in the $\mu$-weight subspace
of~$V_{\omega_i} \otimes L^-_i$.
Since $V_{\omega_i}$ is minuscule for $\mathfrak{sl}_N$
(all weights have multiplicity~$1$), these embeddings give a
filtration
$0 = F_0 \subset F_1 \subset \cdots \subset F_d
= V_{\omega_i}(a) \otimes L^-_i(b)$
(where $d = \dim V_{\omega_i}$) with successive quotients
$F_j / F_{j-1} \cong L^-_i(\mathrm{shift} = \mu_j)$ for
some ordering of the weights.

\emph{Step~2: Block separation.}
Each summand $L^-_i(\mathrm{shift} = \mu)$ carries a distinct
$\ell$-highest weight $\Psi^-_{i,\mu}$ encoding the full
weight vector~$\mu$. For $\mu \neq \nu$ in
$\operatorname{wt}(V_{\omega_i})$, the $\ell$-weight ratio
$\Psi^-_{i,\mu} / \Psi^-_{i,\nu}$ is a monomial~$x^{\mu - \nu}$
with $\mu - \nu \neq 0$ in the weight lattice.
By the block criterion for category~$\mathcal{O}^{\mathrm{sh}}$
\cite{Hernandez05, HJZ25}: two simple highest-$\ell$-weight
modules lie in the same block only if their $\ell$-weight ratio
is a product of shifted simple root contributions. For
generic spectral parameters (avoiding a discrete resonance set
where $\mu - \nu$ aligns with the normalized $R$-matrix
denominator), the monomial~$x^{\mu - \nu}$ admits no such
factorization, so $L^-_i(\mathrm{shift} = \mu)$ and
$L^-_i(\mathrm{shift} = \nu)$ lie in \emph{distinct} blocks.
Hence
\[
\operatorname{Ext}^1_{\mathcal{O}^{\mathrm{sh}}}
\!\bigl(L^-_i(\mathrm{shift} = \mu),\;
 L^-_i(\mathrm{shift} = \nu)\bigr) = 0
\qquad (\mu \neq \nu,\;\text{generic}).
\]

\emph{Step~3: Splitting.}
The filtration from Step~1 has successive quotients in
pairwise distinct blocks (Step~2). Each extension
$0 \to F_{j-1} \to F_j \to
L^-_i(\mathrm{shift} = \mu_j) \to 0$
therefore splits: the right-hand term lies in a block disjoint
from all composition factors of~$F_{j-1}$, so
$\operatorname{Ext}^1(L^-_i(\mathrm{shift} = \mu_j),
F_{j-1}) = 0$.
Induction on~$j$ gives
$V_{\omega_i}(a) \otimes L^-_i(b)
\cong \bigoplus_\mu L^-_i(\mathrm{shift} = \mu)$.
The character identity
(Proposition~\ref{prop:prefundamental-clebsch-gordan}) confirms
that the direct sum accounts for the full character with no
remainder.
\end{proof}

\begin{theorem}[Pro-Weyl recovery of ordinary standards in type~$A$; \ClaimStatusConditional]
\label{thm:pro-weyl-recovery}
\index{Weyl module!pro-Weyl tower}
\index{standard module!pro-Weyl recovery}
For $\fg = \mathfrak{sl}_N$ with $N \ge 2$ and a rational highest
$\ell$-weight~$\Psi$, let
$\Psi_{\leq m}$ denote its truncation to finite Drinfeld divisor
data and $W_m := W(\Psi_{\leq m})$ the local Weyl module.
Then the ordinary standard module is recovered as a derived
inverse limit
\[
M(\Psi) \;\simeq\; R\varprojlim_m\, W_m
\]
inside a separated completion of~$\mathcal{O}$.
Each $W_m$ is universal finite-dimensional highest-weight, hence lies
in $\operatorname{thick}\langle \mathrm{Ev} \rangle$ by the proved
evaluation core (Theorem~\ref{thm:eval-core-identification}).
\end{theorem}

\begin{proof}
The truncation tower $\{W_m\}_{m \geq 0}$ is compatible with the
Drinfeld polynomial data of~$\Psi$, and the transition maps
$W_{m+1} \twoheadrightarrow W_m$ are surjective at each weight. The
inverse system is therefore Mittag--Leffler, so
$R^1\varprojlim W_m = 0$ and the derived inverse limit agrees with the
ordinary inverse limit. By the universal property of local Weyl
modules, a vector in the standard module $M(\Psi)$ depends on only
finitely many coefficients of the rational highest $\ell$-weight, so
the compatible system of truncations identifies
$M(\Psi) \cong \varprojlim_m W_m$. This yields
$M(\Psi) \simeq R\varprojlim_m W_m$ inside the separated completion.
Each $W_m$ is finite-dimensional, hence belongs to
$\operatorname{thick}\langle \mathrm{Ev} \rangle$ by
Theorem~\ref{thm:eval-core-identification}.
\end{proof}

\begin{remark}[Type-A resolution of pro-Weyl recovery]
\label{rem:pro-weyl-recovery-typeA}
In type~$A$ ($\mathfrak{sl}_N$ for all $N \geq 2$),
Theorem~\ref{thm:pro-weyl-recovery} resolves the pro-Weyl packet.
The Weyl module truncation tower
$\{W_m = W(\Psi_{\leq m})\}_{m \geq 0}$ satisfies the
Mittag--Leffler condition: the transition maps
$W_{m+1} \twoheadrightarrow W_m$ are surjective at each weight,
whence $R^1\varprojlim = 0$ and the ordinary standard is recovered as
$M(\Psi) \simeq R\varprojlim_m W_m$ inside the separated completion
$\widehat{\mathcal{O}}$.
The argument is rank-independent (it depends only on the weight
filtration, not on the Dynkin type), and the finite-dimensionality of
each~$W_m$ is guaranteed by the proved evaluation core
(Theorem~\ref{thm:eval-core-identification}).
Extension to arbitrary type remains part of the unresolved
post-core extension/completion packet: one still needs the
shifted-prefundamental generation and completion input beyond the
type-$A$ surface.
\end{remark}

\begin{conjecture}[DK on compacts and completion extension; \ClaimStatusConjectured]
\label{conj:dk-compacts-completion}
\index{Drinfeld--Kohno!derived!compact extension}
\index{Francis--Gaitsgory!pro-nilpotent completion}
The bar-cobar Koszul duality equivalence
$\Phi$~\eqref{eq:dk-eval-core} extends from the finite-dimensional
evaluation core to the full completed shifted-prefundamental core
$\mathcal{G}_{\mathrm{shift}}$, and then by pro-nilpotent completion
in the sense of Francis--Gaitsgory to the completed Yangian category.
\end{conjecture}

\begin{remark}[Two-clause decomposition of Conjecture~\ref{conj:dk-compacts-completion}]
\label{rem:dk-compacts-completion-two-clause}
\index{MC3!compact-completion!two-clause decomposition}%
The conjecture statement is a conjunction of two structurally distinct
items; a proof of either clause alone does not close the MC3 gap.
\begin{enumerate}[label=\textup{(C\arabic*)},nosep]
\item \emph{Compact-core extension of~$\Phi$.} The bar-cobar functor
 of Theorem~\ref{thm:eval-core-identification} extends from the
 finite-dimensional evaluation core to the compact
 shifted-prefundamental core $\mathcal{G}_{\mathrm{shift}}$ inside
 the separated completion, with the three weight-bounded / colimit /
 compactness hypotheses of
 Proposition~\ref{prop:completion-extension-weight-bounded}
 verified in the chosen completed ambient.
\item \emph{Francis--Gaitsgory pro-nilpotent comparison.} The
 resulting compactly-generated enhancement agrees with the
 Francis--Gaitsgory pro-nilpotent completion of
 $D^b(\mathcal{O}^{\mathrm{sh}}_{\leq 0}(Y(\mathfrak{g})))$ in the
 sense of~\cite{GR17} Vol~II Ch.~4, so that the extended equivalence
 is a statement about the same target category that the MC3 comparison
names in the preface.
\end{enumerate}
A proof of~(C1) on its own yields an equivalence on some
compactly-generated enhancement, leaving the identification with the
Francis--Gaitsgory completion open; a proof of~(C2) on its own
provides the completion comparison but not the functor extension. The
two clauses are mechanically independent and require distinct inputs:
(C1) is weight-filtration and $t$-structure compatibility inside the
chosen completion; (C2) is an identification of two presentations of
the completed category.
\end{remark}

\begin{remark}[Type-A DK on compact families]
\label{rem:dk-compacts-completion-typeA}
In type~$A$ ($\mathfrak{sl}_N$ for all $N \geq 2$),
packages~(i)--(iii) of the four-package MC3 list are resolved, but
Conjecture~\ref{conj:dk-compacts-completion} remains the final
unclosed packet. The proved type-$A$ input is:
\begin{enumerate}[label=\textup{(\roman*)},nosep]
\item the evaluation-core DK equivalence
 \textup{(}Theorem~\ref{thm:eval-core-identification}\textup{)};
\item shifted-prefundamental generation on the shifted envelope
 inside the separated completion
 \textup{(}Theorem~\ref{thm:shifted-prefundamental-generation}\textup{)};
\item pro-Weyl recovery of standards
 \textup{(}Theorem~\ref{thm:pro-weyl-recovery}\textup{)}.
\end{enumerate}
What is not yet proved here is the additional compact-core extension
of~$\Phi$ and the comparison between the desired completed category
and the relevant Francis--Gaitsgory/pro-nilpotent completion. Once a
chosen completion is shown to be compactly generated by the shifted
envelope, the localizing extension is formal by
Theorem~\ref{thm:shifted-prefundamental-generation}; what remains is
verifying those hypotheses and extending~$\Phi$ on the compact core.
\end{remark}

\begin{remark}[Finite-window MC3 checks]
\label{rem:mc3-computational-evidence}
\index{MC3!finite-window checks}
Theorems~\ref{thm:baxter-exact-triangles},
\ref{thm:shifted-prefundamental-generation},
\ref{thm:pro-weyl-recovery}, and
Conjecture~\ref{conj:dk-compacts-completion} meet the following
finite-window consistency checks:
\begin{enumerate}[label=\textup{(\roman*)},nosep]
\item \emph{Baxter}: $K_0$-level TQ identities for
 $\mathfrak{sl}_2$ and~$\mathfrak{sl}_3$, together with derived-level
 lifts to distinguished triangles, mapping cones, and the octahedral
 axiom.
\item \emph{Prefundamental}: Hern'andez--Jimbo negative
 prefundamental modules~$L^-_a$ for $Y(\mathfrak{sl}_2)$
 at the character/q-character level. The weight multiplicities
 $\dim(L^-_a)_{-2k} = p(k)$ connect MC3 to partition theory.
\item \emph{Pro-Weyl}: M-level convergence through chain-level Ext
 computations, the derived inverse limit $R\varprojlim W_m$, and
 R$^1$ vanishing via Mittag--Leffler.
\item \emph{DK compacts}: Feigin--Gainutdinov completion gap
 analysis and Francis--Gaitsgory pro-nilpotent approximation.
\item \emph{Thick generation obstructions}: the character of the
 negative prefundamental $L^-_a$ for $Y(\mathfrak{sl}_2)$
 admits the Fock--Verma factorization
 \[
 \operatorname{ch}(L^-)
 \;=\;
 \operatorname{ch}(M(0)) \cdot \operatorname{ch}(F_{\geq 2}),
 \qquad
 F_{\geq 2}
 \;:=\;
 \prod_{n \geq 2} \frac{1}{1 - q^{2n}},
 \]
 identifying $L^-$ as a ``multi-particle dressing'' of the Verma
 module~$M(0)$. At weight~$-2k$, the \emph{excess dimension}
 $p(k) - 1$ (partitions of~$k$ with at least two parts) grows
 sub-exponentially:
 $\log(p(k)-1) \sim \pi\sqrt{2k/3}$ (Hardy--Ramanujan).
 This matches the $E_1$-page growth bounds of
 Proposition~\ref{prop:lqt-e1-subexponential-growth}, confirming
 that the bar complex resolution has compatible combinatorial
 complexity. Verified:
 character containment
 $\operatorname{ch}(M(\lambda))
 \leq \operatorname{ch}(V_\lambda \otimes L^-)$
 for \emph{all}~$\lambda$
 (Corollary~\ref{cor:universal-character-containment});
 sectorwise finiteness
 $\dim\operatorname{Hom}(L^-, V_n) \leq \lfloor n/2 \rfloor + 1$
 for even~$n$, and $= 0$ for odd~$n$;
 iterated Baxter generation $V_1^{\otimes n} \otimes L^-$
 alternates between even and odd weight lattices with
 highest weight increasing by~$1$ at each step;
 Euler characteristics
 $\chi(L^-, V_n) = \sum_{k=0}^{n/2} p(k)$ for even~$n$.
\end{enumerate}
\end{remark}

\begin{remark}[Corrected MC3 open problem]
\label{rem:corrected-mc3-frontier}
\index{MC3!corrected open problem}
Theorem~\ref{thm:eval-core-identification} closes the thick
generation problem:
\begin{equation}\label{eq:eval-core-sandwich}
\mathcal{D}_\hbar^{\mathrm{eval}}
\;=\; D^b(\operatorname{Rep}_{\mathrm{fd}})
\;\subsetneq\; D^b(\mathcal{O}_{Y_\hbar}).
\end{equation}
The left equality is the theorem; the strict inclusion follows from
Lemma~\ref{lem:fd-thick-closure} whenever $\mathcal{O}_{Y_\hbar}$
contains infinite-dimensional objects (Verma or asymptotic modules).
No further enlargement of the thick closure is possible within
ordinary $D^b(\mathcal{O})$.

With the all-types categorical CG and evaluation-generated-core DK
surfaces now separated from the remaining completion packets, the open
problem splits into two layers.
Within ordinary $D^b(\mathcal{O})$, thick generation by evaluation
modules is \emph{completely resolved}: it yields exactly
$D^b(\operatorname{Rep}_{\mathrm{fd}})$. Extending DK beyond
$\operatorname{Rep}_{\mathrm{fd}}$ requires a completed or coderived
enhancement and reopens the generation problem in that enlarged
setting.

\begin{enumerate}[label=\textup{(\alph*)}]
\item If the factorization DK conjecture
 \textup{(}Conjecture~\textup{\ref{conj:full-derived-dk}}\textup{)}
 targets the evaluation-generated core, then the needed core inputs
 are already in place: the type-$A$ module-side evaluation-core DK
 equivalence is proved by~\eqref{eq:dk-eval-core}, and the
 factorization comparison on the evaluation-generated core is proved
 separately by Theorem~\ref{thm:h-level-factorization-kd} and
 Corollary~\ref{cor:dk23-all-types}. On that restricted core
 surface, the remaining ladder proceeds directly to DK-4.
\item If one wishes to include infinite-dimensional
 modules \textup{(}prefundamental representations, asymptotic
 modules in the sense of Hernandez--Zhang~\cite{HJZ25}\textup{)}, then
 a completed, coderived, or ind-completed enhancement
 $\widehat{D}(\mathcal{O})$ is required
 \textup{(}Remark~\textup{\ref{rem:fd-thick-closure-implications}(b)}\textup{)},
 and the generation mechanism shifts from thick closure to
 localizing subcategories with compact generation. The
 structural input for this extension is organized by four
 packages, three of which are resolved in type~$A$ and the
 fourth of which remains open:
 \begin{enumerate}[label=\textup{(\roman*)}]
 \item \emph{Baxter exact triangles}
 \textup{(}Theorem~\textup{\ref{thm:baxter-exact-triangles}}\textup{)}:
 the decategorified TQ relations of Zhang~\cite{Zhang24}
 lift to distinguished triangles in the anti-dominant shifted
 envelope~$\mathcal{O}^{\mathrm{sh}}_{\leq 0}$, with
 finite-dimensional companion modules as the middle term;
 \item \emph{Shifted-prefundamental generation on the shifted envelope}
 \textup{(}Theorem~\textup{\ref{thm:shifted-prefundamental-generation}}\textup{)}:
 every standard object of the anti-dominant shifted envelope
 lies in the thick idempotent-complete closure of
 $\{\text{KR evaluations}\} \cup \{L^-_{i,a-\frac{1}{2}}\}$ inside the
 separated completion, equivalently of the full prefundamental family;
 \item \emph{Pro-Weyl recovery}
 \textup{(}Theorem~\textup{\ref{thm:pro-weyl-recovery}}\textup{)}:
 ordinary standards $M(\Psi)$ are recovered as
 $R\varprojlim_m W_m$ inside a separated completion, where
 $W_m = W(\Psi_{\leq m})$ are local Weyl module truncations;
 \item \emph{DK on compacts, extended by completion}
 \textup{(}Conjecture~\textup{\ref{conj:dk-compacts-completion}}\textup{)}:
 the bar-cobar equivalence holds on the compact shifted-prefundamental
 core and extends to the completed category by
 pro-nilpotent completion in the sense of
 Francis--Gaitsgory.
 \end{enumerate}
 The natural proof order is
 $\mathfrak{sl}_2$ Baxter exact sequence
 $\to$ $\mathfrak{sl}_3$ exact triangle
 $\to$ pro-Weyl completion for one rational highest-$\ell$-weight family
 $\to$ compact generation
 $\to$ factorization/KL extension.
 The formal reductions
 \textup{(}Propositions~\textup{\ref{prop:heart-capture-criterion}}
 and~\textup{\ref{prop:standard-capture-criterion}},
 Corollary~\textup{\ref{cor:sectorwise-localizing-generation}}\textup{)}
 convert the generation problem into sectorwise control of
 standard/Weyl objects in finite truncation sectors.
\item \emph{Type-$A$ resolution.}
 In type~$A$ ($\mathfrak{sl}_N$ for all $N \geq 2$), the four
 post-core MC3 packages reduce to a single remaining
 compact-completion packet
 (Theorem~\ref{thm:mc3-type-a-resolution}).
 The proof proceeds through: (1)~the prefundamental Clebsch--Gordan
 closure
 (Proposition~\ref{prop:prefundamental-clebsch-gordan}:
 $V_n \otimes L^- \cong \bigoplus L^-(\mathrm{shifted})$) lifts the
 Baxter TQ relations to exact triangles;
 (2)~the chromatic/conformal-weight filtration reduces
 shifted-prefundamental generation to countable finite-dimensional
 strata;
 (3)~Mittag--Leffler ($R^1\varprojlim = 0$) on Weyl module
 truncation towers recovers standards as derived inverse limits;
 (4)~spectral sequence degeneration on weight strata, together
 with the Francis--Gaitsgory pro-nilpotent completion formalism,
 isolates the last remaining compact/completed comparison problem:
 one still needs the compact-core extension of the bar-cobar
 equivalence and its comparison with the desired completed
 category.
 %
 For arbitrary simple~$\mathfrak{g}$,
 Theorem~\ref{thm:categorical-cg-all-types} supplies the
 all-types categorical Clebsch--Gordan input, while the later
 post-core pro-Weyl / completion argument becomes
 type-independent only conditionally, once the shifted-
 prefundamental, pro-Weyl, and compact-completion packets of
 Proposition~\ref{conj:mc3-automatic-generalization} are
 separately supplied in that type.
\end{enumerate}
For the factorization DK problem as posed in
Conjecture~\ref{conj:full-derived-dk}, option~(a) is the natural
reading: the $\Eone$-factorization categories are generated by
evaluation objects, and the evaluation-generated core is the
target of the equivalence.
Option~(b) is the deeper project, with the Yangian analogue
of the Hern\'andez--Negu\c{t} ordinary/shifted bridge
\cite{HerNeg24} as the target theorem.
In type~$A$, option~(b) is reduced to a single remaining packet:
Theorem~\ref{thm:mc3-type-a-resolution} proves packages~(i)--(iii)
and isolates Conjecture~\ref{conj:dk-compacts-completion} as the last
step needed to reach the entire completed shifted-prefundamental
category.
\end{remark}

\begin{theorem}[Type-$A$ MC3 reduction to the compact-completion packet; \ClaimStatusConditional]
\label{thm:mc3-type-a-resolution}
\index{MC3!type-A resolution|textbf}
\index{Drinfeld--Kohno!derived!type-A completion}
\index{prefundamental!Clebsch--Gordan closure!type A}
For $\mathfrak{g} = \mathfrak{sl}_N$ with $N \geq 2$, the
post-core MC3 extension problem is reduced to a single remaining
compact-completion packet:
\begin{enumerate}[label=\textup{(\roman*)},nosep]
\item \emph{Baxter exact triangles.}
 The TQ relations lift from $K_0$ to derived exact triangles in the
 completed anti-dominant shifted envelope
 $\widehat{\mathcal{O}}^{\mathrm{sh}}_{\leq 0}$.
 The prefundamental Clebsch--Gordan closure
 $V_n \otimes L^- \cong \bigoplus_{j=0}^{n} L^-(n{-}2j)$
 \textup{(}Proposition~\textup{\ref{prop:prefundamental-clebsch-gordan}}\textup{)}
 provides the structural input.
\item \emph{Shifted-prefundamental generation.}
 Every standard object of
 $\mathcal{O}^{\mathrm{sh}}_{\leq 0}(Y(\mathfrak{g}))$ lies in the
 thick idempotent-complete closure of the Baxter-locus seed
 $\{V_{\omega_i}(a)\} \cup \{L^-_{i,a-\frac{1}{2}}\}$ inside the
 separated completion, and hence the image of
 $D^b(\mathcal{O}^{\mathrm{sh}}_{\leq 0}(Y(\mathfrak{g})))$
 is contained in the localizing subcategory they generate. Under the
 compact/localizing hypotheses of
 Theorem~\ref{thm:shifted-prefundamental-generation}, the same seed
 generates the chosen completed enhancement as well.
 \textup{(}Theorem~\textup{\ref{thm:shifted-prefundamental-generation}}\textup{)},
 proved via the chromatic/conformal-weight filtration reducing to
 countable finite-dimensional strata.
\item \emph{Pro-Weyl recovery.}
 Standards $M(\Psi)$ are recovered as $R\varprojlim_m W_m$ inside the
 separated completion, proved via Mittag--Leffler
 ($R^1\varprojlim = 0$) on Weyl module truncation towers.
\item \emph{Remaining packet.}
The unresolved step is precisely the compact-completion conjecture:
extending the
evaluation-core DK equivalence to the compact shifted-prefundamental
core and verifying the compact/localizing hypotheses needed to compare
 that compact-core equivalence with the desired
 completed/pro-nilpotent category.
\end{enumerate}
In particular, type~$A$ now leaves a single explicit gap beyond the
evaluation-generated core rather than four independent unresolved
packets.
\end{theorem}

\begin{proof}
The four packages are proved in the preceding material:
\begin{enumerate}[label=(\roman*),nosep]
\item The Baxter exact triangles follow from the categorical
 prefundamental Clebsch--Gordan decomposition
 (Proposition~\ref{prop:categorical-cg-typeA}), which upgrades the
 character-level identity
 (Proposition~\ref{prop:prefundamental-clebsch-gordan}) to a
 module-level splitting at generic spectral parameters via the
 block separation of shifted prefundamental modules. Combined
 with the rank-independent singular vector construction
 (Proposition~\ref{prop:baxter-yangian-equivariance}) and naturality
 on $\mathcal{O}_{\mathrm{poly}}$
 (Theorem~\ref{thm:baxter-exact-triangles-opoly}), this gives
 derived exact triangles in
 $\widehat{\mathcal{O}}^{\mathrm{sh}}_{\leq 0}$ at a cofinal
 family of spectral parameters.
\item Shifted-prefundamental generation on the shifted envelope is
 Theorem~\ref{thm:shifted-prefundamental-generation}, proved by the
 four-step argument: evaluation core (Step~1), Baxter-hyperplane
 control of the simple-root prefundamental companion together with
 the generic type-$A$ Baxter package (Step~2), pro-Weyl completion of
 standards (Step~3), and $E_1$-compatibility plus passage to the
 shifted envelope and then, conditionally, to any compactly generated
 completed enhancement (Step~4). The verification of those completion
 hypotheses is the separate remaining packet in~(iv).
\item Pro-Weyl recovery follows from Step~3 of the
 Theorem~\ref{thm:shifted-prefundamental-generation} proof:
 the Mittag--Leffler condition is verified (the transition maps
 $W_{m+1} \twoheadrightarrow W_m$ are surjective at each weight),
 and each $W_m$ lies in
 $\operatorname{thick}\langle\{V_{\omega_i}\}\rangle$ by the
 evaluation core (Theorem~\ref{thm:eval-core-identification}).
\item After~(i)--(iii), the only remaining item from the original MC3
 list is the compact/completed extension packet stated above.
 No further independent type-$A$ obstruction remains on this
 surface.
\end{enumerate}
\end{proof}

\begin{theorem}[Categorical prefundamental CG decomposition, all types;
 \ClaimStatusProvedHere]
\label{thm:mc3-arbitrary-type}% backward compat
\label{thm:categorical-cg-all-types}%
\index{MC3!arbitrary type CG resolution|textbf}%
\index{prefundamental!categorical Clebsch--Gordan!all types|textbf}%
\index{q-character!multiplicity-free|textbf}%
For any simple Lie algebra~$\mathfrak{g}$ and any fundamental
weight~$\omega_i$, for generic spectral parameters~$a,b$:
\begin{equation}\label{eq:prefundamental-cg-general-type}
 V_{\omega_i}(a) \otimes L^-_{i}(b)
 \;\cong\;
 \bigoplus_{m \in \operatorname{ch}_q(V_{\omega_i}(a))}
 L^-_{i}(\mathrm{shift} = m)
\end{equation}
as $Y(\mathfrak{g})$-modules in
$\mathcal{O}^{\mathrm{sh}}$. The direct sum has
$\dim(V_{\omega_i})$ summands, each appearing with
multiplicity~$1$.
\end{theorem}

\begin{proof}
The proof has four steps. Steps~2 and~3 use results
proved for quantum affine algebras $U_q(\hat{\mathfrak{g}})$,
transferred to Yangians via the Gautam--Toledano Laredo
equivalence~\cite{GTL17}; we indicate the transfer at each step.

\emph{Step~1 (Character identity).}
The character-level identity
$\operatorname{ch}(V_{\omega_i} \otimes L^-_i)
= \sum_m \operatorname{ch}(L^-_i(\mathrm{shift} = m))$
is Proposition~\ref{prop:character-cg-all-types}, proved
algebraically for all simple types by the distributive law
of a finite Laurent polynomial (Weyl character) times a formal
power series (prefundamental character).

\emph{Step~2 (Multiplicity-free $\ell$-weights).}
For the quantum affine algebra $U_q(\hat{\mathfrak{g}})$, the
$q$-character of every fundamental evaluation module
$V_{\omega_i}(a)$ has $\dim(V_{\omega_i})$ distinct monomials,
each corresponding to a one-dimensional $\ell$-weight space.
This multiplicity-free property is not contained in a single
reference. For simply-laced types, Nakajima's computation of
$(q,t)$-characters via quiver varieties~\cite{Nakajima04} shows
that the $t = 1$ specialization gives multiplicity-free
$q$-characters for all fundamental Kirillov--Reshetikhin
modules. For classical types, Chari--Moura~\cite{ChariMoura06}
give explicit closed-form $q$-character formulas confirming
multiplicity-freeness. For all types uniformly, the
Frenkel--Mukhin algorithm~\cite{FrenkelMukhin01} correctly
computes the $q$-character of every fundamental module
\textup{(}fundamental modules are \emph{special} in the sense
of~\cite{FrenkelMukhin01}, so the algorithm terminates with the
true $q$-character\textup{)}, and the output has all monomials
distinct.

For the Yangian $Y_\hbar(\mathfrak{g})$, the analogous
multiplicity-free $\ell$-weight property follows from the
Gautam--Toledano Laredo equivalence~\cite{GTL17}
between the categories of finite-dimensional representations
of $U_q(\hat{\mathfrak{g}})$ and $Y_\hbar(\mathfrak{g})$.
This equivalence is an exact faithful functor sending simples to
simples, identifying $\ell$-weight decompositions on both sides
\textup{(}the Yangian Drinfeld generators correspond to
quantum loop Cartan generators under~\cite{GTL17}\textup{)},
and hence preserves $\ell$-weight space dimensions.

\emph{Step~3 (Block separation).}
Each summand $L^-_i(\mathrm{shift} = m_j)$ carries a distinct
highest $\ell$-weight, determined by the $\ell$-weight~$m_j$ of
$V_{\omega_i}(a)$ (Step~2) tensored with the fixed highest
$\ell$-weight of~$L^-_i(b)$. The Hern\'andez block
criterion~\cite{Hernandez05, HJZ25} states that two simple
highest-$\ell$-weight modules in category~$\mathcal{O}$ lie in
the same block only if their $\ell$-weight ratio is a product of
$A_{k,c}^{-1}$-factors. This criterion is purely combinatorial:
it depends only on the Drinfeld polynomial data parametrizing
simple modules, not on the ambient algebra. Since the
Gautam--Toledano Laredo equivalence~\cite{GTL17} identifies the
Drinfeld polynomial data of the Yangian and quantum loop algebra,
the block criterion transfers to the Yangian
category~$\mathcal{O}^{\mathrm{sh}}$.
At generic spectral parameters
\textup{(}avoiding a finite set of values of~$a - b$ where
$A$-factorizations occur\textup{)}, the $\dim(V_{\omega_i})$
highest $\ell$-weights are pairwise non-$A$-linked, hence lie in
pairwise distinct blocks.

\emph{Step~4 (Character accounting forces simplicity).}
The tensor product $V_{\omega_i}(a) \otimes L^-_i(b)$ lies in
$\mathcal{O}^{\mathrm{sh}}$
\textup{(}$V$ is finite-dimensional, $L^-_i$ is
in~$\mathcal{O}^{\mathrm{sh}}$ by~\cite{HJZ25}, and
tensoring with a finite-dimensional module preserves the
category\textup{)}. The block decomposition
in~$\mathcal{O}^{\mathrm{sh}}$ gives a direct sum
$V \otimes L^- = \bigoplus_B (V \otimes L^-)_B$.
By Step~3, the $\dim(V_{\omega_i})$ proposed summands lie in
pairwise distinct blocks $B(m_1), \dotsc, B(m_d)$.
The character identity (Step~1) gives
$\sum_B \operatorname{ch}((V \otimes L^-)_B)
= \sum_j \operatorname{ch}(L^-_i(\mathrm{shift} = m_j))$,
where each term on the right contributes to exactly one block
(Step~3). Matching characters block by block:
$\operatorname{ch}((V \otimes L^-)_{B(m_j)})
= \operatorname{ch}(L^-_i(\mathrm{shift} = m_j))$
for each~$j$, with no room for extra contributions
\textup{(}since the total is accounted for\textup{)}.
A module in~$\mathcal{O}^{\mathrm{sh}}$ whose character
equals that of a simple module~$S$ must be isomorphic to~$S$:
all its Jordan--H\"older factors lie in the same block as~$S$
and have characters dominated by~$\operatorname{ch}(S)$, so the
unique factorization is $[M:S] = 1$. Therefore
$(V \otimes L^-)_{B(m_j)} \cong L^-_i(\mathrm{shift} = m_j)$,
and the block decomposition yields the stated direct sum.
\end{proof}

\begin{remark}[External inputs]
\label{rem:categorical-cg-external-inputs}
Theorem~\ref{thm:categorical-cg-all-types} rests on three
external pillars, all published in refereed journals:
\textup{(a)}~the multiplicity-free $q$-character property
for fundamental representations
\cite{Nakajima04, FrenkelMukhin01, ChariMoura06},
\textup{(b)}~the Hern\'andez block
criterion for category~$\mathcal{O}$
\cite{Hernandez05, HJZ25}, and \textup{(c)}~the
Gautam--Toledano Laredo equivalence~\cite{GTL17}
identifying Drinfeld polynomial data between quantum
loop algebras and Yangians.
The new content is the observation
that these three inputs, combined with the character identity
(Proposition~\ref{prop:character-cg-all-types}), yield the
categorical CG decomposition for \emph{all} types, replacing
the minuscule hypothesis of
Proposition~\ref{prop:categorical-cg-typeA}.
The multiplicity-free property is not contained in a single
reference: for simply-laced types it follows from
$(q,t)$-character theory~\cite{Nakajima04}; for classical types
from explicit formulas~\cite{ChariMoura06}; for all types the
Frenkel--Mukhin algorithm~\cite{FrenkelMukhin01} computes the
$q$-character and confirms multiplicity-freeness.
The block criterion transfers from quantum affine to Yangian
because it is purely combinatorial, depending only on Drinfeld
polynomial data identified by~\cite{GTL17}.
\end{remark}

\begin{corollary}[Three-layer MC3 decomposition after categorical CG closure;
 \ClaimStatusConditional]
\label{cor:mc3-all-types}%
\index{MC3!all simple types|textbf}%
\index{MC3!three-layer split}
MC3 decomposes into three layers.
\begin{itemize}
\item \emph{Layer~1 \textup{(}MC3a, all simple types,
 unconditional\textup{)}}: the categorical Clebsch--Gordan closure
 for prefundamental modules of every simple Lie algebra
 \textup{(}Theorem~\textup{\ref{thm:categorical-cg-all-types}}\textup{)}.
\item \emph{Layer~2 \textup{(}MC3b, all simple types,
 unconditional\textup{)}}: the DK-2/3 bar-cobar equivalence on the
 evaluation-generated core
 \textup{(}Corollary~\textup{\ref{cor:dk23-all-types}}\textup{)}.
\item \emph{Layer~3 \textup{(}MC3c, shifted-envelope generation
 inside~$\widehat{D}(\mathcal{O}^{\mathrm{sh}}_{\leq 0})$\textup{)}}:
 unconditional in type~$A$
 \textup{(}Theorem~\textup{\ref{thm:shifted-prefundamental-generation}}\textup{)},
 and conditional on Conjecture~\textup{\ref{conj:rank-independence-step2}}
 \textup{(}Lemma~L, lift-and-lower\textup{)} for all other simple
 types.
\end{itemize}
Consequently, the remaining MC3 work beyond the evaluation-generated
core is concentrated in the post-CG extension/completion packets:
shifted-prefundamental generation \textup{(}Layer~3 conditional on
Lemma~L outside type~$A$\textup{)}, pro-Weyl recovery, and the
compact / completed DK extension.
\end{corollary}

\begin{proof}
Theorem~\ref{thm:categorical-cg-all-types} gives Layer~1 for all
simple types and Corollary~\ref{cor:dk23-all-types} gives Layer~2.
Layer~3 is Theorem~\ref{thm:shifted-prefundamental-generation} in
type~$A$: it identifies the shifted-envelope generation statement
inside the separated completion. For non-type-$A$ simple types, the
same four-step argument applies once the lift-and-lower input
Conjecture~\ref{conj:rank-independence-step2} is supplied. Beyond
these three layers, the compact-to-completed DK extension is
independent of Lemma~L and tracked separately by
Conjecture~\ref{conj:dk-compacts-completion}.
\end{proof}

\begin{remark}[Francis--Gaitsgory completion mechanism]\label{rem:francis-gaitsgory-mc3}
Once one has an equivalence on the relevant compact core and a
compatible Francis--Gaitsgory-style completion formalism, the passage
to the corresponding cocompletion is categorical and does not depend
on Lie type. What this mechanism does \emph{not} prove by itself is
that the required compact-core equivalence or the needed completion
comparison holds uniformly in every type.
\end{remark}

\begin{proposition}[MC3 completion extension via weight-bounded
 t-structure; \ClaimStatusConditional]
\label{prop:completion-extension-weight-bounded}
\index{MC3!completion extension}
\index{compact generation!weight-bounded t-structure}
\index{Yangian!completion!compact objects}
Let~$\mathfrak{g} = \mathfrak{sl}_N$ with $N \geq 2$.
Let $\mathcal{A} := \mathcal{O}^{\mathrm{sh}}_{\leq 0}
(Y(\mathfrak{g}))$ and let
$\iota\colon D^b(\mathcal{A}) \hookrightarrow
\mathcal{D}^{\mathrm{comp}}$
be a fully faithful exact functor into a presentable stable
$\infty$-category~$\mathcal{D}^{\mathrm{comp}}$.
Suppose:
\begin{enumerate}[label=\textup{(\roman*)},nosep]
\item\label{it:weight-compatible}
 \emph{Weight-bounded t-structure compatibility.}
 The weight decomposition
 $\mathcal{A} = \bigoplus_{\mu} \mathcal{A}_\mu$
 \textup{(}indexed by the weight lattice of~$\mathfrak{g}$,
 with each $\mathcal{A}_\mu$ finite-dimensional\textup{)}
 induces a bounded t-structure on~$D^b(\mathcal{A})$ that
 extends to a right-bounded, weight-separated t-structure
 on~$\mathcal{D}^{\mathrm{comp}}$:
 for every $M \in \mathcal{D}^{\mathrm{comp}}$ and every
 weight~$\mu$, the weight-$\mu$ graded piece
 $\operatorname{gr}_\mu(M)$ has bounded cohomological
 amplitude.
\item\label{it:colimit-generation}
 \emph{Colimit generation.}
 $\mathcal{D}^{\mathrm{comp}}$ is generated under filtered
 colimits by~$\iota(D^b(\mathcal{A}))$.
\item\label{it:compactness-preservation}
 \emph{Compactness of the image.}
 For every object $X \in D^b(\mathcal{A})$, the image
 $\iota(X)$ is compact in~$\mathcal{D}^{\mathrm{comp}}$.
\end{enumerate}
Then:
\begin{enumerate}[label=\textup{(\alph*)},nosep]
\item The compact objects of $\mathcal{D}^{\mathrm{comp}}$ are the
 thick idempotent-complete closure of~$\iota(D^b(\mathcal{A}))$.
\item $\mathcal{D}^{\mathrm{comp}} =
 \operatorname{Loc}\langle \mathcal{G}_{\mathrm{Bax}} \rangle$,
 the localizing subcategory generated by the Baxter-locus seed
 $\mathcal{G}_{\mathrm{Bax}} =
 \{V_{\omega_i}(a)\} \cup \{L^-_{i,a-\frac{1}{2}}\}$.
\item If, in addition, the bar-cobar functor~$\Phi$ of
 Theorem~\textup{\ref{thm:eval-core-identification}} is defined on
 the compact shifted-prefundamental core and restricts there to an
 equivalence compatible with the dual completion, then~$\Phi$
 extends uniquely to a localizing equivalence
 $\mathcal{D}^{\mathrm{comp}} \xrightarrow{\;\sim\;}
 (\mathcal{D}^{\mathrm{comp}})^{\mathrm{op,dual}}$.
\end{enumerate}
\end{proposition}

\begin{proof}
Three steps.

\smallskip\noindent\textit{Step~1 \textup{(}Compact identification\textup{)}.}
By~\ref{it:colimit-generation}, every object of
$\mathcal{D}^{\mathrm{comp}}$ is a filtered colimit of objects
in~$\iota(D^b(\mathcal{A}))$.
By~\ref{it:compactness-preservation}, each such image object is
compact. If $C \in \mathcal{D}^{\mathrm{comp}}$ is compact, then
$\operatorname{id}_C$ factors through a finite stage of any
presenting filtered diagram, so $C$ is a retract of an object
in~$\iota(D^b(\mathcal{A}))$. Therefore
\[
(\mathcal{D}^{\mathrm{comp}})^{\omega}
\;=\;
\overline{\iota(D^b(\mathcal{A}))}^{\,\mathrm{thick,\, idemp}}.
\]
This is the standard argument of
Neeman~\cite[Thm.~2.1]{Neeman92}
and Bondal--Van den Bergh~\cite[Cor.~3.1.2]{BondalVdB03}
in the stable $\infty$-categorical form of
Lurie~\cite[Prop.~5.3.5.10, Prop.~5.5.7.8]{HTT}.

\smallskip\noindent\textit{Step~2 \textup{(}Localizing generation
by $\mathcal{G}_{\mathrm{Bax}}$\textup{)}.}
By Step~1, every compact object lies in the thick closure of
$\iota(D^b(\mathcal{A}))$.
Theorem~\ref{thm:shifted-prefundamental-generation}
\textup{(}Steps~1--3\textup{)} proves that every standard object
of~$\mathcal{A}$ lies in
$\operatorname{thick}\langle \mathcal{G}_{\mathrm{Bax}} \rangle$
inside the separated completion.
The BGG filtration gives every object of~$\mathcal{A}$ a finite
filtration by standards.
Applying the heart-capture criterion
\textup{(}Proposition~\ref{prop:heart-capture-criterion}\textup{)},
every object of~$D^b(\mathcal{A})$, and therefore every compact
object of~$\mathcal{D}^{\mathrm{comp}}$, lies in the thick closure
of~$\mathcal{G}_{\mathrm{Bax}}$.
Since $\mathcal{D}^{\mathrm{comp}}$ is presentable and compactly
generated \textup{(}Step~1\textup{)}, the localizing subcategory
generated by the compact part equals the whole category:
$\mathcal{D}^{\mathrm{comp}} =
\operatorname{Loc}\langle \mathcal{G}_{\mathrm{Bax}} \rangle$.

\smallskip\noindent\textit{Step~3 \textup{(}Localizing DK extension
under the compact-core hypothesis\textup{)}.}
The bar-cobar functor~$\Phi$ is an exact equivalence on the
evaluation-generated core
\textup{(}Theorem~\ref{thm:eval-core-identification}\textup{)}.
The additional hypothesis says precisely that this equivalence has
already been extended to the compact shifted-prefundamental core.
By Step~2, that core generates all compact objects of
$\mathcal{D}^{\mathrm{comp}}$ after idempotent completion.
A localizing exact functor between presentable stable categories
that restricts to an equivalence on compact objects extends uniquely
and is itself an equivalence \textup{(}this is the universal property of
Ind-completion: a functor on compact objects extends uniquely,
Lurie~\cite[Prop.~5.3.5.10]{HTT}\textup{)}.
\end{proof}

\begin{remark}[Characterization of the hypotheses]
\label{rem:completion-extension-hypotheses}
\index{MC3!completion extension!hypothesis analysis}
The three hypotheses isolate precisely the content of
Conjecture~\textup{\ref{conj:dk-compacts-completion}}.
Hypothesis~\ref{it:weight-compatible} is structural: it holds for the
Ind-completion $\operatorname{Ind}(D^b(\mathcal{A}))$
\textup{(}by construction\textup{)} and for any Francis--Gaitsgory
pro-nilpotent completion compatible with the weight grading.
Hypothesis~\ref{it:colimit-generation} holds by definition for any
Ind-completion. The substantive content is
hypothesis~\ref{it:compactness-preservation}: every bounded complex
with finite-length cohomology in~$\mathcal{A}$ must remain compact in
the chosen enhancement.

For the Ind-completion
$\mathcal{D}^{\mathrm{comp}} =
\operatorname{Ind}(D^b(\mathcal{A}))$,
all three hypotheses hold tautologically and the proposition gives
$\operatorname{Ind}(D^b(\mathcal{A})) =
\operatorname{Loc}\langle \mathcal{G}_{\mathrm{Bax}} \rangle$.
This is the unconditional result: the Ind-completion of the shifted
category is generated by the Baxter-locus seed. The DK equivalence on
that Ind-completion still requires the compact-core extension
hypothesis in Proposition~\ref{prop:completion-extension-weight-bounded}(c).

For the Francis--Gaitsgory pro-nilpotent completion, the issue is that
the completion functor $D^b(\mathcal{A}) \to \widehat{D}_{FG}$ may
fail to preserve compactness if the pro-nilpotent tower introduces
non-bounded inverse systems.
Concretely, hypothesis~\ref{it:compactness-preservation} reduces to:
for every finite-length $M \in \mathcal{A}$ and every filtered system
$\{N_\alpha\}$ in~$\widehat{D}_{FG}$,
\[
\varinjlim_\alpha \operatorname{Hom}(\iota(M), N_\alpha)
\;\xrightarrow{\;\sim\;}\;
\operatorname{Hom}\bigl(\iota(M),
\varinjlim_\alpha N_\alpha\bigr).
\]
By the weight-bounded t-structure, this is equivalent to the statement
that the weight-$\mu$ Ext groups
$\operatorname{Ext}^*(\iota(M), N_\alpha)_\mu$
stabilize for each~$\mu$ along the filtered system. The weight
finiteness of~$\mathcal{A}$ \textup{(}each weight space is
finite-dimensional\textup{)} reduces this to a Mittag--Leffler
condition on the weight-graded pieces of the Ext tower.

The pro-Weyl tower $\{W_m\}$ already satisfies Mittag--Leffler
\textup{(}Theorem~\ref{thm:pro-weyl-recovery}\textup{)}: the
transition maps are surjective at each weight. If the pro-nilpotent
completion is built from towers with the same Mittag--Leffler
property, hypothesis~\ref{it:compactness-preservation} follows.
The precise form of the conjecture is therefore: the pro-nilpotent
filtration on the completed category is weightwise Mittag--Leffler.
\end{remark}

\begin{remark}[Reduction of \textup{Conjecture~\ref{conj:dk-compacts-completion}}]
\label{rem:dk-compacts-reduction}
\index{MC3!compact-completion!reduction}
Proposition~\ref{prop:completion-extension-weight-bounded} converts
Conjecture~\ref{conj:dk-compacts-completion} from an abstract
categorical statement \textup{(}``extend the bar-cobar equivalence to
the completed category''\textup{)} into a concrete, checkable
Mittag--Leffler condition on the pro-nilpotent tower. The reduction
proceeds through two independent channels:

\medskip\noindent\emph{Channel~A \textup{(}Neeman/BvdB localizing
generation\textup{)}.} If $\{P_i\}$ compactly generate
$D^b(\mathcal{A})$ as a thick subcategory and the completion
preserves compactness, then $\{P_i\}$ also compactly generate
$\mathcal{D}^{\mathrm{comp}}$ as a localizing subcategory.
The key theorem is
Bondal--Van den Bergh~\cite[Cor.~3.1.2]{BondalVdB03}: for a
compactly generated triangulated category, every compact object is a
retract of an object in the thick closure of the generators.
The Ind-completion $\operatorname{Ind}(D^b(\mathcal{A}))$ is the
universal such category, and the comparison with the
Francis--Gaitsgory completion requires only the compactness
preservation hypothesis.

\medskip\noindent\emph{Channel~B \textup{(}Direct Baxter-locus
finiteness\textup{)}.}
On the Baxter locus $b = a - 1/2$, the short exact sequences from
the Baxter singular vectors give explicit filtrations.
Every compact object $C$ has bounded cohomological amplitude
\textup{(}by definition of compactness in a stable category with a
bounded t-structure, $C$ has finitely many nonzero cohomology
objects\textup{)}, and each cohomology object has a finite standard
filtration \textup{(}by the BGG theorem for highest-weight
categories\textup{)}. Each standard is resolved by a finite Baxter
filtration \textup{(}Step~2 of
Theorem~\ref{thm:shifted-prefundamental-generation}\textup{)}.
The composite filtration is therefore finite for every compact
object: the finiteness of the standard filtration at each
cohomological degree combines with the boundedness of the
cohomological amplitude to give a globally finite iterated filtration.

The two channels are independent and converge to the same conclusion.
Channel~A is categorical and type-independent once the hypotheses are
verified; Channel~B is constructive and gives explicit control over
the filtration length. Channel~A requires the abstract
compactness-preservation hypothesis; Channel~B uses the explicit
Baxter filtrations to verify that hypothesis weight-by-weight.
\end{remark}

\subsection{Landscape of the MC3 extension}
\label{sec:mc3-extension-landscape}

\begin{conjecture}[Type-independence of the remaining MC3 completion mechanisms;
 \ClaimStatusConjectured]
\label{conj:mc3-automatic-generalization}
\label{conj:mc3-proof-step-analysis}% backward compatibility (renamed from prop:)
Let~$\mathfrak{g}$ be simple. Assume that, beyond the all-types
categorical CG theorem and the evaluation-generated-core DK
equivalence, the following packets are available in the relevant
shifted/completed category:
\begin{enumerate}[label=\textup{(\roman*)}]
\item shifted-prefundamental generation on the relevant compact core;
\item the Mittag--Leffler / pro-Weyl hypotheses on the Weyl truncation
 towers;
\item a compatible compact-completion formalism carrying compact-core
 equivalences to the desired completed category.
\end{enumerate}
Then, once package~(i) and DK-2/3 on the evaluation-generated core are
in place, the remaining pro-Weyl and completion arguments of the
type-$A$ proof apply verbatim: they are categorical and do not use
additional type-$A$ combinatorics.
\end{conjecture}

\begin{remark}[Type-independent completion mechanism]
After Theorem~\ref{thm:categorical-cg-all-types} and
Corollary~\ref{cor:dk23-all-types}, the open inputs are the
Mittag--Leffler argument on Weyl truncation towers and the formal
passage from a compact-core equivalence to the chosen completed
category. These manipulations depend on the hypotheses listed above,
not on the type-$A$-specific minuscule/block-separation input used to
establish package~(i).
\end{remark}

\begin{remark}[MC3 package~\textup{(i)} resolved for all types]
\label{rem:mc3-cg-all-types-resolved}
The categorical prefundamental Clebsch--Gordan decomposition
\textup{(}package~\textup{(i)} of MC3\textup{)} is now proved
for all simple types
(Theorem~\ref{thm:categorical-cg-all-types}), using
multiplicity-free $q$-characters
\cite{Nakajima04, FrenkelMukhin01, ChariMoura06}
in place of the minuscule hypothesis.
\end{remark}

\begin{remark}[Type~$B$ folding and uniform CG closure]
\label{rem:mc3-type-b-folding}
Folding from type~$A_{2n-1}$ to type~$B_n$ gives one route to the
type-$B$ prefundamental Clebsch--Gordan identities. The uniform
statement is stronger: Theorem~\ref{thm:categorical-cg-all-types}
proves the CG closure for every simple type directly from the
multiplicity-free $\ell$-weight property.
\end{remark}

\begin{remark}[Difficulty gradient across Dynkin types (resolved)]
\label{rem:mc3-difficulty-gradient}
The multiplicity-free $\ell$-weight argument
(Theorem~\ref{thm:categorical-cg-all-types}) resolves the CG
decomposition \textup{(}package~\textup{(i)}\textup{)} uniformly
for all types:
\begin{center}
\begin{tabular}{cccc}
\toprule
Type & Minuscule count & Key input & CG conclusion \\
\midrule
$A_n$ & $n$ & minuscule (all fund.\ reps) & \textbf{Proved} \\
$B_n$ & $1$ & $\ell$-weight separation & \textbf{Proved} \\
$C_n$ & $1$ & $\ell$-weight separation & \textbf{Proved} \\
$D_n$ & $3$ & $\ell$-weight separation & \textbf{Proved} \\
$E_6$ & $2$ & $\ell$-weight separation & \textbf{Proved} \\
$E_7$ & $1$ & $\ell$-weight separation & \textbf{Proved} \\
$E_8$ & $0$ & $\ell$-weight separation & \textbf{Proved} \\
$F_4$ & $0$ & $\ell$-weight separation & \textbf{Proved} \\
$G_2$ & $0$ & $\ell$-weight separation & \textbf{Proved} \\
\bottomrule
\end{tabular}
\end{center}
The minuscule hypothesis (all ordinary weight multiplicities
equal~$1$) was used in the type~$A$ proof only to ensure block
separation. The strictly weaker condition (that all
$\ell$-weight spaces are one-dimensional) suffices and holds
for every fundamental representation of every simple Lie algebra.
The multiplicity-free property is not contained in a single
reference: for simply-laced types it follows from quiver variety
geometry~\cite{Nakajima04}; for classical types from explicit
$q$-character formulas~\cite{ChariMoura06}; for all types the
Frenkel--Mukhin algorithm~\cite{FrenkelMukhin01} computes the
$q$-character and confirms multiplicity-freeness.
For types $G_2$, $F_4$, $E_8$ with zero minuscule
representations, the $\ell$-weight separation handles the
zero-weight multiplicity (which is $1$, $2$, $8$ respectively
for the smallest fundamental).
\end{remark}

\begin{proposition}[$E_8$ relevant-root uniformity]
\label{prop:e8-root-uniformity}
\ClaimStatusProvedHere
The relevant positive root counts per simple root node of $E_8$ are
\[
[78,\, 98,\, 106,\, 104,\, 97,\, 83,\, 57,\, 92]
\]
(range $57$--$106$ out of $120$ total). This distribution is \emph{uniform}: no node has fewer than $47\%$ of total positive roots as relevant, precluding a low-complexity bootstrap node.
\end{proposition}

\begin{proposition}[Character-level Clebsch--Gordan for all simple types]
\label{prop:character-cg-all-types}
\ClaimStatusProvedHere
The prefundamental character-level Clebsch--Gordan identity
\[
\operatorname{ch}(V_\omega \otimes L^-) \;=\; \sum_{j} \operatorname{ch}(L^-(\mu_j))
\]
holds for all fundamental weights $\omega$ of every simple Lie algebra $\mathfrak{g}$, including types $F_4$, $G_2$, $E_6$, $E_7$, $E_8$.
\end{proposition}

\begin{proof}
The character $\operatorname{ch}(V_\omega) = \sum_\mu
\operatorname{mult}_\omega(\mu)\, e^\mu$ is a finite Laurent
polynomial, and $\operatorname{ch}(L^-)$ is a formal power series.
Their product distributes:
$\operatorname{ch}(V_\omega) \cdot \operatorname{ch}(L^-)
= \sum_\mu \operatorname{mult}_\omega(\mu)\,
e^\mu \cdot \operatorname{ch}(L^-)
= \sum_\mu \operatorname{mult}_\omega(\mu)\,
\operatorname{ch}(L^-(\mathrm{shift} = \mu))$.
No convergence issues arise because $V_\omega$ has finitely many
weights. This argument is purely algebraic and type-independent.
\end{proof}

\subsection{Five chart-local generation mechanisms}
\label{sec:mc3-five-family-strengthening}

The rank-$1$ Baxter hyperplane $b = a - \tfrac12$ is a
coordinate condition on one rational type-$A$ singular-vector chart.
It is not a universal spectral condition on every Yangian-like
category. The five mechanisms below have different scopes. One is the
proved finite-dimensional evaluation core. The others are conditional
generation criteria inside a named completed, twisted, elliptic, or
graded ambient category. None of them identifies the desired
Francis--Gaitsgory completion or extends the bar-cobar equivalence to
the compact shifted-prefundamental core; that remains the separate
packet of Conjecture~\ref{conj:dk-compacts-completion}.

\begin{proposition}[Asymptotic prefundamentals replace the Baxter locus for $Y_\hbar(\mathfrak{sl}_N)$;
 \ClaimStatusConditional]
\label{prop:mc3-asymptotic-prefundamentals-five-family}
\index{MC3!asymptotic prefundamentals}
\index{thick generation!asymptotic route}
Let $\mathcal{C}_{\mathrm{pf}} \subset \widehat{\mathcal{O}}(Y_\hbar(\mathfrak{sl}_N))$
be the monoidal subcategory generated by the family of asymptotic
prefundamental modules $\{L^\pm_{\omega_i, a}\}_{1 \le i \le N-1,\, a \in \mathbb{C}}$
obtained as the Hernandez--Jimbo limit of Kirillov--Reshetikhin modules.
Assume:
\begin{enumerate}[label=\textup{(\alph*)},nosep]
\item the Kirillov--Reshetikhin limits exist as compact objects in
 the chosen separated completion;
\item the Frenkel--Hernandez QQ-system identities lift from
 $K_0$ to finite exact filtrations in that completion;
\item tensoring with finite-dimensional Kirillov--Reshetikhin
 modules preserves compactness.
\end{enumerate}
Then the localizing subcategory generated by the full asymptotic
prefundamental family is independent of the Baxter hyperplane: every
boundary module $L^-_{i,b}$ belongs to
\(\operatorname{Loc}\langle L^\pm_{\omega_j,a}\rangle\) in the chosen
completion. This is a criterion internal to the chosen completion,
not an identification of that completion with the
Francis--Gaitsgory pro-nilpotent completion.
\end{proposition}

\begin{proof}
The Hernandez--Jimbo limit $\lim_{k \to \infty} W^{(i)}_{k,a}$ exists
as a pro-object in the ambient used in~\cite{HJZ25}, and its
$q$-character has the Frenkel--Hernandez product
factorisation~\cite{FrenkelHernandez15}
\[
 \chi_q\!\big(L^\pm_{\omega_i, a}\big) \;=\; \Psi_{i,a}^{\pm 1}
 \prod_{j \neq i} Q_{j}^{\pm \delta_{i,j}}(a),
\]
where the $Q_j$ are the QQ-system generators. Under hypotheses
\textup{(a)--(c)}, these $K_0$ identities are realized by exact
filtrations stable under compact tensoring. The boundary negative
prefundamental appears as the $q=0$ limit of the same KR tower, so it
lies in the generated localizing subcategory. No rank-$1$ Baxter
singular vector is used in this argument.
\end{proof}

\begin{proposition}[Reflection-equation singular vectors for non-type-$A$;
 \ClaimStatusConditional]
\label{prop:mc3-reflection-equation-five-family}
\index{MC3!reflection equation!non-type-A}
\index{thick generation!reflection-equation route}
Let $Y^{\mathrm{tw}}_\hbar(\mathfrak{so}_n)$ or $Y^{\mathrm{tw}}_\hbar(\mathfrak{sp}_n)$
be the twisted Yangian of Molev--Ragoucy type~\cite{MolevRagoucy02,
Molev07} associated with a reflection matrix $K(u)$ satisfying the
boundary Yang--Baxter equation
\[
 R_{12}(u-v)\,K_1(u)\,R_{21}(u+v)\,K_2(v) \;=\;
 K_2(v)\,R_{12}(u+v)\,K_1(u)\,R_{21}(u-v).
\]
Let $\mathcal{C}_{\mathrm{RE}}$ be the monoidal subcategory generated
by evaluation modules $V_{\omega_1, a_i}$. Assume that the
Molev--Ragoucy classification realizes every finite-dimensional
simple in the chosen reflection-equation block as a subquotient of
generic fusion products of these vector evaluation modules, and that
the Sklyanin determinant
\[
 \operatorname{sdet} K(u) \;\in\; Z\!\big(Y^{\mathrm{tw}}_\hbar\big)[[u^{-1}]]
\]
separates the blocks under consideration. Then
\(\mathcal{C}_{\mathrm{RE}}\) thickly generates the bounded derived
category of that finite-dimensional block.
\end{proposition}

\begin{proof}
The reflection equation gives intertwiners on generic fusion
products, and centrality of the Sklyanin determinant
\cite{Sklyanin88, Molev07} fixes the block. By the stated
classification hypothesis, every simple in the block is a subquotient
of an object of \(\mathcal{C}_{\mathrm{RE}}\). Lemma~\ref{lem:composition-thick-generation}
then promotes containment of simples to containment of the whole
bounded derived block. The conclusion is about the twisted Yangian
block just named; it is not a statement about the ordinary Yangian
category~\(\mathcal O\).
\end{proof}

\begin{theorem}[Finite-dimensional evaluation core from Drinfeld polynomials;
 \ClaimStatusProvedHere]
\label{thm:mc3-universal-multiplicity-free-five-family}
\index{MC3!universal thick generation|textbf}
\index{thick generation!multiplicity-free $\ell$-weights}
Let~$\mathfrak{g}$ be any simple Lie algebra and $Y_\hbar(\mathfrak{g})$
its Yangian. The thick subcategory of
\[
D^b(\operatorname{Rep}_{\mathrm{fd}}(Y_\hbar(\mathfrak{g})))
\]
generated by fundamental evaluation modules is the whole bounded
derived category of finite-dimensional modules.
Multiplicity-free $\ell$-weights and the Bethe subalgebra give
generic diagonalization on finite-dimensional evaluation modules, but they do
not enlarge this theorem to the full category~\(\mathcal O\) or to a
completed prefundamental category.
\end{theorem}

\begin{proof}
By the Drinfeld classification of finite-dimensional Yangian modules,
every simple finite-dimensional \(Y_\hbar(\mathfrak g)\)-module is
the irreducible quotient of an ordered tensor product of fundamental
evaluation modules with prescribed Drinfeld polynomials. Hence every
simple object of \(\operatorname{Rep}_{\mathrm{fd}}\) lies in the
thick closure of the evaluation seed. Lemma~\ref{lem:composition-thick-generation}
then gives the whole bounded derived category. The
multiplicity-free \(q\)-character results
\cite{Nakajima04,FrenkelMukhin01,ChariMoura06} refine the generic
Bethe diagonalization picture; they are not the logical source of
generation beyond the finite-dimensional core.
\end{proof}

\begin{proposition}[Elliptic Yangian $E_{\rho,\eta}(\mathfrak{sl}_N)$: theta-divisor failure locus;
 \ClaimStatusConditional]
\label{prop:mc3-elliptic-theta-divisor-five-family}
\index{MC3!elliptic Yangian}
\index{thick generation!theta-divisor route}
Let $E_{\rho,\eta}(\mathfrak{sl}_N)$ be the elliptic Yangian over an
elliptic curve $\mathcal{E}_\rho = \mathbb{C}/(\mathbb{Z} + \rho\mathbb{Z})$
with deformation parameter~$\eta$, and let $\mathcal{C}_{\mathrm{ell}}$
be the monoidal category generated by evaluation modules at generic
parameters. The dynamical Yang--Baxter equation (face-vertex
correspondence of Felder~\cite{FelderVarchenko96}) identifies the
transfer-matrix algebra with the elliptic Bethe subalgebra
\[
 B_{\mathrm{ell}}(u) \;\subset\;
 E_{\rho,\eta}(\mathfrak{sl}_N)\big[[u^{-1}]\big].
\]
Assume the elliptic Bethe eigenvectors are complete, the joint
spectrum is simple, and fusion products have finite composition
series on
\(\mathcal{E}_\rho^{N-1}\setminus\Theta\). Then the evaluation
modules at parameters outside the theta divisor thickly generate the
corresponding finite-dimensional elliptic block. The divisor
\(\Theta\) is the resonance locus where this criterion does not
apply.
\end{proposition}

\begin{proof}
The elliptic Bethe ansatz of Felder--Varchenko~\cite{FelderVarchenko96}
produces joint eigenvectors of $B_{\mathrm{ell}}(u)$ parametrised by
solutions of the elliptic Bethe equations; the simplicity of the joint
spectrum holds outside the zero locus of the Jacobian theta function
of the equations, which is a theta divisor. Under the stated
completeness and finite-length hypotheses, simple spectrum separates
the simple objects in the block, and Lemma~\ref{lem:composition-thick-generation}
converts this separation into thick generation. The proof stops at
the elliptic finite-dimensional block.
\end{proof}

\begin{proposition}[Super-Yangian $Y_\hbar(\mathfrak{gl}_{m|n})$: parity-balance constraint;
 \ClaimStatusConditional]
\label{prop:mc3-super-parity-balance-five-family}
\index{MC3!super-Yangian}
\index{thick generation!parity-balance route}
For the super-Yangian $Y_\hbar(\mathfrak{gl}_{m|n})$ with
$\mathbb{Z}/2$-grading $A = A_{\bar 0} \oplus A_{\bar 1}$, the Nazarov
quantum Berezinian~\cite{Nazarov91} factorises the Baxter operator as
\[
 Q(u) \;=\; Q^+(u)\,\otimes\,Q^-(u),
\]
with $Q^\pm(u)$ acting on the even and odd sectors respectively. Thick
generation of the evaluation-module subcategory
$\mathcal{C}_{m|n} \subset \mathcal{O}(Y_\hbar(\mathfrak{gl}_{m|n}))$
holds in the finite-dimensional parity-balanced block, provided the
graded Drinfeld classification and the super-Shapovalov determinant
give finite-length fusion filtrations there. The balance condition is
\[
 b^+ - b^- \;=\; a - m + n,
\]
and it is a block constraint, not a universal Baxter hyperplane for
ordinary Yangians.
\end{proposition}

\begin{proof}
The Berezinian of Nazarov~\cite{Nazarov91} separates the even and odd
Cartan series, so the graded highest-\(\ell\)-weight data splits into
the two Baxter factors \(Q^\pm\). Under the stated classification and
determinant hypotheses, every simple in the parity-balanced
finite-dimensional block is a subquotient of a tensor product of
evaluation modules in the same block. Lemma~\ref{lem:composition-thick-generation}
then gives the derived thick-generation statement.
\end{proof}

\begin{remark}[Five charts and their scopes]
\label{rem:mc3-unified-mechanism-five-family}
\index{MC3!five-family unification}
\index{Baxter constraint!chart-local}
The five mechanisms have different logical strength:
\begin{center}
\renewcommand{\arraystretch}{1.15}
\begin{tabular}{lll}
\toprule
Family & Mechanism & Scope \\
\midrule
$Y_\hbar(\mathfrak{sl}_N)$ rational & QQ-system / asymptotic KR limit & conditional completion criterion \\
twisted $B/C/D$ Yangians & reflection equation, $\operatorname{sdet}K(u)$ & conditional finite block \\
all simple $\mathfrak g$ & Drinfeld polynomials & proved finite-dimensional core \\
$E_{\rho,\eta}(\mathfrak{sl}_N)$ elliptic & elliptic Bethe / DYBE & conditional off $\Theta$ \\
$Y_\hbar(\mathfrak{gl}_{m|n})$ super & Berezinian parity split & conditional graded block \\
\bottomrule
\end{tabular}
\end{center}
The type-$A$ Baxter equation is therefore chart-local. It supplies
one explicit singular-vector proof on the rational stratum, while the
finite-dimensional evaluation core is already controlled by
Drinfeld-polynomial generation. The post-core completed DK statement
still requires the compact-core extension of \(\Phi\) and the
completion comparison named in
Conjecture~\ref{conj:dk-compacts-completion}.
\end{remark}

\section{Monopole operators as affine Grassmannian modifications}
\label{sec:monopole-affine-grassmannian}
\index{monopole operator!affine Grassmannian}
\index{affine Grassmannian!monopole operators}

The Coulomb branch algebras of
\S\ref{sec:coulomb-branch} connect to the geometric theory of monopole
operators, following \cite{CDG2023}, \S7.2 and \S7.4.

\subsection{Setup: Dirichlet boundary and monopole sectors}
\label{sec:monopole-dirichlet-setup}
\index{Dirichlet boundary condition!monopole sectors}

\providecommand{\GrG}{\operatorname{Gr}_G}

Consider a 3d $\mathcal{N}=2$ gauge theory with gauge group~$G$ on
the half-space $\Sigma \times \mathbb{R}_{\geq 0}$, equipped with
Dirichlet boundary conditions on $\partial(\Sigma \times
\mathbb{R}_{\geq 0}) = \Sigma$. The boundary operator algebra
receives two types of contribution:
\begin{enumerate}[label=(\roman*),nosep]
\item \emph{Perturbative modes}: generated by boundary values of the
 bulk fields, forming a vertex algebra on~$\Sigma$ (the perturbative
 boundary VOA).
\item \emph{Monopole operators}: disorder operators that modify the
 boundary condition at a point $p \in \Sigma$. A monopole operator
 of charge~$\lambda$ replaces the trivial $G$-bundle near~$p$ with a
 bundle whose transition function is the cocharacter
 $z \mapsto z^\lambda$.
\end{enumerate}

\begin{definition}[Monopole operator as Grassmannian modification]
\label{def:monopole-grassmannian}
\index{monopole operator!definition}
\index{affine Grassmannian!as monopole moduli}
A monopole operator of magnetic charge
$\lambda \in X_*(T) \cong \pi_1(\GrG)$ is a modification of the
$G$-bundle at a marked point, i.e.\ a point in the affine
Grassmannian
\begin{equation}\label{eq:affine-grassmannian-def}
 \GrG \;=\; G(\!(z)\!)/G[\![z]\!].
\end{equation}
The connected components $\GrG^\lambda$, indexed by dominant
cocharacters $\lambda \in X_*(T)^+$, each contribute a sector to the
boundary algebra.
\end{definition}

\subsection{Geometric decomposition of the Hilbert space}
\label{sec:monopole-hilbert-decomposition}
\index{Hilbert space!monopole decomposition}

\begin{proposition}[Hilbert space decomposition;
\ClaimStatusProvedElsewhere]
\label{prop:monopole-hilbert-decomp}
\index{Costello--Dimofte--Gaiotto}
\textup{(Costello--Dimofte--Gaiotto \cite{CDG2023}, \S7.2.)}
In the holomorphic-topological twist, the Hilbert space of the
Dirichlet boundary theory decomposes over connected components of
the affine Grassmannian:
\begin{equation}\label{eq:monopole-hilbert-decomp}
 \mathcal{H}_{\mathrm{Dir}}
 \;=\;
 \bigoplus_{\lambda \in X_*(T)^+}
 \mathcal{H}^\lambda,
\end{equation}
where $\mathcal{H}^\lambda$ is the space of states in the monopole
sector of charge~$\lambda$. Correlation functions of monopole
operators in the sector~$\lambda$ are holomorphic functions on the
Schubert cell $\GrG^\lambda$.
\end{proposition}

\begin{proposition}[Dirichlet boundary character;
\ClaimStatusProvedElsewhere]
\label{prop:dirichlet-character}
\index{Dirichlet boundary condition!character formula}
\textup{(Costello--Dimofte--Gaiotto \cite{CDG2023}, \S7.4.)}
The character of the Dirichlet boundary algebra takes the form
\begin{equation}\label{eq:dirichlet-character}
 \chi_{\mathrm{Dir}}(q,\mathbf{x})
 \;=\;
 \sum_{\lambda \in X_*(T)^+}
 q^{E(\lambda)}
 \;\cdot\;
 \chi\bigl(T_{\GrG^\lambda}\bigr)
 \;\cdot\;
 \chi_{\mathrm{matter}}(\lambda,\mathbf{x})
 \;\cdot\;
 \chi_{\mathcal{L}}(\lambda),
\end{equation}
where $E(\lambda)$ is the energy (conformal dimension) of the monopole
sector, $\chi(T_{\GrG^\lambda})$ is the character of the tangent space
to the Schubert cell, $\chi_{\mathrm{matter}}$ encodes the matter
representation contribution, and $\chi_{\mathcal{L}}$ is the line
bundle twist.
\end{proposition}

\subsection{Connection to Coulomb branch algebras}
\label{sec:monopole-coulomb-connection}
\index{Coulomb branch!and monopole operators}

\begin{remark}[Coulomb branch as Grassmannian algebra]
\label{rem:coulomb-grassmannian}
\index{Coulomb branch!Braverman--Finkelberg--Nakajima}
The Coulomb branch algebra $\mathcal{A}(G,N)$ of
Definition~\ref{def:coulomb-branch} is precisely the algebra of
holomorphic functions on a resolution of (a subvariety of) the affine
Grassmannian. The BFN construction
(Theorem~\ref{thm:bfn}) produces
$\mathcal{A}(G,N) \cong H^{G(\!(z)\!)}_*(\mathcal{R}_{G,N})$
as Borel--Moore homology of the variety of triples; the monopole
perspective identifies the same algebra as generated by monopole
operators with their OPE determined by the Grassmannian geometry.
The two viewpoints are related by the Mirkovi\'c--Vilonen
correspondence: weight functors on $\operatorname{Perv}(\GrG)$ are
identified with the monopole charge decomposition
\eqref{eq:monopole-hilbert-decomp}.
\end{remark}

\begin{remark}[Bar complex and monopole sectors]
\label{rem:bar-monopole}
\index{bar complex!monopole sector decomposition}
The bar complex $\bar{B}^{\mathrm{ch}}$ of the Dirichlet boundary
VOA inherits the decomposition \eqref{eq:monopole-hilbert-decomp}:
the bar differential preserves the monopole grading
$\lambda \in X_*(T)^+$, so that
$\bar{B}^{\mathrm{ch}}(\mathcal{H}_{\mathrm{Dir}})$ is an algebraic
direct sum on every finite monopole-charge window. The completed
object used by bar-cobar is the corresponding filtered
charge-completion, equivalently the inverse limit over finite
windows. In that completed/coderived ambient, bar-cobar inversion
(Theorem~B) recovers the chosen Dirichlet boundary algebra from its
bar coalgebra. It does not identify the bar coalgebra with the
chiral derived centre, and it does not replace the BFN convolution
construction of the Coulomb branch.
\end{remark}


\section{Gauge theory Koszul duals as shifted cotangent loop algebras}
\label{sec:gauge-shifted-cotangent}
\index{Koszul dual!gauge theory}
\index{shifted cotangent!loop algebra}

\begin{proposition}[Koszul dual of gauge boundary chiral algebra;
\ClaimStatusProvedElsewhere]
\label{prop:gauge-koszul-dual-shifted-cotangent}
\index{dg-shifted Yangian!as gauge Koszul dual}
\textup{(Dimofte--Niu--Py \cite{DNP2025}, \S4.2.)}
For a 3d $\mathcal{N}=2$ gauge theory with gauge algebra~$\fg$ and
Neumann boundary conditions, the Koszul dual algebra of the boundary
chiral algebra is the \emph{shifted cotangent of the loop algebra}:
\begin{equation}\label{eq:gauge-koszul-shifted-cotangent}
 \cA^!
 \;\cong\;
 T^*[-1]\,\fg[\lambda]^*,
\end{equation}
where $\fg[\lambda]^* = \prod_{n \geq 0} \fg^*\,\lambda^n$ is the
topological dual of the polynomial current algebra, and $T^*[-1]$
denotes the $(-1)$-shifted cotangent bundle (i.e.\ the cotangent
fibers sit in cohomological degree~$1$).
\end{proposition}

\begin{proof}
Dimofte--Niu--Py identify the boundary Koszul dual for Neumann
boundary conditions with the shifted cotangent of the current
algebra~\cite[\S4.2]{DNP2025}. In the pure gauge case the perturbative
bar complex is the relative Chevalley--Eilenberg complex for
$\fg(\!(z)\!)$ over $\fg[\![z]\!]$; dualizing its cohomology gives
\[
\operatorname{Sym}\bigl(\fg^*[z^{-1}]z^{-1}[-1]\bigr)
\cong T^*[-1]\fg[\lambda]^* .
\]
Matter fields tensor in the corresponding boundary contribution, while
the gauge current sector remains the shifted cotangent loop factor. The
BV pairing of the 3d theory supplies the $(-1)$-shifted symplectic form
on this dual algebra.
\end{proof}

\begin{remark}[The dg-shifted Yangian is the gauge Koszul dual]
\label{rem:dg-yangian-gauge-koszul}
\index{dg-shifted Yangian!identification}
Proposition~\ref{prop:gauge-koszul-dual-shifted-cotangent} identifies
the dg-shifted Yangian of the monograph
(Definition~\ref{def:dg-shifted-yangian}) with the Koszul dual of
the gauge boundary chiral algebra. This is the content of the
identification ``line operators $=$ $\cA^!$-modules''
(Vol~II): the category of
line operators in the 3d gauge theory is equivalent to the module
category of the shifted cotangent loop algebra.
Concretely, $T^*[-1]\fg[\lambda]^*$ is the
genus-$0$ piece of the modular dg-shifted Yangian $\Ymod_T$
(Definition~\ref{def:modular-yangian-pro}): the pro-MC element
$R_T^{\mathrm{mod}}(z;\hbar)$ lives in
$\operatorname{MC}(\Ymod_T)$, and its genus-$0$ projection is
the spectral $r$-matrix $r(z)$ acting on
$T^*[-1]\fg[\lambda]^*$.
\end{remark}

\begin{example}[Pure $\fg$-gauge theory]
\label{ex:pure-gauge-koszul-dual}
\index{Koszul dual!pure gauge theory}
For pure $\fg$-gauge theory (no matter), the boundary chiral algebra
is $\widehat{V}_k(\fg)$ and the Koszul dual is
\begin{equation}\label{eq:pure-gauge-koszul}
 \cA^!
 \;=\;
 T^*[-1]\,\fg[\lambda]^*
 \;=\;
 \fg[\lambda]^* \;\oplus\; \fg[\lambda][-1],
\end{equation}
with the differential encoding the coadjoint action. The evaluation
representation $\fg[\lambda]^* \to \fg^*$ (setting $\lambda = a$) gives
the simple modules of $\cA^!$, i.e.\ the elementary line operators of
the gauge theory evaluated at spectral parameter~$a$.
\end{example}

\begin{remark}[Connection to the Coulomb branch]
\label{rem:shifted-cotangent-coulomb}
The shifted cotangent identification
\eqref{eq:gauge-koszul-shifted-cotangent} connects to the Coulomb
branch story of \S\ref{sec:coulomb-branch} through a comparison
map from the perturbative degree-zero part of the Koszul dual
\(\cA^!\) to the local Coulomb-branch chart. Equality with the BFN
Coulomb branch algebra
\(\mathcal{A}(G,N)=H^{G_{\mathcal O}\times\mathbb C^\times}_*
(\mathcal R_{G,N})\)
requires the non-perturbative monopole sectors and their convolution
product; it is not a formal consequence of taking \(H^0(\cA^!)\)
for the perturbative boundary algebra. Thus the objects remain
separate:
\[
\cA,\qquad \barB(\cA),\qquad
\cA^{\mathrm i}=H^\bullet\barB(\cA),\qquad
\cA^!=(\cA^{\mathrm i})^\vee,\qquad
Z^{\mathrm{der}}_{\mathrm{ch}}(\cA).
\]
The full dg structure of $\cA^!$ (equivalently, the genus-$0$ piece
of the modular Yangian $\Ymod_T$,
Definition~\ref{def:modular-yangian-pro}) retains the
perturbative higher homotopy information that is invisible in the
underived Coulomb-branch chart.
The shadow algebra $\cA^{\mathrm{sh}}$
(Definition~\ref{def:shadow-algebra}) of the gauge boundary chiral
algebra then captures the modular invariants: $\kappa$, $\Delta$, and
the contact hierarchy, all as graded projections of the shadow
obstruction tower $\Theta_{\widehat{V}_k(\fg)}^{\leq r}$
(Definition~\ref{def:shadow-postnikov-tower}) at finite order.
\end{remark}

%% ================================================================
%% THE TYPE-A BAXTER--REES COMPACTIFICATION
%% ================================================================

\section{The type-$A$ Baxter--Rees compactification}
\label{sec:typeA-baxter-rees}
\index{Baxter--Rees compactification}
\index{Yangian!Baxter--Rees family}
\index{weightwise MC4}

The representation theory of $Y_\hbar(\fsl_N)$ admits a natural
compactification that unifies three apparently different regimes
into a single algebraic family. Consider the chain of
Kirillov--Reshetikhin modules
\begin{equation}\label{eq:KR-chain-baxter-rees}
 W_0 \xhookrightarrow{F_{1,0}} W_1 \xhookrightarrow{F_{2,1}}
 W_2 \xhookrightarrow{} \cdots
\end{equation}
for a simple node $i \in \{1,\dots,N-1\}$. Write
$W_\infty := \varinjlim_l W_l$ for the inductive limit. Define
$\mathcal{E}_i := W_\infty \otimes_{\CC} \CC[q]$ and, for every
generator~$t$ of the Cartan-doubled Yangian
$Y_\infty(\fsl_N)$, set
\begin{equation}\label{eq:baxter-rees-action}
 \rho_q(t) := C_t + q\,D_t
 \;\in\; \operatorname{End}_{\CC[q]}(\mathcal{E}_i),
\end{equation}
where $C_t$ and $D_t$ are the boundary and deformation coefficients
extracted from the affine-linear stage formula
\begin{equation}\label{eq:baxter-rees-stage}
 t\,F_{k,l} = F_{k,l+1}(C_t^l + q\,D_t^l)\big|_{q = k^{-1}}
\end{equation}
valid for all sufficiently large $k > l$. The fiber at $q = a \neq 0$
recovers (after Cartan renormalization) the ordinary asymptotic module
$L(\Psi_{i,a^{-1}}/\Psi_{i,0})$; the boundary fiber $q = 0$ is the
negative prefundamental module $L^-_{i,0}$. This is the
\emph{Baxter--Rees family}: a flat algebraic family over
$\mathbb{A}^1_q$ interpolating between evaluation-type modules at generic
fibers and the prefundamental boundary at the origin.

\subsection{Three scales: compact, completion, noncompact}

The Yangian $Y_\hbar(\fsl_N)$ decomposes into three
concentric scales, and the Baxter--Rees family makes the passage
between them explicit.

\begin{enumerate}
 \item \textbf{Compact core.}
 The idempotent-complete tensor subcategory
 $\cC_{\mathrm{poly}}$ generated by the vector representation
 $V = \CC^N$ under tensor products, direct summands, and spectral
 shifts (see~\S\ref{sec:coulomb-branch}).
 Every object is a finite direct sum of Schur functors $S_\lambda V$.
 The ordered factorization structure is already visible here: a
 single binary seed $R(u) = r_{V,V}(u) \in \End(V^{\otimes 2})((u^{-1}))$
 determines the entire additive spectral kernel by propagation
 to all Schur-functor pairs.

 \item \textbf{Completion scale.}
 The principal RTT mode tower
 $\{Y^{\mathrm{RTT}}_{\le M}(\fsl_N)\}_{M \ge 1}$
 obtained by truncating the RTT generators $t^{(r)}_{ij}$ at mode
 level~$M$. This tower does not stabilize in ordinary homology
 degree (Proposition~\ref{prop:yangian-failure-unweighted}), but it
 does stabilize \emph{weightwise}
 (Theorem~\ref{thm:yangian-weightwise-MC4}).

 \item \textbf{Noncompact boundary.}
 Asymptotic and negative prefundamental modules, constructed as
 derived telescope fibers of the KR chain
 \eqref{eq:KR-chain-baxter-rees}, then assembled into the
 Baxter--Rees family $\mathcal{E}_i$ over $\mathbb{A}^1_q$.
\end{enumerate}

\begin{remark}[Relation to the gauge Koszul dual]
\label{rem:three-scales-gauge}
\index{three-scale picture}
The three scales are the representation-theoretic avatar of the
three-layer structure appearing on the gauge Koszul side
(Remark~\ref{rem:dg-yangian-gauge-koszul}): the compact core maps to
evaluation line operators, the completion scale to the full
perturbative line-operator category, and the noncompact boundary to
the Stokes/asymptotic sector of the gauge theory.
\end{remark}

\subsection{Seed rigidity on the polynomial core}

The compact core is rigid rather than mysterious.

\begin{definition}[Additive spectral kernel on the polynomial core;
\ClaimStatusProvedHere]
\label{def:yangian-additive-spectral-kernel}
\index{additive spectral kernel}
An \emph{additive spectral kernel} on $\cC_{\mathrm{poly}}$ is a
family of meromorphic operators
$r_{X,Y}(u) \in \End(X \otimes Y)((u^{-1}))$, natural in
$X,Y \in \cC_{\mathrm{poly}}$, satisfying the unit conditions
$r_{\mathbf{1},X} = 0 = r_{X,\mathbf{1}}$ and the additivity law
\begin{equation}\label{eq:yangian-additive-kernel-law}
 r_{X \otimes X',Y}(u) = r_{X,Y}(u)^{13} + r_{X',Y}(u)^{23}.
\end{equation}
\end{definition}

\begin{theorem}[Propagation from the vector seed;
\ClaimStatusProvedHere]
\label{thm:yangian-vector-seed-propagation}
\index{seed propagation}
Let $r$ be an additive $\mathfrak{gl}_N$-equivariant spectral kernel
on $\cC_{\mathrm{poly}}$ with vector seed
$R(u) := r_{V,V}(u) = f(u)\,P + g(u)\,\id_{V^{\otimes 2}}$,
where $P$ is the flip.
Then for every pair of Schur functors
$X = S_\lambda V$, $Y = S_\mu V$:
\begin{equation}\label{eq:yangian-schur-envelope}
 r_{X,Y}(u) = f(u)\,\Omega^{\fsl_N}_{X,Y}
 + g(u)\,|\lambda|\,|\mu|\,\id_{X \otimes Y},
\end{equation}
where $\Omega^{\fsl_N}_{X,Y} = \sum_{a,b}\rho_X(E_{ab})
\otimes \rho_Y(E_{ba})$ is the cross-Casimir.
In particular, the entire kernel is uniquely determined by the
single binary seed~$R(u)$.
\end{theorem}

\begin{proof}
By induction on tensor-power decomposition. The propagation formula
for $r_{V^{\otimes m},V^{\otimes n}}(u)$ yields
$\sum_{i=1}^m \sum_{j=1}^n R_{i,m+j}(u)$. The identity
$\sum_{a,b} E_{ab} \otimes E_{ba} = P$ on $V \otimes V$ then
identifies the double sum with $f(u)\,\Omega_{V^{\otimes m},
V^{\otimes n}} + g(u)\,mn\,\id$. Restricting to Schur-functor
direct summands via Young idempotents gives
\eqref{eq:yangian-schur-envelope}.
\end{proof}

\begin{corollary}[Compact-core rigidity; \ClaimStatusProvedHere]
\label{cor:compact-core-rigidity}
In the strict binary regime (all $A_\infty$-operations $m_k$ with
$k \ge 3$ vanish), the pairwise seed $r_{V,V}(u)$ plus the classical
Yang--Baxter equation force the Maurer--Cartan equation on all
ordered $n$-point configurations simultaneously: if
$R_n(z_1,\dots,z_n) := \sum_{i<j} r_{ij}(z_i - z_j)$, then
$\MC(R_n) = 0$ for every $n \ge 2$.
\end{corollary}

\begin{proof}
Expanding $\MC(R_n) = dR_n + \tfrac12[R_n,R_n]$, the triple
interactions $[r_{ij},r_{ik}] + [r_{ij},r_{jk}] + [r_{ik},r_{jk}]$
are killed by the classical Yang--Baxter equation, and the pairwise
terms $dr_{ij} + \tfrac12[r_{ij},r_{ij}]$ by the binary MC equation.
\end{proof}

\subsection{Weightwise MC4 for the principal RTT tower}

The completion scale requires a refinement of the bar-cobar
convergence statement.

\begin{proposition}[Failure of unweighted stabilization;
\ClaimStatusProvedHere]
\label{prop:yangian-failure-unweighted}
The inverse system
$\{\bar{B}(Y^{\mathrm{RTT}}_{\le M}(\fsl_N))\}_{M \ge 1}$
does not stabilize in ordinary homology degree: the new
mode-$(M+1)$ generators survive in $I/I^2$ at each stage, so the
maps on $H_1$ are never isomorphisms.
\end{proposition}

The cure is a weight filtration. Assign $\wt(t^{(r)}_{ij}) = r$ and
extend multiplicatively. Every reduced bar tensor
$[a_1|\cdots|a_p]$ of total weight~$w$ has each factor of weight at
most~$w$, so for $M \ge w$ no factor involves a mode larger than~$M$.

\begin{theorem}[Weightwise MC4 for the principal RTT tower;
\ClaimStatusProvedHere]
\label{thm:yangian-weightwise-MC4}
\index{MC4!weightwise}
Assume each finite stage
$Y^{\mathrm{RTT}}_{\le M}(\fsl_N)$ lies in the finite-stage
bar-cobar regime. Define the weight-completed bar object
\begin{equation}
 \widehat{\bar{B}}_{\mathrm{wt}}\bigl(Y_\hbar(\fsl_N)\bigr)
 :=
 \prod_{w \ge 0}\varprojlim_M
 \bar{B}^{(w)}\bigl(Y^{\mathrm{RTT}}_{\le M}(\fsl_N)\bigr).
\end{equation}
Then:
\begin{enumerate}
 \item the weight-completed bar object is a separated, complete,
 conilpotent dg coalgebra;
 \item the bar differential is continuous for the product topology;
 \item the completed cobar construction satisfies
 \begin{equation}
 \widehat\Omega\bigl(
 \widehat{\bar{B}}_{\mathrm{wt}}(Y_\hbar(\fsl_N))
 \bigr)
 \xrightarrow{\;\sim\;}
 \widehat{Y}^{\mathrm{RTT}}_{\mathrm{wt}}(\fsl_N).
 \end{equation}
\end{enumerate}
\end{theorem}

\begin{proof}
The inverse system in each fixed weight~$w$ is eventually constant
(at any stage $M \ge w$). Each weight piece of the limit is
canonically the stable value, and the bar differential preserves
total weight, hence acts independently on each stable piece. The
bar-cobar quasi-isomorphism stabilizes at finite stage $M \ge w$
for each weight~$w$; taking the product over~$w$ yields the
stated quasi-isomorphism.
\end{proof}

\begin{remark}[Why weightwise is the correct form of MC4]
\label{rem:yangian-correct-MC4}
The theorem explains why the MC4 comparison
(\S\ref{sec:coulomb-branch}, Remark~\ref{rem:shifted-cotangent-coulomb})
succeeds for $Y_\hbar(\fsl_N)$ despite the failure of unweighted
stabilization: the bar-cobar resolution converges in the weight
topology, and every finite-weight question is decided at a finite
stage of the RTT tower.
\end{remark}

\subsection{The Baxter--Rees family}

\begin{theorem}[Algebraicity of the Baxter--Rees family;
\ClaimStatusProvedHere]
\label{thm:yangian-baxter-rees-algebraicity}
\index{Baxter--Rees family!algebraicity}
The assignment $t \mapsto \rho_q(t) = C_t + q\,D_t$ extends to a
$\CC[q]$-linear representation
\begin{equation}\label{eq:baxter-rees-representation}
 \rho_q \colon Y_\infty(\fsl_N)
 \longrightarrow
 \End_{\CC[q]}(\mathcal{E}_i).
\end{equation}
Hence $\mathcal{E}_i$ is a flat algebraic family of
$Y_\infty(\fsl_N)$-modules over $\mathbb{A}^1_q$.
\end{theorem}

\begin{proof}
For any defining relation $\mathscr{R}(t_1,\dots,t_m) = 0$,
substituting $t_a \mapsto C_{t_a}^l + q\,D_{t_a}^l$ and restricting
to the image of~$W_l$ yields a polynomial $\mathscr{R}_l(q)$ in~$q$
with values in a finite-dimensional space. Evaluating at
$q = k^{-1}$ for infinitely many integers~$k$ gives zero by the
stage formula~\eqref{eq:baxter-rees-stage}, so
$\mathscr{R}_l(q) \equiv 0$. Flatness is immediate because
$\mathcal{E}_i$ is free over $\CC[q]$.
\end{proof}

\begin{theorem}[Generic and boundary fibers; \ClaimStatusProvedHere]
\label{thm:yangian-generic-boundary-fibers}
\index{Baxter--Rees family!fibers}
For every $a \in \CC^\times$, the fiber
$(\mathcal{E}_i)_a := \mathcal{E}_i \otimes_{\CC[q]} \CC_{q=a}$ is
canonically isomorphic to the Cartan-renormalized asymptotic module
$\theta_{i,a^{-1}}^* L(\Psi_{i,a^{-1}}/\Psi_{i,0})$.
At the boundary,
\begin{equation}\label{eq:baxter-rees-boundary-fiber}
 (\mathcal{E}_i)_0 \;\cong\; L^-_{i,0}.
\end{equation}
After spectral shift by~$x$,
$\tau_x^*(\mathcal{E}_i)_0 \cong L^-_{i,x}$.
\end{theorem}

\begin{proof}
The generic fiber action is the specialization of
$\rho_q(t) = C_t + q\,D_t$ at $q = a$, which, by construction
of the renormalized chart, coincides with the asymptotic action
at parameter $a^{-1}$. The $q = 0$ statement is exactly the boundary
action $t \mapsto C_t$ producing the negative prefundamental module.
\end{proof}

The next statement converts the spectral parameter~$z$ of the RTT
formalism into the boundary chart parameter~$q$.

\begin{proposition}[Derived realization of the Baxter--Rees family;
\ClaimStatusProvedHere]
\label{prop:baxter-rees-derived-realization}
Let
\begin{equation}
 \mathfrak{T}_i^{\mathrm{Bax}}
 :=
 \Cone\!\Bigl(
 \bigoplus_{l \ge 0} W_l \otimes \CC[q]
 \xrightarrow{\;\id - F\;}
 \bigoplus_{l \ge 0} W_l \otimes \CC[q]
 \Bigr),
\end{equation}
with $Y_\infty(\fsl_N)$-action on the $l$-th summand given by
$C_t^l + q\,D_t^l$. Then
$H^{-1}(\mathfrak{T}_i^{\mathrm{Bax}}) = 0$ and
$H^0(\mathfrak{T}_i^{\mathrm{Bax}}) \cong \mathcal{E}_i$.
\end{proposition}

\begin{proof}
Injectivity of the transition maps gives $H^{-1} = 0$; the
inductive limit identifies $H^0$ with
$W_\infty \otimes \CC[q] = \mathcal{E}_i$.
\end{proof}

\subsection{The cumulant tower and $H^2$-reduction}

Once the Baxter--Rees family is placed in the ordered modular bar
framework, the higher-genus extension problem is controlled by a
single deformation dg Lie algebra. The key tool is the
\emph{cumulant tower}: for a pairwise MC background
$\rho \in L_{\{1,2\}}^1$ and any ordered set~$I$, the pairwise field
$R_I(\rho) = \sum_{i<j} \rho_{ij}$ has aggregate curvature
$K_I(\rho) = \MC_I(R_I(\rho))$, and the \emph{ordered cumulant}
\begin{equation}\label{eq:ordered-cumulant}
 \fC_I(\rho) := \pr_I^{\mathrm{full}}\,K_I(\rho)
 \;\in\; L_I\langle I\rangle
\end{equation}
is its full-support projection. A Boolean M\"obius inversion relates
curvatures and cumulants, and in the strict dg Lie regime the
cumulant $\fC_I(\rho)$ vanishes for $|I| \ge 4$: the entire
higher-point Maurer--Cartan system collapses to the two-leg and
three-leg sectors
(cf.\ Definition~\ref{def:modular-yangian-pro}).

\begin{definition}[Exact-support deformation dg Lie algebra;
\ClaimStatusProvedHere]
\label{def:yangian-exact-support-dg-lie}
\index{exact-support deformation!dg Lie algebra}
Fix a pairwise MC background~$\rho$. Define
\begin{equation}
 \fg_{\mathrm{ex}}(\rho)
 :=
 \prod_{|I| \ge 3} L_I\langle I\rangle,
\end{equation}
with differential $\delta_\rho$ and bracket $[-,-]_{\mathrm{ex}}$
given by full-support projection of the twisted differential and
bracket. Then $(\fg_{\mathrm{ex}}(\rho), \delta_\rho,
[-,-]_{\mathrm{ex}})$ is a complete dg Lie algebra, and $x \in
\fg_{\mathrm{ex}}(\rho)^1$ is Maurer--Cartan if and only if
$\cA_I(x) := R_I + \widetilde{x}_I$ is Maurer--Cartan on every
ordered set~$I$.
\end{definition}

The decisive reduction theorem places the first obstruction squarely
in the three-leg sector.

\begin{theorem}[$H^2$-reduction to the three-leg sector;
\ClaimStatusProvedHere]
\label{thm:yangian-H2-reduction}
\index{$H^2$-reduction theorem}
Assume a prime-form/Hodge graph realization with connective
coefficient complexes. Then:
\begin{equation}\label{eq:H2-three-leg}
 H^2\bigl(\fg_{\mathrm{ex}}(\rho)\bigr)
 \;\cong\;
 H^2\bigl(L_{\{1,2,3\}}\langle\{1,2,3\}\rangle,\,\delta_\rho\bigr).
\end{equation}
Equivalently, the first obstruction to extending the pairwise MC
background $\rho$ to a full modular completion is purely
three-legged.
\end{theorem}

\begin{proof}
Filter $\fg_{\mathrm{ex}}(\rho)$ by geometric degree (= number of
prime-form edges). On the $E_1$-page of the resulting spectral
sequence, total degree~$2$ forces geometric degree $p \ge |I| - 1$
(connected graphs on $|I|$ vertices have $\ge |I|-1$ edges) and
coefficient degree $q \ge 0$ (connectiveness). Hence
$|I| - 1 \le 2$ and $|I| \ge 3$, so $|I| = 3$, $p = 2$, $q = 0$.
The full-support differential is componentwise in the ordered set~$I$.
There is no degree-$1$ full-support term for any \(|I|\ge3\), since a
connected full-support graph on three vertices already has two edges.
Thus there are no incoming boundaries to degree~$2$. For
\(|I|\ge4\), the degree-$2$ cochain group itself vanishes by the same
edge count. The only degree-$2$ cochains of
\(\fg_{\mathrm{ex}}(\rho)\) are therefore the three-vertex
full-support cochains, with the differential induced by
\(\delta_\rho\) on
\(L_{\{1,2,3\}}\langle\{1,2,3\}\rangle\). This proves
\eqref{eq:H2-three-leg}.
\end{proof}

Under modular Arnold normalization (the higher-genus upgrade of the
classical Arnold relation $d\log(z_1 - z_2) \wedge d\log(z_1 - z_3)
\sim d\log(z_1 - z_2) \wedge d\log(z_2 - z_3)$), the three Y-graph
$2$-forms represent a single cohomology class $[\Omega_{123}]$, and
the three-leg obstruction becomes
\begin{equation}\label{eq:modular-YB-identification}
 H^2\bigl(\fg_{\mathrm{ex}}(\rho)\bigr)
 \;\cong\;
 [\Omega_{123}] \cdot
 H^0\!\bigl(\Span\{\mathrm{YB}(\rho)\},\,d_{\mathrm{coeff}}\bigr),
\end{equation}
where
$\mathrm{YB}(\rho) = [\rho_{12},\rho_{13}] + [\rho_{12},\rho_{23}]
+ [\rho_{13},\rho_{23}]$ is the coefficient Yang--Baxter element.
In particular, the modular completion is unobstructed in degree~$2$
whenever $[\mathrm{YB}(\rho)] = 0$.

\begin{remark}[Three-leg criterion]
\label{rem:yangian-frontier-shift}
This converts the modular dg-shifted Yangian existence problem from a
tower of unrelated higher-genus equations to a single explicit
cohomological criterion: kill the three-leg Yang--Baxter class.
Everything higher begins only after this class vanishes, and if it
does, the gauge-normalized modular completion is given by an explicit
planar-tree expansion
(Theorem~\ref{thm:yangian-H2-reduction}).
\end{remark}

\subsection{The boundary Kodaira--Spencer class}

The Baxter--Rees family carries a canonical first-order invariant
encoding the infinitesimal deformation of the prefundamental module
$L^-_{i,0}$ along the boundary chart.

\begin{definition}[Baxter--Kodaira--Spencer class;
\ClaimStatusProvedHere]
\label{def:yangian-baxter-KS-class}
\index{Kodaira--Spencer class!Baxter boundary}
Write $M_i^- := (\mathcal{E}_i)_0 \cong L^-_{i,0}$. The reduction
of $\mathcal{E}_i$ modulo $q^2$ yields an exact sequence
\begin{equation}\label{eq:baxter-KS-sequence}
 0 \longrightarrow M_i^- \longrightarrow
 \mathcal{E}_i/q^2\mathcal{E}_i \longrightarrow
 M_i^- \longrightarrow 0,
\end{equation}
where the left inclusion is multiplication by~$q$. The resulting
self-extension class
$\kappa_i \in \Ext^1_{Y_\infty(\fsl_N)}(M_i^-,\,M_i^-)$
is the \emph{Baxter--Kodaira--Spencer class}.
\end{definition}

\begin{proposition}[Concrete cocycle; \ClaimStatusProvedHere]
\label{prop:yangian-baxter-KS-cocycle}
The class $\kappa_i$ is represented by the $1$-cocycle
$\delta_i(t) := D_t$ relative to the boundary action
$\rho_0(t) = C_t$:
\begin{equation}\label{eq:baxter-KS-cocycle}
 \rho_q(t) = \rho_0(t) + q\,\delta_i(t) \pmod{q^2}.
\end{equation}
\end{proposition}

\begin{proof}
The cocycle condition
$\delta_i(ab) = \rho_0(a)\,\delta_i(b) + \delta_i(a)\,\rho_0(b)$
follows from comparing the $q$-coefficients of
$\rho_q(ab) = \rho_q(a)\,\rho_q(b)$. This is exactly the condition
describing the associated square-zero self-extension.
\end{proof}

The Kodaira--Spencer class is the RTT relation rewritten as a
boundary deformation: at $q = 0$ the RTT generators act on the
prefundamental module by the boundary operators $C_t$, and the linear
correction $q\,D_t$ is the first-order deformation lifting $L^-_{i,0}$
along the Baxter--Rees chart. When the Baxter--Rees family carries
polynomial $R$-matrices and Theta-associators (cf.\ the shifted Yangian
theory of Hernandez--Zhang~\cite{HZ24} and Zhang~\cite{ZhangTheta24}),
the weightwise polynomial continuation principle extends both braiding
and associativity data to the boundary. The resulting \emph{formal
braided boundary germ} at each simple root equips $L^-_{i,0}$ with
boundary intertwiners
$\mathscr{R}_{M,i}^{\partial}(u) \colon M \otimes L^-_{i,0} \to
L^-_{i,0} \otimes M$ satisfying the linearized Yang--Baxter equation
around the boundary, and the tangent tensors
$\dot{\mathscr{R}}_{M,i}(u) := (d/dq)|_{q=0}\,\mathscr{R}_{M,i}(u,q)$
are cocycles for the linearized intertwining and pentagon relations.

\begin{remark}[Connection to the modular Yangian]
\label{rem:yangian-baxter-rees-modular}
\index{modular dg-shifted Yangian!Baxter--Rees connection}
The Baxter--Rees compactification is the representation-theoretic
shadow of the modular dg-shifted Yangian $\Ymod_T$
(Definition~\ref{def:modular-yangian-pro}). The three-scale picture
(compact core, completion, noncompact boundary) corresponds to
the three layers of the pro-MC element $R_T^{\mathrm{mod}}(z;\hbar)$:
its genus-$0$ projection is the spectral $r$-matrix (compact core),
its RTT mode expansion is the completion scale (weightwise MC4), and
its boundary germ at each simple root is the Baxter--Kodaira--Spencer
deformation (noncompact boundary). The $H^2$-reduction theorem
\eqref{eq:H2-three-leg} identifies the modular Yang--Baxter class as
the sole obstruction to assembling these three scales into a single
modular completion, making the type-$A$ derived Drinfeld--Kohno
comparison a concrete cohomological verification problem rather
than a slogan.
\end{remark}

\section{Synthesis}
\label{sec:yangcomp-supplement}

\begin{remark}[The fourth taxonomic slot: K3 chiral Hall--Drinfeld double]
\label{rem:yangcomp-1}
The computations collected here survey three classical quantum-group
taxa -- Drinfeld--Jimbo $U_q(\mathfrak g)$ at finite type, Yangians
$Y_\hbar(\mathfrak g)$ via Drinfeld 1985, and RTT algebras
$U^{\mathrm{RTT}}(R)$. The K3 chiral Hall--Drinfeld double
\[
\mathbf{H}_{\Delta_5} \;=\; \mathcal{D}_\hbar\!\left(\mathcal{Y}^{\mathrm{Hall}}_\hbar(\mathrm{CoHA}_{K3\times E}),\,
\widetilde{\Phi}^{\mathrm{Sieg\text{-}Bor}}_{\mathrm{Sp}_4}[\Phi_{10}^{\mathrm{un}}/\eta^{24}],\,
R_{\mathrm{Sieg,dyn}}\right)
\]
occupies a fourth taxonomic slot: its underlying Miki generators are
Hall-theoretic rather than Cartan-theoretic, its coproduct is twisted
by the Siegel--Borcherds associator $\widetilde{\Phi}^{\mathrm{Sieg\text{-}Bor}}_{\mathrm{Sp}_4}$
encoded by $\Phi_{10}^{\mathrm{un}}/\eta^{24}$, and its $R$-matrix is
Siegel-elliptic-dynamical rather than rational, trigonometric, or
elliptic in the classical sense of Belavin--Drinfeld. At Lie/Hopf
level the Hall side is abelian in its Miki-toroidal pieces; the
non-abelian BKM bracket appears only after vertex-operator closure on
the K3 Fock space. It is not a Drinfeld Yangian
\(Y_\hbar(\mathfrak g_{\Delta_5})\). The explicit
Miki-generator computations in this chapter are the
\(\hbar\)-formal shadow of this fourth taxon. Under the CY-3 specialisation
\(\mathrm{CoHA}_{K3\times E}\to\mathrm{CoHA}_{K3}\), selected Hall/CoHA
charts admit RTT-generator shadows comparable with the classical
Yangian generators tabulated here; this is a Hall degeneration, not a
Drinfeld-Yangian identification. See Vol.~III,
Chapter~\ref{V3-ch:k3-chiral-bialgebra-platonic} and Schiffmann--Vasserot 2012,
Davison 2017, Nakajima 2001.
\[
\Phi_{10}^{\mathrm{un}}=\Delta_5^2,\qquad
D_5=64^{-1}\Delta_5,\qquad
\Phi_{10}^{\mathrm{OP}}=D_5^2=4096^{-1}\Delta_5^2.
\]
Thus the Hall-double associator uses the unnormalised Igusa square;
Oberdieck--Pandharipande reduced DT theory instead uses
$Z_{\mathrm{OP/DT}}=-D_5^{-2}=-4096\Delta_5^{-2}$. The
$K3\times E$ scalars remain separated as
\[
\Spec_{\kappa_\bullet}(K3\times E)=\{0,3,5,24\}
\]
for categorical, Heisenberg-chiral, BKM, and fibre entries
respectively.
\end{remark}

\begin{remark}[Siegel-elliptic dynamical $R$-matrix and the universal identity]
\label{rem:yangcomp-2}
The dynamical shift $R_{\mathrm{Sieg,dyn}}(z_1,z_2;\tau,w)$ arises as
the quasi-classical limit of the trigonometric $R$-matrix pulled back
along Pasol--Zagier 2013's universal Igusa form $\Phi_{10}^{\mathrm{un}}$ over the
Siegel upper half-space $\mathfrak H_2$. At the Lusztig specialisation
$\hbar^2 = -1/\ell$ with $\ell = 8$ the universal identity
$\hbar^2 \cdot K^{\kappa_{\mathrm{ch}}} = -1$ couples three faces of $8$:
Mukai doubling $K^{\kappa_{\mathrm{ch}}} = 2c_+(\mathrm{Mukai}(K3)) = 8$,
monodromy order of $L^{\Delta_5}$ around the Humbert divisor $H_1\subset\bar{\mathcal A}_2$
(via Bruinier 2002 Prop.~5.1 Heegner Chern-class reciprocity), and the
Lusztig root-of-unity parameter $\ell = 8$. This is why the na\"ive
$Y_\hbar(\mathfrak{sl}_2)$ computations in
\S\ref{sec:sl2-yangian-explicit} admit a curved lift only at $\hbar^2 = -1/8$,
a constraint invisible from the classical Yangian viewpoint.
See Bruinier 2002 (arXiv:math/0108079), Lusztig 1990, Mukai 1987.
\end{remark}

\section{PBW $\rho$-orbit truncation at \texorpdfstring{$\zeta_8$}{zeta8}}
\label{sec:u-zeta-8-pbw-truncation}

At the Lusztig small quantum group $u_{\zeta_8}(\mathfrak{sl}_N)$ the
De Concini--Kac PBW basis
$\{\prod_{\alpha\in R^+} E_\alpha^{k_\alpha}\}_{0\le k_\alpha\le \ell-1}$
with $\ell = 8$ inherits a Borcherds $\rho$-grading through the BKM
envelope of~$\Phi$ on the Igusa side of the bridge. The number
$d(N)$ below is the dimension of this single $\rho$-orbit
PBW-weight truncation. It is not the ordinary vector-space dimension
of the small quantum group and not a quantum-loop/Yangian evaluation
dimension. In the normalization used here the computed pre-wall
values are
\[
d(1)=2,\qquad d(2)=22,\qquad d(3)=238,\qquad d(4)=366,
\]
with increments $20,216,128$ on the Borcherds orbit.

\begin{theorem}[Formal PBW increment past the De Concini--Kac wall
$N = \ell/2 = 4$;
\ClaimStatusProvedHere]
\label{thm:u-zeta-8-PBW-wall-crossing}
\index{De Concini-Kac!wall crossing}%
\index{small quantum group!PBW truncation}%
Assume that the post-wall $\rho$-orbit truncation is governed by the
De Concini--Kac linked tail seeded by
$\delta(4)=d(4)-d(3)=128=2^7$ and by the fourfold decay
$\delta(N+1)=\delta(N)/4$ for $N=4,5,6$. Then the
PBW-truncation increment $\delta(N) := d(N) - d(N-1)$ satisfies
\begin{equation}\label{eq:u-zeta-8-PBW-increment}
\delta(N) \;=\; 2^{15 - 2N},\qquad N\in\{4,5,6,7\},
\end{equation}
equivalently
\begin{equation}\label{eq:u-zeta-8-PBW-d-values}
d(5) = 398,\qquad d(6) = 406,\qquad d(7) = 408,
\end{equation}
and the formal geometric tail has rational closure
\[
d(4)+\sum_{N\ge 5}2^{15-2N}
=366+\frac{128}{3}
=\frac{1226}{3}.
\]
The verified integer window therefore terminates at the plateau value
$d(7)=408$.
\end{theorem}

\begin{proof}
The recurrence gives
$\delta(5)=128/4=32$, $\delta(6)=8$, and $\delta(7)=2$, which is
exactly~\eqref{eq:u-zeta-8-PBW-increment}. Adding these increments to
$d(4)=366$ gives~\eqref{eq:u-zeta-8-PBW-d-values}. The rational
closure is the geometric sum
$366+32(1+1/4+1/16+\cdots)=366+128/3$.
\end{proof}

\begin{remark}[Plateau and the Lusztig Frobenius kernel;
\ClaimStatusProvedHere]
\label{rem:u-zeta-8-PBW-plateau}
The verified integer plateau is the statement
$d(5)=398$, $d(6)=406$, $d(7)=408$ in the
$\rho$-orbit truncation. It is not a literal
$\mathbb{F}_8$-dimension of a Frobenius-kernel quotient: the
Lusztig Frobenius is a root-of-unity comparison over a coefficient
field containing $\zeta_8$, whereas the present number records only
the Borcherds-weight orbit after the De Concini--Kac linkage wall.
Thus the plateau belongs to the $\rho$-graded truncation model, not
to the full small quantum group.
\end{remark}
